\chapter{Grupuri}

\section{Definiții preliminare}

\subsection{Legi de compoziție internă}

\begin{definitie}
O \emph{lege de compoziție internă} pe mulțimea nevidă $M$ este o funcție $*:M\times M\to M$, $(x,y)\mapsto x*y$. O submulțime $H\subseteq M$ este \emph{parte stabilă} dacă $x,y\in H\Rightarrow x*y\in H$; în acest caz $*$ induce o lege de compoziție pe $H$.
\end{definitie}

\begin{definitie}
Legea $*$ este \emph{asociativă} dacă $(x*y)*z=x*(y*z)$ pentru orice $x,y,z$, respectiv \emph{comutativă} dacă $x*y=y*x$ pentru orice $x,y$. Un element $e\in M$ este \emph{element neutru} dacă $x*e=e*x=x$ pentru orice $x\in M$. Un element $x$ este \emph{simetrizabil} dacă există $x'\in M$ cu $x*x'=x'*x=e$; $x'$ se numește \emph{simetricul} lui $x$. O structură $(M,*)$ asociativă și cu element neutru se numește \emph{monoid}.
\end{definitie}

\begin{definitie}[semigrup]
O pereche $(S,\cdot)$ formată dintr-o mulțime nevidă și o lege de compoziție \emph{asociativă} se numește \emph{semigrup}; dacă legea este și comutativă, semigrupul se numește \emph{comutativ} (abelian). Un semigrup cu element neutru este exact un \emph{monoid} (definit mai sus). Astfel orice grup este monoid și orice monoid este semigrup (grupuri $\subset$ monoizi $\subset$ semigrupuri, ca familii de structuri): trecerea de la semigrup la monoid adaugă axioma elementului neutru, iar cea de la monoid la grup adaugă simetrizabilitatea tuturor elementelor.
\end{definitie}

\begin{propozitie}[unicitate]\label{prop:unicitate}
\emph{(i)} Dacă legea are element neutru, acesta este unic. \emph{(ii)} Într-un monoid, simetricul unui element simetrizabil este unic. \emph{(iii)} Dacă $x,y$ sunt simetrizabile într-un monoid, atunci $x*y$ este simetrizabil și $(x*y)'=y'*x'$ (atenție la ordine!), iar $(x')'=x$.
\end{propozitie}

\begin{proof}
(i) Dacă $e$ și $e'$ sunt neutre, $e=e*e'=e'$. (ii) Dacă $x'$ și $x''$ sunt simetrice ale lui $x$, atunci
$x''=x''*e=x''*(x*x')=(x''*x)*x'=e*x'=x'$. (iii) Calcul direct:
$(x*y)*(y'*x')=x*(y*y')*x'=x*x'=e$ și analog în cealaltă ordine; egalitatea $(x')'=x$ rezultă din simetria definiției.
\end{proof}

\begin{observatie}[notații și puteri]
În notație \emph{multiplicativă} scriem $xy$, $1$ sau $e$ pentru neutru și $x^{-1}$ pentru simetric (numit \emph{invers}); în notație \emph{aditivă} scriem $x+y$, $0$ și $-x$ (numit \emph{opus}). Definim $x^0=e$, $x^{n+1}=x^n\cdot x$ și, pentru elemente simetrizabile, $x^{-n}=(x^{-1})^n$. Într-un monoid au loc $x^{m+n}=x^m x^n$ și $(x^m)^n=x^{mn}$ pentru toți exponenții pentru care puterile au sens. \textbf{Capcană:} în general $(xy)^n\neq x^n y^n$; egalitatea are loc pentru orice $n$ dacă $xy=yx$ (inducție).
\end{observatie}

\begin{exemplu}
$\NN$ este parte stabilă a lui $(\ZZ,+)$, dar în $(\NN,+)$ singurul element simetrizabil este $0$. În monoidul $(\ZZ,\cdot)$ elementele simetrizabile sunt $\pm 1$.
\end{exemplu}

\subsection{Grupuri: definiție și exemple fundamentale}

\begin{definitie}
$(G,\cdot)$ este \emph{grup} dacă legea este asociativă, are element neutru și orice element este simetrizabil. Grupul este \emph{abelian} (comutativ) dacă legea este comutativă. \emph{Ordinul grupului} este cardinalul $|G|$.
\end{definitie}

\begin{exemplu}[exemplele fundamentale]\label{ex:grupuri}
\begin{enumerate}[label=(\alph*),nosep]
\item Grupuri aditive: $(\ZZ,+)$, $(\QQ,+)$, $(\RR,+)$, $(\CC,+)$ --- abeliene, infinite.
\item Grupuri multiplicative: $(\QQ^*,\cdot)$, $(\RR^*,\cdot)$, $(\CC^*,\cdot)$, $((0,\infty),\cdot)$.
\item $(\ZZ_n,+)$ --- grupul claselor de resturi modulo $n$, abelian, $|\ZZ_n|=n$.
\item $U_n=\{z\in\CC\mid z^n=1\}=\{1,\varepsilon,\varepsilon^2,\dots,\varepsilon^{n-1}\}$, unde $\varepsilon=\cos\frac{2\pi}{n}+i\sin\frac{2\pi}{n}$: \emph{grupul rădăcinilor de ordinul $n$ ale unității}, abelian, cu $n$ elemente, față de înmulțire.
\item $(S_n,\circ)$ --- grupul permutărilor mulțimii $\{1,\dots,n\}$, $|S_n|=n!$, \emph{neabelian} pentru $n\ge 3$.
\item $GL_2(\RR)$ --- grupul matricelor $2\times 2$ inversabile, cu înmulțirea; neabelian.
\item $G=\RR\setminus\{1\}$ cu $x\circ y=x+y-xy$: neutrul este $0$, iar simetricul lui $x$ este $\frac{x}{x-1}$ (verificare directă). Vom vedea în §\ref{sec:morfisme} că $(G,\circ)\simeq(\RR^*,\cdot)$.
\end{enumerate}
\end{exemplu}

\begin{propozitie}[reguli de calcul în grup]\label{prop:simplificare}
Într-un grup $G$: \emph{(i)} au loc \emph{simplificările} $ax=ay\Rightarrow x=y$ și $xa=ya\Rightarrow x=y$; \emph{(ii)} ecuațiile $ax=b$ și $xa=b$ au soluții unice, $x=a^{-1}b$, respectiv $x=ba^{-1}$.
\end{propozitie}

\begin{proof}
(i) Înmulțim la stânga (respectiv la dreapta) cu $a^{-1}$. (ii) Verificare directă; unicitatea rezultă din (i).
\end{proof}

\begin{observatie}
Din (ii): în tabla unui grup finit, fiecare element apare \textbf{exact o dată} pe fiecare linie și pe fiecare coloană (proprietatea de ,,sudoku``). Este primul test de verificat când se cere să se decidă dacă o tablă dată definește un grup.
\end{observatie}

\begin{propozitie}\label{prop:x2e}
Dacă $x^2=e$ pentru orice $x\in G$, atunci $G$ este abelian.
\end{propozitie}

\begin{proof}
Ipoteza spune că fiecare element este propriul invers. Atunci $xy=(xy)^{-1}=y^{-1}x^{-1}=yx$.
\end{proof}

\subsection{Subgrupuri. Ordinul unui element}

\begin{definitie}
$H\subseteq G$ este \emph{subgrup} al grupului $G$ (notăm $H\le G$) dacă $H$ este parte stabilă și $(H,\cdot)$ este grup cu legea indusă. Subgrupurile $\{e\}$ și $G$ se numesc \emph{improprii}.
\end{definitie}

\begin{propozitie}[criteriul de subgrup]\label{prop:criteriu}
Fie $H\subseteq G$ nevidă. Atunci $H\le G$ dacă și numai dacă $xy^{-1}\in H$ pentru orice $x,y\in H$.
\end{propozitie}

\begin{proof}
Necesitatea este clară. Reciproc: luând $y=x$ obținem $e=xx^{-1}\in H$; apoi $x^{-1}=ex^{-1}\in H$ pentru orice $x\in H$; în fine $xy=x(y^{-1})^{-1}\in H$, deci $H$ este parte stabilă, iar asociativitatea se moștenește.
\end{proof}

\begin{propozitie}[cazul finit]\label{prop:finitstabil}
Dacă $H\subseteq G$ este \textbf{finită}, nevidă și stabilă la operație, atunci $H\le G$.
\end{propozitie}

\begin{proof}
Fie $x\in H$. Puterile $x,x^2,x^3,\dots$ rămân în $H$, care este finită, deci există $i<j$ cu $x^i=x^j$. Simplificând în $G$, $x^{j-i}=e\in H$. Dacă $j-i=1$, atunci $x=e$ și $x^{-1}=e\in H$; altfel $x^{-1}=x^{j-i-1}\in H$. Concluzia rezultă din Propoziția~\ref{prop:criteriu}.
\end{proof}

\begin{lema}[un grup nu e reuniunea a două subgrupuri proprii]\label{lema:reuniune}
Dacă $H,K\le G$ și $H\cup K=G$, atunci $H=G$ sau $K=G$.
\end{lema}

\begin{proof}
Presupunem, prin absurd, că $H\neq G$ și $K\neq G$. Cum $H\cup K=G$ dar $K\neq G$, există $h\in H\setminus K$; simetric, există $k\in K\setminus H$. Elementul $hk$ aparține lui $G=H\cup K$, deci $hk\in H$ sau $hk\in K$.
\begin{itemize}[nosep,leftmargin=1.4em]
  \item Dacă $hk\in H$, atunci $k=h^{-1}(hk)\in H$, contrazicând $k\notin H$.
  \item Dacă $hk\in K$, atunci $h=(hk)k^{-1}\in K$, contrazicând $h\notin K$.
\end{itemize}
Ambele cazuri sunt imposibile, deci presupunerea a fost falsă: $H=G$ sau $K=G$.
\end{proof}

\begin{observatie}
Rezultatul este optim: un grup \emph{poate} fi reuniunea a \textbf{trei} subgrupuri proprii --- de pildă grupul lui Klein $\ZZ_2\times\ZZ_2$ este reuniunea celor trei subgrupuri de ordin $2$ ale sale. Lema arată doar că două nu ajung niciodată.
\end{observatie}

\begin{definitie}
Pentru $x\in G$, mulțimea $\langle x\rangle=\{x^k\mid k\in\ZZ\}$ este un subgrup, numit \emph{subgrupul ciclic generat de $x$}. \emph{Ordinul} lui $x$, notat $\ord(x)$, este cel mai mic $n\ge 1$ cu $x^n=e$, dacă există; altfel $x$ are ordin infinit.
\end{definitie}

\begin{teorema}[proprietățile ordinului]\label{teo:ordin}
Fie $x\in G$ cu $\ord(x)=n$ finit. Atunci:
\emph{(i)} $x^m=e\iff n\mid m$;\quad
\emph{(ii)} $x^i=x^j\iff i\equiv j \pmod n$;\quad
\emph{(iii)} $\langle x\rangle=\{e,x,\dots,x^{n-1}\}$ are exact $n$ elemente, adică $|\langle x\rangle|=\ord(x)$.
Dacă $x$ are ordin infinit, toate puterile $x^k$, $k\in\ZZ$, sunt distincte.
\end{teorema}

\begin{proof}
(i) Dacă $n\mid m$, evident $x^m=e$. Reciproc, scriem $m=nq+r$ cu $0\le r<n$; atunci $e=x^m=(x^n)^q x^r=x^r$, iar minimalitatea lui $n$ forțează $r=0$. (ii) $x^i=x^j\iff x^{i-j}=e$, apoi (i). (iii) Din (ii), elementele $e,x,\dots,x^{n-1}$ sunt distincte, iar orice putere $x^k$ coincide cu $x^{k \bmod n}$. Cazul infinit: $x^i=x^j$ cu $i\neq j$ ar da $x^{|i-j|}=e$, contradicție.
\end{proof}

\begin{propozitie}[ordinul unei puteri]\label{prop:ordputere}
Dacă $\ord(x)=n$ și $k\in\ZZ$, atunci $\ord(x^k)=\dfrac{n}{(n,k)}$, unde $(n,k)$ este c.m.m.d.c.
\end{propozitie}

\begin{proof}
Fie $d=(n,k)$. Avem $(x^k)^m=e\iff n\mid km\iff \frac{n}{d}\mid \frac{k}{d}\,m\iff \frac{n}{d}\mid m$, ultima echivalență deoarece $\bigl(\frac{n}{d},\frac{k}{d}\bigr)=1$ (lema lui Gauss din aritmetică). Cel mai mic astfel de $m\ge1$ este $\frac{n}{d}$.
\end{proof}

\begin{propozitie}[ordinul: invers și conjugare]\label{prop:ordinv}
În orice grup $G$: \emph{(i)} $\ord(x^{-1})=\ord(x)$; \emph{(ii)} $\ord(g^{-1}xg)=\ord(x)$ pentru orice $g\in G$; \emph{(iii)} $\ord(xy)=\ord(yx)$ pentru orice $x,y\in G$.
\end{propozitie}

\begin{proof}
(i) $(x^{-1})^k=(x^k)^{-1}$, deci $x^k=e\iff(x^{-1})^k=e$: cele două elemente au exact aceleași puteri egale cu $e$. (ii) $(g^{-1}xg)^k=g^{-1}x^kg$ (produsul se telescopează), deci $(g^{-1}xg)^k=e\iff x^k=e$. (iii) $yx=x^{-1}(xy)x$, apoi aplicăm (ii). Toate afirmațiile rămân valabile și pentru ordin infinit.
\end{proof}

\begin{teorema}[ordinul unui produs de elemente care comută]\label{teo:ordprodus}
Fie $a,b\in G$ cu $ab=ba$, $\ord a=m$, $\ord b=n$, și fie $d=(m,n)$, $\ell=[m,n]$ (c.m.m.m.c.). Atunci
\[
\ord(ab)\ \Big|\ \ell
\qquad\text{și}\qquad
\frac{\ell}{d}\ \Big|\ \ord(ab).
\]
În particular, dacă $(m,n)=1$, atunci $\ord(ab)=mn$.
\end{teorema}

\begin{proof}
\emph{Prima divizibilitate}: $(ab)^{\ell}=a^{\ell}b^{\ell}=e$, deoarece $m\mid\ell$ și $n\mid\ell$ (aici este esențial $ab=ba$); apoi Teorema~\ref{teo:ordin}(i).

\emph{A doua}: fie $k=\ord(ab)$ și $c=a^k$. Din $(ab)^k=a^kb^k=e$ rezultă $c=a^k=b^{-k}$, deci $c\in\langle a\rangle\cap\langle b\rangle$. Atunci $c^m=(a^m)^k=e$ și $c^n=(b^n)^{-k}=e$, deci $\ord(c)$ divide și $m$, și $n$, prin urmare divide $d=(m,n)$; în particular $c^d=e$. Aceasta înseamnă $a^{kd}=e$ și $b^{kd}=e$, de unde $m\mid kd$ și $n\mid kd$, adică (simplificând cu $d$, care divide $m$ și $n$)
\[
\frac md\ \Big|\ k\qquad\text{și}\qquad\frac nd\ \Big|\ k.
\]
Observăm acum că $\bigl(\frac md,\frac nd\bigr)=1$: dacă un prim $r$ le-ar divide pe amândouă, atunci $rd\mid m$ și $rd\mid n$, deci $rd\mid(m,n)=d$ --- absurd. Un număr divizibil cu doi factori primi între ei este divizibil cu produsul lor, deci
\[
\frac md\cdot\frac nd=\frac{mn}{d^2}=\frac{\ell}{d}\ \Big|\ k.
\]
\emph{Cazul coprim}: pentru $d=1$ avem $\ell=mn$, iar cele două divizibilități dau $mn\mid\ord(ab)\mid mn$.
\end{proof}

\begin{corolar}[{element de ordin $[m,n]$}]\label{cor:lcm}
Dacă un grup \textbf{abelian} conține elemente de ordine $m$ și $n$, atunci conține un element de ordin $[m,n]$.
\end{corolar}

\begin{proof}
Scriem $\ell=[m,n]=p_1^{\alpha_1}\cdots p_r^{\alpha_r}$. Pentru fiecare $i$, puterea $p_i^{\alpha_i}$ divide $m$ sau $n$ (exponentul lui $p_i$ din $\ell$ este atins în unul dintre ele); dacă, de pildă, $p_i^{\alpha_i}\mid m$, atunci elementul $c_i=a^{m/p_i^{\alpha_i}}$ are ordinul $\frac{m}{(m,\,m/p_i^{\alpha_i})}=p_i^{\alpha_i}$ (Propoziția~\ref{prop:ordputere}). Obținem $c_1,\dots,c_r$ cu ordine $p_1^{\alpha_1},\dots,p_r^{\alpha_r}$, două câte două prime între ele; aplicând repetat cazul coprim din Teorema~\ref{teo:ordprodus}, produsul $c_1c_2\cdots c_r$ are ordinul $p_1^{\alpha_1}\cdots p_r^{\alpha_r}=\ell$.
\end{proof}

\begin{observatie}
(a) Comutativitatea este esențială: în $S_3$, transpozițiile $\sigma=(1\,2)$ și $\tau=(2\,3)$ au ordinul $2$, dar $\sigma\tau$ este un ciclu de lungime $3$, deci $\ord(\sigma\tau)=3\nmid[2,2]$. (b) Ambele margini din Teorema~\ref{teo:ordprodus} se ating: pentru $b=a^{-1}$ avem $\ord(ab)=\ord(e)=1=\ell/d$, iar în cazul coprim $\ord(ab)=\ell$. Între aceste margini, $\ord(ab)$ nu este determinat doar de $m$ și $n$.
\end{observatie}

\begin{propozitie}[intersecția de subgrupuri]\label{prop:intersectie}
Intersecția oricărei familii de subgrupuri ale lui $G$ este un subgrup al lui $G$.
\end{propozitie}

\begin{proof}
$e$ aparține fiecărui subgrup, deci și intersecției; dacă $x,y$ aparțin fiecărui subgrup din familie, atunci $xy^{-1}$ la fel, și aplicăm criteriul (Propoziția~\ref{prop:criteriu}).
\end{proof}

% --- grupuri mici + exponent, plasate la Ordinul unui element ---
\subsubsection*{Exponentul unui grup}

\begin{definitie}[exponentul grupului]
\emph{Exponentul} unui grup $G$, notat $\exp(G)$, este cel mai mic $n\ge1$ cu proprietatea $x^n=e$ pentru \textbf{orice} $x\in G$, dacă un astfel de $n$ există. Pentru $G$ finit, exponentul există (de exemplu $n=|G|$ convine, prin Corolarul~\ref{cor:lagrange}).
\end{definitie}

\begin{propozitie}[proprietățile exponentului]\label{prop:exponent}
Fie $G$ un grup finit. Atunci:
\emph{(i)} $\exp(G)$ este c.m.m.m.c.-ul ordinelor elementelor lui $G$ și, pentru $m\in\NN^*$: $x^m=e$ pentru orice $x$ $\iff\exp(G)\mid m$;
\emph{(ii)} $\exp(G)$ divide $|G|$;
\emph{(iii)} $\exp(G)$ și $|G|$ au \textbf{aceiași divizori primi};
\emph{(iv)} dacă $G$ este \textbf{abelian}, există un element de ordin $\exp(G)$; în consecință, $\exp(G)$ este cel mai mare ordin de element din $G$.
\end{propozitie}

\begin{proof}
(i) $x^m=e$ pentru toți $x$ $\iff\ord(x)\mid m$ pentru toți $x$ (Teorema~\ref{teo:ordin}(i)) $\iff L\mid m$, unde $L$ este c.m.m.m.c.-ul ordinelor; cel mai mic asemenea $m$ este chiar $L$.
(ii) $x^{|G|}=e$ pentru orice $x$ (Corolarul~\ref{cor:lagrange}), apoi (i).
(iii) Dacă un prim $q$ divide $\exp(G)$, atunci $q\mid|G|$ din (ii). Reciproc, dacă $p\mid|G|$, teorema lui Cauchy (Teorema~\ref{teo:cauchy}) dă un element de ordin $p$, deci $p\mid\exp(G)$ prin (i).
(iv) Enumerăm $G=\{x_1,\dots,x_k\}$ și aplicăm repetat Corolarul~\ref{cor:lcm}: obținem succesiv elemente de ordin $[\ord x_1,\ord x_2]$, apoi $[\ord x_1,\ord x_2,\ord x_3]$ etc., iar la final un element de ordin $L=\exp(G)$.
\end{proof}

\begin{corolar}[criteriu de ciclicitate]\label{cor:expciclic}
Un grup abelian finit $G$ este ciclic dacă și numai dacă $\exp(G)=|G|$.
\end{corolar}

\begin{proof}
Dacă $G=\langle a\rangle$, atunci $\exp(G)\ge\ord(a)=|G|$, iar $\exp(G)\mid|G|$ dă egalitatea. Reciproc, prin Propoziția~\ref{prop:exponent}(iv) există un element de ordin $\exp(G)=|G|$, care generează $G$.
\end{proof}

\begin{observatie}
Ipoteza de comutativitate la (iv) și la corolar este esențială: $\exp(S_3)=[1,2,3]=6=|S_3|$, dar $S_3$ nu are niciun element de ordin $6$ (nu este ciclic). Invers, chiar și în cazul abelian exponentul poate fi strict mai mic decât ordinul: grupul lui Klein are $\exp=2$ și $4$ elemente. Punctul (iv) este exact ideea din a doua demonstrație a Teoremei~\ref{teo:ciclicitate}: într-un subgrup finit al lui $(K^*,\cdot)$, ordinul maxim coincide cu exponentul.
\end{observatie}

\begin{definitie}[subgrupul generat de o mulțime]
Fie $(G,\cdot)$ un grup și $X\subseteq G$. \emph{Subgrupul generat de $X$} este
\[
\langle X\rangle=\bigcap\{H\mid X\subseteq H,\ H\le G\},
\]
intersecția tuturor subgrupurilor lui $G$ care conțin $X$.
\end{definitie}

\begin{propozitie}[descrierea subgrupului generat]
$\langle X\rangle$ este cel mai mic subgrup al lui $G$ (față de incluziune) care conține $X$, iar pentru $X\neq\varnothing$
\[
\langle X\rangle=\{x_1^{k_1}x_2^{k_2}\cdots x_n^{k_n}\mid n\in\NN^*,\ x_i\in X,\ k_i\in\ZZ\},
\]
mulțimea tuturor produselor finite de puteri întregi ale elementelor lui $X$. În particular $\langle x\rangle=\{x^k\mid k\in\ZZ\}$.
\end{propozitie}
\begin{proof}
Intersecția unei familii de subgrupuri este subgrup (conține $e$; e închisă la $xy^{-1}$), deci $\langle X\rangle\le G$, conține $X$ și este inclus în orice subgrup ce conține $X$ --- adică e cel mai mic. Fie $M$ mulțimea produselor finite din enunț. Orice subgrup ce conține $X$ conține și $M$ (închidere la produse și inverse), deci $M\subseteq\langle X\rangle$. Reciproc, $M$ este el însuși subgrup (produsul a două produse finite e produs finit; inversul lui $x_1^{k_1}\cdots x_n^{k_n}$ este $x_n^{-k_n}\cdots x_1^{-k_1}$) și conține $X$, deci $\langle X\rangle\subseteq M$.
\end{proof}

\section{Subgrupuri remarcabile}

Grupăm aici subgrupurile atașate canonic unui element sau unei submulțimi --- centralizator, centru, normalizator, subgrup normal --- împreună cu proprietățile lor de bază, folosite constant în restul fișei.

\begin{definitie}[centralizator; centru]
Pentru $a\in G$, \emph{centralizatorul} lui $a$ este
\[
C_G(a)=\{x\in G\mid xa=ax\};
\]
pentru o submulțime $M\subseteq G$, \emph{centralizatorul} lui $M$ este
$C_G(M)=\{x\in G\mid xm=mx,\ \forall m\in M\}=\bigcap_{m\in M}C_G(m)$.
\emph{Centrul} grupului este $Z(G)=C_G(G)=\{a\in G\mid ax=xa,\ \forall x\in G\}$.
\end{definitie}

\begin{propozitie}\label{prop:centralizator}
\emph{(i)} $C_G(a)\le G$ pentru orice $a$; în consecință $C_G(M)\le G$ și $Z(G)=\bigcap_{a\in G}C_G(a)\le G$.
\emph{(ii)} $\langle a\rangle\le C_G(a)$, iar $a\in Z(G)\iff C_G(a)=G$.
\emph{(iii)} $G$ este abelian $\iff Z(G)=G$.
\end{propozitie}

\begin{proof}
(i) $e\in C_G(a)$. Fie $x,y\in C_G(a)$. Din $ya=ay$, înmulțind la stânga și la dreapta cu $y^{-1}$, obținem $ay^{-1}=y^{-1}a$, deci
\[
(xy^{-1})a=x(y^{-1}a)=x(ay^{-1})=(xa)y^{-1}=(ax)y^{-1}=a(xy^{-1}),
\]
adică $xy^{-1}\in C_G(a)$; aplicăm criteriul. Restul rezultă din Propoziția~\ref{prop:intersectie}. (ii) și (iii): direct din definiții ($a$ comută cu propriile puteri).
\end{proof}

\begin{definitie}[normalizator]
Pentru $x\in G$ și $M\subseteq G$ notăm $xMx^{-1}=\{xmx^{-1}\mid m\in M\}$ (\emph{conjugata} lui $M$). \emph{Normalizatorul} lui $M$ în $G$ este
\[
N_G(M)=\{x\in G\mid xMx^{-1}=M\},
\]
echivalent, mulțimea acelor $x$ cu $xM=Mx$ (egalitate de mulțimi).
\end{definitie}

\begin{propozitie}\label{prop:normalizator}
\emph{(i)} $N_G(M)\le G$ pentru orice submulțime $M$.
\emph{(ii)} Dacă $H\le G$, atunci $H\le N_G(H)$ și $C_G(H)\le N_G(H)$.
\end{propozitie}

\begin{proof}
(i) $e\in N_G(M)$. Dacă $xMx^{-1}=M$, aplicând $x^{-1}\,(\cdot)\,x$ ambilor membri obținem $M=x^{-1}Mx$, deci $x^{-1}\in N_G(M)$. Pentru închidere:
\[
(xy)M(xy)^{-1}=x\,(yMy^{-1})\,x^{-1}=xMx^{-1}=M.
\]
(ii) Pentru $h\in H$: $hHh^{-1}\subseteq H$ (închiderea lui $H$); aplicând același lucru lui $h^{-1}$ rezultă $h^{-1}Hh\subseteq H$, echivalent $H\subseteq hHh^{-1}$, deci egalitate și $h\in N_G(H)$. Dacă $x\in C_G(H)$, atunci $xhx^{-1}=h$ pentru orice $h\in H$, deci $xHx^{-1}=H$.
\end{proof}

\begin{observatie}
Dacă $M$ este \textbf{finită}, incluziunea $xMx^{-1}\subseteq M$ implică deja egalitatea (funcția $m\mapsto xmx^{-1}$ este injectivă de la $M$ în $M$, deci surjectivă); pentru mulțimi infinite o singură incluziune nu este suficientă.
\end{observatie}

\begin{definitie}[subgrup normal]\label{def:normal}
Un subgrup $H\le G$ se numește \emph{normal} (notăm $H\trianglelefteq G$) dacă $xHx^{-1}=H$ pentru orice $x\in G$; echivalent, $N_G(H)=G$; echivalent, $xH=Hx$ pentru orice $x\in G$. Astfel, $N_G(H)$ este cel mai mare subgrup al lui $G$ în care $H$ este normal.
\end{definitie}

\begin{propozitie}[criteriul de normalitate]\label{prop:normalcrit}
$H\trianglelefteq G$ dacă și numai dacă $xHx^{-1}\subseteq H$ pentru \textbf{orice} $x\in G$ (este suficientă o singură incluziune, cerută însă pentru toți $x$).
\end{propozitie}

\begin{proof}
Aplicând incluziunea și lui $x^{-1}$ obținem $x^{-1}Hx\subseteq H$, echivalent $H\subseteq xHx^{-1}$; deci egalitate.
\end{proof}

\begin{exemplu}[subgrupuri normale și nenormale]\label{ex:normal}
\begin{enumerate}[label=(\alph*),nosep]
\item Într-un grup abelian orice subgrup este normal ($xhx^{-1}=h$).
\item $Z(G)\trianglelefteq G$ și, mai general, orice subgrup al centrului este normal.
\item Orice subgrup de \textbf{indice 2} este normal: pentru $x\notin H$, clasa la stânga $xH$ și clasa la dreapta $Hx$ sunt amândouă egale cu $G\setminus H$ (clasele la dreapta partiționează $G$ din motive simetrice cu Lema~\ref{lema:clase}), deci $xH=Hx$; pentru $x\in H$ egalitatea este evidentă.
\item În $S_3$, subgrupul $H=\{e,(1\,2)\}$ \textbf{nu} este normal: $(1\,3)(1\,2)(1\,3)^{-1}=(2\,3)\notin H$.
\item Nucleul oricărui morfism de grupuri este subgrup normal (Teorema~\ref{teo:nucleu}).
\end{enumerate}
\end{exemplu}

\begin{lema}[Bézout pentru exponenți și subgrupuri]\label{lema:bezoutexp}
Fie $H\le G$, $x\in G$ și $n,p\in\NN^*$. Dacă $x^n\in H$ și $x^p\in H$, atunci $x^{(n,p)}\in H$. În particular, dacă $(n,p)=1$, atunci $x\in H$.
\end{lema}

\begin{proof}
Fie $d=(n,p)$; relația lui Bézout dă $u,v\in\ZZ$ cu $un+vp=d$. Atunci
\[
x^{d}=x^{un+vp}=(x^n)^u\,(x^p)^v\in H,
\]
deoarece $H$ este închis la puteri întregi (inclusiv negative) și la produse.
\end{proof}

\begin{observatie}
Lema se aplică de obicei cu $H=Z(G)$ (sau cu un centralizator): dacă $x^n\in Z(G)$ și $x^p\in Z(G)$, atunci $x^{(n,p)}\in Z(G)$; în particular, pentru $(n,p)=1$ rezultă $x\in Z(G)$. O aplicație tipică este Corolarul~\ref{cor:nn1}.
\end{observatie}

\begin{exemplu}
În $(\ZZ_n,+)$, ordinul lui $\cl{k}$ este $\dfrac{n}{(n,k)}$ (versiunea aditivă a Propoziției~\ref{prop:ordputere}, cu $\cl k=k\cdot\cl 1$).
\end{exemplu}

\section{Morfisme și izomorfisme de grupuri}\label{sec:morfisme}

\begin{definitie}
Fie $(G,\cdot)$ și $(G',*)$ grupuri. O funcție $f:G\to G'$ este \emph{morfism} dacă
$f(xy)=f(x)*f(y)$ pentru orice $x,y\in G$. Un morfism bijectiv este \emph{izomorfism} (scriem $G\simeq G'$). Un morfism $G\to G$ este \emph{endomorfism}; un izomorfism $G\to G$ este \emph{automorfism}.
\end{definitie}

\begin{propozitie}\label{prop:morfism}
Dacă $f:G\to G'$ este morfism, atunci $f(e)=e'$, $f(x^{-1})=f(x)^{-1}$ și $f(x^n)=f(x)^n$ pentru orice $n\in\ZZ$.
\end{propozitie}

\begin{proof}
$f(e)=f(e\cdot e)=f(e)*f(e)$ și simplificăm în $G'$. Apoi $f(x)*f(x^{-1})=f(xx^{-1})=e'$, deci $f(x^{-1})=f(x)^{-1}$ prin unicitatea inversului. Ultima afirmație: inducție pentru $n\ge0$, iar pentru $n<0$ se combină cu precedenta.
\end{proof}

\begin{propozitie}\label{prop:compunere}
Compunerea a două morfisme este morfism, iar inversa unui izomorfism este izomorfism (deci $\simeq$ este relație de echivalență).
\end{propozitie}

\begin{proof}
Prima parte, calcul direct. Pentru a doua: fie $u=f^{-1}(a)$, $v=f^{-1}(b)$; din $f(uv)=f(u)*f(v)=a*b$ rezultă $f^{-1}(a*b)=uv=f^{-1}(a)f^{-1}(b)$.
\end{proof}

\begin{definitie}
\emph{Nucleul} și \emph{imaginea} morfismului $f:G\to G'$ sunt
$\Ker f=\{x\in G\mid f(x)=e'\}$ și $\Ima f=f(G)$.
\end{definitie}

\begin{propozitie}[imagini directe și inverse]\label{prop:imagini}
Fie $f:G\to G'$ un morfism. \emph{(i)} Pentru orice subgrup $K\le G$, imaginea $f(K)=\{f(x)\mid x\in K\}$ este subgrup al lui $G'$; în particular $\Ima f=f(G)\le G'$. \emph{(ii)} Pentru orice subgrup $K'\le G'$, preimaginea $f^{-1}(K')=\{x\in G\mid f(x)\in K'\}$ este subgrup al lui $G$.
\end{propozitie}

\begin{proof}
(i) $e'=f(e)\in f(K)$; pentru $x,y\in K$, $f(x)f(y)^{-1}=f(xy^{-1})\in f(K)$ (Propoziția~\ref{prop:morfism}), și aplicăm criteriul de subgrup. (ii) Analog: $e\in f^{-1}(K')$, iar din $f(x),f(y)\in K'$ rezultă $f(xy^{-1})=f(x)f(y)^{-1}\in K'$.
\end{proof}

\begin{teorema}[nucleul]\label{teo:nucleu}
Fie $f:G\to G'$ un morfism. Atunci:
\emph{(i)} $\Ker f$ este subgrup \textbf{normal} al lui $G$;\quad
\emph{(ii)} $f$ este injectiv $\iff\Ker f=\{e\}$.
\end{teorema}

\begin{proof}
(i) $\Ker f=f^{-1}(\{e'\})$ este subgrup (Propoziția~\ref{prop:imagini}(ii)). Normalitatea: pentru $x\in G$ și $k\in\Ker f$,
\[
f(xkx^{-1})=f(x)\,f(k)\,f(x)^{-1}=f(x)\,e'\,f(x)^{-1}=e',
\]
deci $x(\Ker f)x^{-1}\subseteq\Ker f$ pentru orice $x$, iar Propoziția~\ref{prop:normalcrit} încheie.
(ii) Dacă $f$ este injectiv și $f(x)=e'=f(e)$, atunci $x=e$. Reciproc, dacă $\Ker f=\{e\}$ și $f(x)=f(y)$, atunci $f(xy^{-1})=e'$, deci $xy^{-1}=e$, adică $x=y$.
\end{proof}

\begin{propozitie}[morfismele și ordinul]\label{prop:morford}
Dacă $f:G\to G'$ este morfism și $x$ are ordin finit, atunci $\ord f(x)\mid\ord x$. Dacă $f$ este injectiv, $\ord f(x)=\ord x$ (inclusiv în cazul infinit).
\end{propozitie}

\begin{proof}
$f(x)^{\ord x}=f(x^{\ord x})=e'$, deci $\ord f(x)\mid \ord x$ (Teorema~\ref{teo:ordin}(i)). Dacă $f$ este injectiv și $f(x)^m=e'=f(e)$, atunci $x^m=e$, deci $\ord x\mid m$; luând $m=\ord f(x)$ obținem egalitatea. În cazul infinit, dacă $f(x)$ ar avea ordin finit $m$, ar rezulta $x^m=e$, contradicție.
\end{proof}

\begin{observatie}[invarianți la izomorfism]\label{obs:invarianti}
Un izomorfism păstrează: cardinalul, comutativitatea, ciclicitatea, ordinul fiecărui element (Propoziția~\ref{prop:morford}), numărul soluțiilor ecuației $x^k=e$ și numărul elementelor de un ordin dat. \textbf{Metoda standard} de a arăta că $G\not\simeq G'$: găsiți un astfel de invariant care diferă.
\end{observatie}

\begin{tehbox}{De ce ,,izomorf`` înseamnă ,,aceeași structură``}
Un izomorfism $f:G\to G'$ este, în esență, o \emph{redenumire} a elementelor: fiecare $x\in G$ primește numele nou $f(x)\in G'$, iar condiția de morfism $f(xy)=f(x)\,f(y)$ spune că tabla operației se copiază întocmai --- compui întâi și traduci, sau traduci întâi și compui: obții același lucru. De aceea \textbf{orice proprietate formulată exclusiv în limbajul operației de grup se transportă prin $f$}; orice teoremă despre operația lui $G$ devine, cuvânt cu cuvânt, o teoremă despre $G'$. Concret, dicționarul:
\begin{enumerate}[leftmargin=*, itemsep=1pt]
\item \textbf{Elementul neutru merge în elementul neutru.} Din $f(e)=f(e\cdot e)=f(e)f(e)$, simplificând în $G'$, rezultă $f(e_G)=e_{G'}$. ,,Rolul`` de neutru nu ține de numele elementului, ci de comportarea lui la operație --- iar comportarea se traduce.
\item \textbf{Inversele rămân perechi.} $f(x)\,f(x^{-1})=f(xx^{-1})=e_{G'}$, deci $f(x^{-1})=f(x)^{-1}$: perechile element--invers se corespund exact.
\item \textbf{Puterile și ordinele se păstrează.} Prin inducție $f(x^n)=f(x)^n$ pentru orice $n\in\ZZ$; cum $f$ e injectivă, $x^n=e\iff f(x)^n=e'$, deci $\ord(f(x))=\ord(x)$ (Propoziția~\ref{prop:morford}). Fiind și bijecție, $G$ și $G'$ au \emph{exact același număr de elemente de fiecare ordin} --- același ,,spectru`` de ordine.
\item \textbf{Ecuațiile au același număr de soluții.} $x$ verifică $x^k=a$ $\iff$ $f(x)$ verifică $y^k=f(a)$: bijecția $x\mapsto f(x)$ duce mulțimea soluțiilor unei ecuații exact în mulțimea soluțiilor ecuației traduse. În particular $\#\{x\in G:x^k=e\}=\#\{y\in G':y^k=e'\}$ pentru orice $k$.
\item \textbf{Subgrupurile se corespund unu-la-unu.} Imaginea unui subgrup e subgrup (Propoziția~\ref{prop:imagini}), iar $H\mapsto f(H)$ este o \emph{bijecție} între subgrupurile lui $G$ și cele ale lui $G'$, cu inversa $K\mapsto f^{-1}(K)$. Ea păstrează: incluziunile, cardinalul ($|f(H)|=|H|$), indicele ($f$ duce clasele $xH$ în clasele $f(x)f(H)$), \emph{normalitatea} ($H\trianglelefteq G\iff f(H)\trianglelefteq G'$, căci $f(gxg^{-1})=f(g)f(x)f(g)^{-1}$) și ciclicitatea, cu $f(\langle x\rangle)=\langle f(x)\rangle$. Întreaga ,,hartă`` a subgrupurilor este deci aceeași.
\item \textbf{Obiectele canonice se corespund.} $f(Z(G))=Z(G')$ și $f\big(C_G(a)\big)=C_{G'}\big(f(a)\big)$ (din $ax=xa\iff f(a)f(x)=f(x)f(a)$, surjectivitatea acoperind tot $G'$); conjugarea se transportă prin $f(gxg^{-1})=f(g)\,f(x)\,f(g)^{-1}$, deci clasele de conjugare merg în clase de conjugare de aceleași cardinale; comutatorii merg în comutatori.
\item \textbf{Proprietățile globale se moștenesc.} $G$ abelian $\iff G'$ abelian; $G$ ciclic $\iff G'$ ciclic, generatorii corespunzându-se ($G=\langle x\rangle\iff G'=\langle f(x)\rangle$) --- în particular cele două grupuri au același număr de generatori.
\item \textbf{Pentru grupuri finite: aceeași tablă.} Scriind tabla lui Cayley a lui $G$ și înlocuind fiecare element prin numele lui nou $f(x)$, obținem exact tabla lui $G'$: două grupuri finite sunt izomorfe dacă și numai dacă tablele lor coincid după o renumerotare a elementelor.
\end{enumerate}
\textbf{Ce NU vede izomorfismul:} natura concretă a elementelor (numere, matrice, permutări, clase de resturi) și notația operației (aditivă sau multiplicativă). $(\RR,+)$ și $((0,\infty),\cdot)$ sunt ,,același grup``, deși elementele lor arată complet diferit --- logaritmul e doar dicționarul dintre ele. \textbf{Atenție:} o simplă \emph{bijecție} nu transportă structura --- $\ZZ_4$ și grupul lui Klein au ambele $4$ elemente, dar nu sunt izomorfe ($\ZZ_4$ are un element de ordin $4$, Klein nu); toată forța stă în condiția de morfism. \textbf{Morala de concurs:} pentru $G\simeq G'$ construiești explicit dicționarul $f$ și verifici morfism $+$ bijecție (sau transport de structură); pentru $G\not\simeq G'$ cauți în lista de mai sus un invariant care diferă.
\end{tehbox}

\begin{exemplu}[izomorfisme și non-izomorfisme]\label{ex:noniso}
\begin{enumerate}[label=(\alph*),nosep]
\item $\exp:(\RR,+)\to((0,\infty),\cdot)$ este izomorfism: $e^{x+y}=e^xe^y$, bijectivă cu inversa $\ln$.
\item Pentru $G=\RR\setminus\{1\}$ cu $x\circ y=x+y-xy$ (Exemplul~\ref{ex:grupuri}(g)), funcția $f:G\to\RR^*$, $f(x)=1-x$, este izomorfism: $f(x\circ y)=1-x-y+xy=(1-x)(1-y)=f(x)f(y)$, bijectivă. (\emph{Transport de structură} --- tipar frecvent în manuale și la olimpiadă.)
\item $(\RR^*,\cdot)\not\simeq(\CC^*,\cdot)$: ecuația $x^4=1$ are $2$ soluții în $\RR^*$ și $4$ în $\CC^*$.
\item $(\RR,+)\not\simeq(\RR^*,\cdot)$: în $\RR^*$, elementul $-1$ are ordinul $2$, dar în $(\RR,+)$ singurul element de ordin finit este $0$.
\item $(\QQ,+)\not\simeq(\ZZ,+)$: $\ZZ$ este ciclic, $\QQ$ nu. Într-adevăr, dacă $\QQ=\langle a/b\rangle$ cu $a\neq0$, atunci $\frac{a}{2b}=k\cdot\frac ab$ ar da $k=\frac12\notin\ZZ$.
\end{enumerate}
\end{exemplu}

\begin{propozitie}[transport de structură]\label{prop:transport}
Fie $(G,\cdot)$ un grup, $M$ o mulțime nevidă și $\varphi:M\to G$ o funcție \textbf{bijectivă}. Definim pe $M$ legea
\[
x*y=\varphi^{-1}\bigl(\varphi(x)\cdot\varphi(y)\bigr).
\]
Atunci $(M,*)$ este grup, iar $\varphi:(M,*)\to(G,\cdot)$ este izomorfism. Elementul neutru este $\varphi^{-1}(e)$, simetricul lui $x$ este $\varphi^{-1}\bigl(\varphi(x)^{-1}\bigr)$, iar dacă $G$ este abelian, $(M,*)$ este abelian.
\end{propozitie}

\begin{proof}
Prin construcție, $\varphi(x*y)=\varphi(x)\varphi(y)$. Asociativitatea:
\[
\varphi\bigl(x*(y*z)\bigr)=\varphi(x)\bigl(\varphi(y)\varphi(z)\bigr)=\bigl(\varphi(x)\varphi(y)\bigr)\varphi(z)=\varphi\bigl((x*y)*z\bigr),
\]
și $\varphi$ este injectivă. Neutrul: $x*\varphi^{-1}(e)=\varphi^{-1}\bigl(\varphi(x)\,e\bigr)=x$ și simetric la stânga; simetricul: $x*\varphi^{-1}\bigl(\varphi(x)^{-1}\bigr)=\varphi^{-1}(e)$. Așadar $(M,*)$ este grup, iar $\varphi$ este morfism bijectiv.
\end{proof}

\begin{exemplu}[legile ,,de manual``]\label{ex:transport}
Legile de pe $\RR$ de forma $x*y=xy-m(x+y)+m^2+m$ sunt transporturi prin funcții afine: pentru $\varphi(x)=x-m$ avem
\[
\begin{gathered}
\varphi(x)\varphi(y)=(x-m)(y-m)=xy-m(x+y)+m^2,\\
\text{deci}\qquad \varphi^{-1}\bigl(\varphi(x)\varphi(y)\bigr)=x*y.
\end{gathered}
\]
Prin urmare $\bigl(\RR\setminus\{m\},*\bigr)\simeq(\RR^*,\cdot)$, cu neutrul $\varphi^{-1}(1)=m+1$ și simetricul lui $x$ egal cu $m+\frac{1}{x-m}$; toate proprietățile (asociativitate, comutativitate, elemente simetrizabile) vin \emph{gratuit} din Propoziția~\ref{prop:transport}. Exemplul~\ref{ex:noniso}(b) este tot un transport, prin $\varphi(x)=1-x$. \textbf{Strategia de concurs}: dată o lege ,,inventată``, căutați $\varphi$ afină, $\varphi(x)=\alpha x+\beta$, cu $\varphi(x*y)=\varphi(x)\varphi(y)$ (identificare de coeficienți); analog pentru transport din $(\RR,+)$.
\end{exemplu}

\begin{observatie}[endomorfismele lui $(\ZZ_n,+)$ și $(\QQ,+)$]\label{obs:endo}
(a) Orice endomorfism $f$ al lui $(\ZZ_n,+)$ este de forma $f(\cl k)=\cl k\cdot\cl a$, unde $\cl a=f(\cl 1)$ (din aditivitate, $f(\cl k)=k\,f(\cl 1)$); $f$ este automorfism dacă și numai dacă $\cl a$ este inversabil în inelul $\ZZ_n$, adică $(a,n)=1$.
(b) Orice endomorfism $f$ al lui $(\QQ,+)$ este $f(x)=cx$ cu $c=f(1)$: din aditivitate $f(nx)=nf(x)$ pentru $n\in\ZZ$, apoi $q\,f\!\left(\frac pq\right)=f(p)=p\,f(1)$, deci $f\!\left(\frac pq\right)=\frac pq\,f(1)$. (Aceeași schemă rezolvă ecuația funcțională Cauchy pe $\QQ$.)
\end{observatie}

\begin{definitie}
Pentru grupurile $G$ și $H$ notăm $\Hom(G,H)$ mulțimea morfismelor de la $G$ la $H$, $\End(G)=\Hom(G,G)$ mulțimea \emph{endomorfismelor} și $\Aut(G)$ mulțimea \emph{automorfismelor} lui $G$.
\end{definitie}

\begin{propozitie}[{structura lui $\Hom(G,H)$}]\label{prop:homgrup}
Dacă $H$ este \textbf{abelian}, atunci $\Hom(G,H)$ este grup abelian față de \emph{produsul punctual}
\[
(f\cdot g)(x)=f(x)\,g(x)
\]
(atenție: nu compunerea!). Elementul neutru este morfismul trivial $x\mapsto e_H$, iar simetricul lui $f$ este $x\mapsto f(x)^{-1}$.
\end{propozitie}

\begin{proof}
Închiderea: $(f\cdot g)(xy)=f(x)f(y)\,g(x)g(y)=f(x)g(x)\,f(y)g(y)$, unde comutativitatea lui $H$ a permis mutarea lui $f(y)$ peste $g(x)$; deci $f\cdot g$ este morfism. Funcția $\bar f(x)=f(x)^{-1}$ este morfism: $\bar f(xy)=\bigl(f(x)f(y)\bigr)^{-1}=f(y)^{-1}f(x)^{-1}=f(x)^{-1}f(y)^{-1}$, din nou prin comutativitate. Asociativitatea, comutativitatea și proprietatea neutrului se verifică punctual, fiind moștenite din $H$.
\end{proof}

\begin{observatie}
Fără comutativitatea lui $H$, produsul punctual a două morfisme poate să nu mai fie morfism: pentru $f=g=1_{S_3}$, produsul punctual este $x\mapsto x^2$, care nu este endomorfism al lui $S_3$ (Propoziția~\ref{prop:patrat}, $S_3$ nefiind abelian).
\end{observatie}

\begin{propozitie}[{$\Hom(\ZZ,H)$}]\label{prop:homZ}
Pentru orice grup $H$, funcția $f\mapsto f(1)$ este o \textbf{bijecție} de la $\Hom(\ZZ,H)$ la $H$: orice morfism de la $(\ZZ,+)$ la $H$ este de forma $f_a(k)=a^k$, pentru un unic $a\in H$. Dacă $H$ este abelian, aceasta este izomorfism de grupuri:
\[
\Hom(\ZZ,H)\simeq H.
\]
\end{propozitie}

\begin{proof}
Orice morfism satisface $f(k)=f(k\cdot 1)=f(1)^k$ (Propoziția~\ref{prop:morfism}), deci este determinat de $a=f(1)$; reciproc, pentru orice $a\in H$, $f_a(k)=a^k$ este morfism, din $a^{k+l}=a^ka^l$. Așadar $f\mapsto f(1)$ este bijecție. Dacă $H$ este abelian, $(f\cdot g)(1)=f(1)g(1)$, deci bijecția respectă operațiile.
\end{proof}

\begin{teorema}[{$\Hom(\ZZ_m,\ZZ_n)$}]\label{teo:homZmZn}
Fie $d=(m,n)$. Atunci
\[
\Hom(\ZZ_m,\ZZ_n)\simeq\ZZ_{d}.
\]
Mai precis: morfismele sunt exact $f_{\cl a}(\cl x)=x\cdot\cl a$, unde $\cl a$ parcurge subgrupul
$K=\{\cl a\in\ZZ_n\mid m\cdot\cl a=\cl 0\}=\bigl\langle\cl{n/d}\bigr\rangle$, ciclic de ordin $d$, iar $f\mapsto f(\cl 1)$ este izomorfism de la $\Hom(\ZZ_m,\ZZ_n)$ (cu adunarea punctuală) la $K$.
\end{teorema}

\begin{proof}
\emph{Determinare}: orice morfism satisface $f(\cl x)=x\cdot f(\cl 1)$ (aditivitate). \emph{Buna definire}: valoarea $\cl a=f(\cl 1)$ trebuie să satisfacă $m\cl a=f(\cl m)=f(\cl 0)=\cl 0$; reciproc, dacă $m\cl a=\cl 0$, formula $f_{\cl a}(\cl x)=x\cl a$ este bine definită (pentru $x\equiv x'\pmod m$, diferența $(x-x')\cl a$ este multiplu de $m\cl a=\cl 0$) și este evident morfism. Așadar $f\mapsto f(\cl 1)$ este bijecție de la $\Hom(\ZZ_m,\ZZ_n)$ la $K$, și respectă adunarea punctuală: $(f+g)(\cl 1)=f(\cl 1)+g(\cl 1)$.

\emph{Structura lui $K$}: condiția $m\cl a=\cl 0$ este ecuația ,,$y^m=e$`` în notație aditivă, care în grupul ciclic $(\ZZ_n,+)$ are exact $(m,n)=d$ soluții (Propoziția~\ref{prop:solutii}). $K$ este subgrup (verificare directă sau Propoziția~\ref{prop:imagini}(ii)) al grupului ciclic $\ZZ_n$, deci ciclic (Teorema~\ref{teo:subgrupciclic}); având ordinul $d$, este unicul subgrup de ordin $d$, anume $\langle\cl{n/d}\rangle$, și $K\simeq\ZZ_d$ (Teorema~\ref{teo:clasif}).
\end{proof}

\begin{corolar}
$\Hom(\ZZ_m,\ZZ_n)$ conține doar morfismul nul $\iff(m,n)=1$. Pentru $m=n$ regăsim $\End(\ZZ_n)$ cu $n$ elemente, în acord cu Observația~\ref{obs:endo}(a).
\end{corolar}

\begin{propozitie}[{$\End(G)$ ca monoid; $\Aut(G)=U(\End(G))$}]\label{prop:endmonoid}
$(\End(G),\circ)$ este monoid (în general necomutativ), cu elementul neutru $1_G$, iar
\[
\Aut(G)=U\bigl(\End(G)\bigr):
\]
automorfismele sunt exact elementele inversabile ale acestui monoid. În particular, $(\Aut(G),\circ)$ este grup.
\end{propozitie}

\begin{proof}
Compunerea a două endomorfisme este endomorfism (Propoziția~\ref{prop:compunere}), este asociativă și are neutrul $1_G$: monoid. Dacă $f\in\Aut(G)$, inversa $f^{-1}$ este morfism (tot Propoziția~\ref{prop:compunere}), deci $f^{-1}\in\End(G)$ și $f$ este inversabil în monoid. Reciproc, dacă $f\circ g=g\circ f=1_G$, atunci $f$ este bijectivă, deci automorfism. În fine, elementele inversabile ale oricărui monoid formează grup față de operația indusă --- aceeași verificare ca la grupul unităților unui inel: $(f\circ g)^{-1}=g^{-1}\circ f^{-1}$.
\end{proof}

\begin{definitie}[automorfisme interioare]
Pentru $g\in G$, funcția $i_g:G\to G$, $i_g(x)=gxg^{-1}$, se numește \emph{automorfismul interior} asociat lui $g$. Notăm $\Inn(G)=\{i_g\mid g\in G\}$.
\end{definitie}

\begin{teorema}[{$\Inn(G)\trianglelefteq\Aut(G)$}]\label{teo:inn}
\emph{(i)} Fiecare $i_g$ este automorfism al lui $G$, iar $\Phi:G\to\Aut(G)$, $\Phi(g)=i_g$, este morfism de grupuri cu
\[
\Ima\Phi=\Inn(G)\qquad\text{și}\qquad\Ker\Phi=Z(G);
\]
în particular $\Inn(G)\le\Aut(G)$ și $i_g=i_h\iff h^{-1}g\in Z(G)$.
\emph{(ii)} Pentru orice $f\in\Aut(G)$ și orice $g\in G$,
\[
f\circ i_g\circ f^{-1}=i_{f(g)}.
\]
În consecință, $\Inn(G)$ este subgrup \textbf{normal} al lui $\Aut(G)$.
\end{teorema}

\begin{proof}
(i) $i_g(xy)=gxyg^{-1}=(gxg^{-1})(gyg^{-1})=i_g(x)i_g(y)$, deci $i_g$ este endomorfism; calcul direct: $(i_g\circ i_h)(x)=g(hxh^{-1})g^{-1}=(gh)x(gh)^{-1}$, adică $i_g\circ i_h=i_{gh}$, deci $\Phi$ este morfism, iar $i_g\circ i_{g^{-1}}=i_e=1_G$ arată că $i_g$ este bijectiv (automorfism). Imaginea unui morfism este subgrup (Propoziția~\ref{prop:imagini}), deci $\Inn(G)\le\Aut(G)$. Nucleul: $i_g=1_G\iff gxg^{-1}=x$ pentru orice $x\iff g\in Z(G)$; apoi $i_g=i_h\iff i_{h^{-1}g}=1_G\iff h^{-1}g\in Z(G)$.

(ii) Pentru orice $x\in G$:
\[
\bigl(f\circ i_g\circ f^{-1}\bigr)(x)=f\bigl(g\,f^{-1}(x)\,g^{-1}\bigr)=f(g)\,x\,f(g)^{-1}=i_{f(g)}(x).
\]
Deci $f\,\Inn(G)\,f^{-1}\subseteq\Inn(G)$ pentru orice $f\in\Aut(G)$, iar criteriul de normalitate (Propoziția~\ref{prop:normalcrit}) încheie.
\end{proof}

\begin{lema}[centru trivial]\label{lema:centrutrivial}
Dacă $Z(G)=\{e\}$, atunci $Z\bigl(\Aut(G)\bigr)=\{1_G\}$.
\end{lema}

\begin{proof}
Fie $f\in Z(\Aut(G))$. În particular, $f$ comută cu fiecare automorfism interior: $f\circ i_g=i_g\circ f$, adică $f\circ i_g\circ f^{-1}=i_g$. Dar din Teorema~\ref{teo:inn}(ii), $f\circ i_g\circ f^{-1}=i_{f(g)}$; deci $i_{f(g)}=i_g$, ceea ce înseamnă $g^{-1}f(g)\in Z(G)=\{e\}$, adică $f(g)=g$, pentru \textbf{orice} $g\in G$. Așadar $f=1_G$.
\end{proof}

\begin{observatie}
(a) Din $\Ker\Phi=Z(G)$ și Teorema~\ref{teo:nucleu}: dacă $Z(G)=\{e\}$, morfismul $g\mapsto i_g$ este injectiv, deci $G\simeq\Inn(G)\le\Aut(G)$ --- grupurile cu centru trivial se ,,scufundă`` în propriul grup de automorfisme. (b) Pentru grupuri abeliene $\Inn(G)=\{1_G\}$, iar automorfismele interesante sunt cele ,,exterioare``: de pildă $\Aut(\ZZ_n,+)\simeq\bigl(U(\ZZ_n),\cdot\bigr)$, prin $\cl a\mapsto f_{\cl a}$, cu $f_{\cl a}\circ f_{\cl b}=f_{\cl a\cl b}$ (Observația~\ref{obs:endo}(a)); iar din Propoziția~\ref{prop:homZ}, $\End(\ZZ,+)$ corespunde lui $\ZZ$ ($f_a(k)=ak$) și $\Aut(\ZZ,+)=\{\pm1_\ZZ\}\simeq\ZZ_2$.
\end{observatie}

\subsection*{Comutativitate prin aplicații de tip putere}

\begin{propozitie}\label{prop:patrat}
\emph{(i)} $f:G\to G$, $f(x)=x^2$, este endomorfism dacă și numai dacă $G$ este abelian. \emph{(ii)} La fel pentru $f(x)=x^{-1}$.
\end{propozitie}

\begin{proof}
(i) $f(xy)=f(x)f(y)\iff xyxy=x^2y^2\iff yx=xy$ (simplificări la stânga cu $x$ și la dreapta cu $y$). (ii) $(xy)^{-1}=x^{-1}y^{-1}\iff y^{-1}x^{-1}=x^{-1}y^{-1}$, adică inversele comută, echivalent cu comutativitatea.
\end{proof}

\begin{teorema}[exponenți consecutivi]\label{teo:consec}
Dacă există $n\in\NN^*$ astfel încât $(xy)^k=x^k y^k$ pentru $k\in\{n,n+1,n+2\}$ și pentru orice $x,y\in G$, atunci $G$ este abelian.
\end{teorema}

\begin{proof}
Din $(xy)^{n+1}=(xy)^n(xy)$ obținem $x^{n+1}y^{n+1}=x^ny^nxy$; înmulțind la stânga cu $x^{-n}$ și la dreapta cu $y^{-1}$ rezultă
\[
xy^n=y^nx\qquad\text{pentru orice }x,y\in G,
\]
adică $y^n\in Z(G)$ pentru orice $y$. Aplicând același raționament perechii $(n+1,n+2)$, obținem $y^{n+1}\in Z(G)$ pentru orice $y$. Atunci
\[
(xy)\,y^n=x\,y^{n+1}=y^{n+1}x=y\,(y^n x)=y\,(x y^n)=(yx)\,y^n,
\]
și simplificând la dreapta cu $y^n$ rezultă $xy=yx$.
\end{proof}

\begin{teorema}[endomorfismul $x\mapsto x^n$]\label{teo:endon}
Fie $n\ge2$ și presupunem că $f:G\to G$, $f(x)=x^n$, este endomorfism, adică $(xy)^n=x^ny^n$ pentru orice $x,y\in G$. Atunci:
\emph{(i)} $(yx)^{n-1}=x^{n-1}y^{n-1}$ pentru orice $x,y\in G$;
\emph{(ii)} dacă, în plus, $f$ este \textbf{surjectiv}, atunci $x^{n-1}\in Z(G)$ pentru orice $x\in G$.
\end{teorema}

\begin{proof}
(i) Scriind produsul alternat și grupând altfel factorii,
\[
(xy)^n=\underbrace{(xy)(xy)\cdots(xy)}_{n}=x\,\underbrace{(yx)(yx)\cdots(yx)}_{n-1}\,y=x\,(yx)^{n-1}y.
\]
Egalând cu $x^ny^n$ și simplificând la stânga cu $x$, la dreapta cu $y$: $(yx)^{n-1}=x^{n-1}y^{n-1}$.

(ii) Schimbând $x$ cu $y$ în (i): $(xy)^{n-1}=y^{n-1}x^{n-1}$. Atunci
\[
x^ny^n=(xy)^n=(xy)\,(xy)^{n-1}=xy\cdot y^{n-1}x^{n-1}=x\,y^n\,x^{n-1},
\]
și simplificând la stânga cu $x$ obținem
\[
x^{n-1}y^n=y^nx^{n-1}\qquad\text{pentru orice }x,y\in G:
\]
elementul $x^{n-1}$ comută cu orice putere a $n$-a. Dar $f$ fiind surjectiv, \emph{orice} $g\in G$ se scrie $g=y^n$, deci $x^{n-1}g=gx^{n-1}$ pentru toți $g$, adică $x^{n-1}\in Z(G)$.
\end{proof}

\begin{corolar}\label{cor:nn1}
Dacă $x\mapsto x^n$ și $x\mapsto x^{n+1}$ sunt ambele endomorfisme \textbf{surjective} ale lui $G$ (cu $n\ge2$), atunci $G$ este abelian.
\end{corolar}

\begin{proof}
Din Teorema~\ref{teo:endon}(ii), $x^{n-1}\in Z(G)$ și $x^{n}\in Z(G)$ pentru orice $x$. Cum $(n-1,n)=1$, Lema~\ref{lema:bezoutexp} aplicată subgrupului $H=Z(G)$ dă $x\in Z(G)$ pentru orice $x$, adică $Z(G)=G$.
\end{proof}

\begin{observatie}
Comparați cu Teorema~\ref{teo:consec}: acolo se cer \emph{trei} exponenți consecutivi, dar fără surjectivitate; aici doar doi, în schimbul surjectivității. Ambele scheme (plus identitatea (i), valabilă fără nicio ipoteză suplimentară) sunt punctul de plecare standard pentru problemele de tip ,,$(xy)^n=x^ny^n$``.
\end{observatie}

\section{Grupuri ciclice}

\begin{definitie}
$G$ este \emph{ciclic} dacă există $x\in G$ (numit \emph{generator}) cu $G=\langle x\rangle$. Exemple: $(\ZZ,+)=\langle 1\rangle$, $(\ZZ_n,+)=\langle\cl 1\rangle$, $U_n=\langle\varepsilon\rangle$. Orice grup ciclic este abelian.
\end{definitie}

\begin{teorema}[clasificarea grupurilor ciclice]\label{teo:clasif}
Orice grup ciclic infinit este izomorf cu $(\ZZ,+)$; orice grup ciclic de ordin $n$ este izomorf cu $(\ZZ_n,+)$.
\end{teorema}

\begin{proof}
Fie $G=\langle x\rangle$. Cazul infinit: $f:\ZZ\to G$, $f(k)=x^k$, este morfism ($x^{k+l}=x^kx^l$), surjectiv prin definiție și injectiv deoarece puterile sunt distincte. Cazul finit: $f:\ZZ_n\to G$, $f(\cl k)=x^k$, este \emph{bine definită} și injectivă datorită echivalenței $x^i=x^j\iff i\equiv j\pmod n$ (Teorema~\ref{teo:ordin}(ii)), surjectivă și morfism ca mai sus.
\end{proof}

\begin{teorema}[subgrupurile lui $\ZZ$]\label{teo:subgrupZ}
Subgrupurile lui $(\ZZ,+)$ sunt exact mulțimile $n\ZZ=\{nk\mid k\in\ZZ\}$, cu $n\in\NN$.
\end{teorema}

\begin{proof}
Fiecare $n\ZZ$ este subgrup (criteriul: $nk-nl=n(k-l)$). Reciproc, fie $H\le\ZZ$. Dacă $H=\{0\}$, luăm $n=0$. Altfel $H$ conține un element pozitiv (odată cu $x$ conține $-x$); fie $n$ cel mai mic element pozitiv din $H$. Evident $n\ZZ\subseteq H$. Pentru $h\in H$, împărțirea cu rest $h=nq+r$, $0\le r<n$, dă $r=h-nq\in H$, deci $r=0$ prin minimalitate. Așadar $H=n\ZZ$.
\end{proof}

\begin{teorema}[subgrupurile grupurilor ciclice]\label{teo:subgrupciclic}
Fie $G=\langle x\rangle$ ciclic de ordin $n$. Atunci: \emph{(i)} orice subgrup al lui $G$ este ciclic, de forma $\langle x^m\rangle$ cu $m\mid n$; \emph{(ii)} pentru fiecare divizor $d$ al lui $n$ există \textbf{exact un} subgrup de ordin $d$, anume $\langle x^{n/d}\rangle$.
\end{teorema}

\begin{proof}
(i) Fie $H\le G$, $H\neq\{e\}$, și $m=\min\{k\ge 1\mid x^k\in H\}$. Ca în demonstrația Teoremei~\ref{teo:subgrupZ} (împărțire cu rest la exponent), orice $x^k\in H$ are $m\mid k$; în particular, din $x^n=e\in H$ rezultă $m\mid n$ și $H=\langle x^m\rangle$. (ii) $\langle x^{n/d}\rangle$ are ordinul $\frac{n}{(n,\,n/d)}=d$ (Propoziția~\ref{prop:ordputere}). Unicitatea: un subgrup de ordin $d$ este $\langle x^m\rangle$ cu $m\mid n$, de ordin $\frac{n}{m}=d$, deci $m=\frac{n}{d}$ este determinat.
\end{proof}

\begin{propozitie}[generatorii]\label{prop:generatori}
Dacă $G=\langle x\rangle$ are ordinul $n$, atunci $x^k$ este generator dacă și numai dacă $(k,n)=1$. În particular, $G$ are exact $\varphi(n)$ generatori, unde $\varphi$ este indicatorul lui Euler.
\end{propozitie}

\begin{proof}
$x^k$ generează $G$ $\iff\ord(x^k)=n\iff\frac{n}{(n,k)}=n\iff(n,k)=1$.
\end{proof}

\begin{propozitie}[numărul soluțiilor lui $y^d=e$]\label{prop:solutii}
Într-un grup ciclic de ordin $n$, ecuația $y^d=e$ are exact $(d,n)$ soluții.
\end{propozitie}

\begin{proof}
Scriem $y=x^t$, $0\le t\le n-1$. Atunci $y^d=e\iff n\mid td\iff\frac{n}{(n,d)}\mid t$, ceea ce dă exact $(n,d)$ valori ale lui $t$ în intervalul considerat.
\end{proof}

\begin{propozitie}[grupuri fără subgrupuri proprii]\label{prop:farasub}
Dacă $G\neq\{e\}$ are ca subgrupuri doar $\{e\}$ și $G$, atunci $G$ este finit, de ordin prim.
\end{propozitie}

\begin{proof}
Fie $x\neq e$; atunci $\langle x\rangle=G$, deci $G$ este ciclic. Dacă $G$ ar fi infinit, $G\simeq\ZZ$ (Teorema~\ref{teo:clasif}) ar avea subgrupul propriu $2\ZZ$, contradicție. Deci $|G|=n$; dacă $n=ab$ cu $1<a<n$, subgrupul $\langle x^a\rangle$ are ordinul $b\notin\{1,n\}$, contradicție. Așadar $n$ este prim.
\end{proof}

\section{Grupul permutărilor: cicluri și signatură}\label{sec:permutari}

Fixăm $n\ge2$. În $S_n$ produsul $\sigma\tau=\sigma\circ\tau$ se calculează de la dreapta la stânga: întâi $\tau$, apoi $\sigma$.

\begin{definitie}[cicluri; transpoziții]\label{def:ciclu}
Fie $k\ge2$ și $a_1,\dots,a_k\in\{1,\dots,n\}$ distincte. \emph{Ciclul} $(a_1\,a_2\,\dots\,a_k)\in S_n$ este permutarea care duce $a_1\mapsto a_2\mapsto\dots\mapsto a_k\mapsto a_1$ și fixează celelalte elemente; $k$ este \emph{lungimea} ciclului, iar $\{a_1,\dots,a_k\}$ este \emph{suportul} lui. Ciclurile de lungime $2$ se numesc \emph{transpoziții}. Două cicluri sunt \emph{disjuncte} dacă suporturile lor sunt disjuncte.
\end{definitie}

\begin{teorema}[descompunerea în cicluri disjuncte]\label{teo:cicluri}
Orice permutare $\sigma\in S_n$, $\sigma\neq e$, se scrie ca produs de cicluri disjuncte două câte două, în mod unic până la ordinea factorilor, iar ciclurile disjuncte comută între ele. Dacă lungimile acestor cicluri sunt $k_1,\dots,k_r$, atunci
\[
\ord(\sigma)=[k_1,\dots,k_r].
\]
\end{teorema}

\begin{proof}
Relația $a\sim b\iff b=\sigma^j(a)$ pentru un $j\in\ZZ$ este de echivalență pe $\{1,\dots,n\}$; clasele ei se numesc \emph{orbitele} lui $\sigma$. Fie $a$ fixat și $k\ge1$ cel mai mic cu $\sigma^k(a)=a$ (există, căci $\sigma$ are ordin finit). Elementele $a,\sigma(a),\dots,\sigma^{k-1}(a)$ sunt distincte (din $\sigma^i(a)=\sigma^j(a)$ cu $0\le i<j<k$ ar rezulta $\sigma^{j-i}(a)=a$), iar orice $\sigma^j(a)$ este unul dintre ele (împărțim $j$ cu rest la $k$); deci ele formează orbita lui $a$, pe care $\sigma$ acționează exact ca ciclul $(a\ \sigma(a)\ \dots\ \sigma^{k-1}(a))$. Produsul ciclurilor asociate orbitelor cu cel puțin două elemente coincide cu $\sigma$ pe fiecare orbită, deci peste tot. Unicitatea: în orice astfel de descompunere, suporturile ciclurilor sunt exact orbitele cu cel puțin două elemente, iar pe fiecare dintre ele ciclul este determinat de $\sigma$.

Dacă $\gamma,\delta$ sunt cicluri disjuncte, atunci pe suportul lui $\gamma$ ambele produse $\gamma\delta$ și $\delta\gamma$ acționează ca $\gamma$ (căci $\delta$ fixează acele elemente, iar $\gamma$ le trimite tot în suportul său), simetric pe suportul lui $\delta$, iar în rest ambele sunt identitatea; deci $\gamma\delta=\delta\gamma$.

Un ciclu $\gamma=(a_1\,\dots\,a_k)$ are ordinul $k$: $\gamma^j(a_1)=a_{1+j}$ (indici modulo $k$), deci $\gamma^j\neq e$ pentru $1\le j<k$, iar $\gamma^k=e$. Dacă $\sigma=\gamma_1\cdots\gamma_r$ cu cicluri disjuncte, acestea comută, deci $\sigma^m=\gamma_1^m\cdots\gamma_r^m$; fiecare $\gamma_i^m$ mișcă doar elemente din suportul lui $\gamma_i$, iar suporturile sunt disjuncte, deci $\sigma^m=e\iff\gamma_i^m=e$ pentru orice $i\iff k_i\mid m$ pentru orice $i\iff[k_1,\dots,k_r]\mid m$.
\end{proof}

\begin{definitie}[signatura]\label{def:sgn}
O \emph{inversiune} a permutării $\sigma\in S_n$ este o pereche $(i,j)$ cu $i<j$ și $\sigma(i)>\sigma(j)$. Notând $\mathrm{inv}(\sigma)$ numărul inversiunilor, \emph{signatura} (semnul) lui $\sigma$ este $\sgn(\sigma)=(-1)^{\mathrm{inv}(\sigma)}$. Permutarea $\sigma$ este \emph{pară} dacă $\sgn(\sigma)=1$ și \emph{impară} dacă $\sgn(\sigma)=-1$.
\end{definitie}

\begin{teorema}[proprietățile signaturii]\label{teo:sgn}
\emph{(i)} $\sgn(\sigma)=\displaystyle\prod_{1\le i<j\le n}\frac{\sigma(j)-\sigma(i)}{j-i}$.
\emph{(ii)} $\sgn(\sigma\tau)=\sgn(\sigma)\sgn(\tau)$; deci $\sgn:S_n\to(\{\pm1\},\cdot)$ este morfism de grupuri, iar nucleul său $A_n$ (\emph{grupul altern}, format din permutările pare) este subgrup normal al lui $S_n$.
\emph{(iii)} Orice transpoziție are signatura $-1$, iar un ciclu de lungime $k$ are signatura $(-1)^{k-1}$. În consecință, dacă descompunerea lui $\sigma$ în cicluri disjuncte are ciclurile de lungimi $k_1,\dots,k_r$, atunci $\sgn(\sigma)=(-1)^{(k_1-1)+\dots+(k_r-1)}$.
\end{teorema}

\begin{proof}
(i) Când $\{i,j\}$ parcurge perechile neordonate de elemente distincte, $\{\sigma(i),\sigma(j)\}$ le parcurge pe aceleași, bijectiv; deci produsul modulelor numărătorilor este egal cu produsul numitorilor, iar produsul are câte un factor negativ pentru fiecare inversiune.

(ii) Din (i),
\[
\sgn(\sigma\tau)=\prod_{i<j}\frac{\sigma(\tau(j))-\sigma(\tau(i))}{\tau(j)-\tau(i)}\cdot\prod_{i<j}\frac{\tau(j)-\tau(i)}{j-i}.
\]
Al doilea produs este $\sgn(\tau)$. În primul, fiecare factor depinde doar de perechea neordonată $\{\tau(i),\tau(j)\}$ (schimbând ordinea, numărătorul și numitorul își schimbă amândouă semnul), iar aceasta parcurge toate perechile $\{a,b\}$ cu $a<b$; deci primul produs este $\prod_{a<b}\frac{\sigma(b)-\sigma(a)}{b-a}=\sgn(\sigma)$. Nucleul unui morfism este subgrup normal (Teorema~\ref{teo:nucleu}).

(iii) Transpoziția $(1\,2)$ are o singură inversiune, deci signatura $-1$. O transpoziție oarecare $(a\,b)$ este egală cu $\pi(1\,2)\pi^{-1}$, pentru orice $\pi\in S_n$ cu $\pi(1)=a$, $\pi(2)=b$ (verificare directă pe elemente), deci, prin (ii), $\sgn(a\,b)=\sgn(\pi)\cdot(-1)\cdot\sgn(\pi)^{-1}=-1$. Pentru un ciclu, verificăm direct pe fiecare $a_i$ că
\[
(a_1\,a_2\,\dots\,a_k)=(a_1\,a_k)(a_1\,a_{k-1})\cdots(a_1\,a_2),
\]
un produs de $k-1$ transpoziții, deci signatura lui este $(-1)^{k-1}$. Ultima afirmație rezultă din (ii) și Teorema~\ref{teo:cicluri}.
\end{proof}

\begin{exemplu}
În $S_3$: $e$ și ciclurile de lungime $3$, $(1\,2\,3)$ și $(1\,3\,2)$, sunt pare (ordin $1$, respectiv $3$), iar cele trei transpoziții sunt impare (ordin $2$); așadar $A_3=\{e,(1\,2\,3),(1\,3\,2)\}$. În $S_4$, permutarea $(1\,2)(3\,4)$ are ordinul $[2,2]=2$ și este pară, iar $(1\,2\,3\,4)$ are ordinul $4$ și este impară.
\end{exemplu}

\section{Marile teoreme din teoria grupurilor}\label{sec:mari}

\emph{Notă}: Lagrange și Cauchy fac parte din nucleul programei; acțiunile de grupuri, teoremele lui Sylow și teorema lui Frobenius o depășesc, dar sunt reunite aici --- alături de Cayley --- ca \emph{marile teoreme cu nume} din teoria grupurilor.

\subsection*{Teorema lui Lagrange}

\begin{definitie}
Pentru $H\le G$ și $x\in G$, mulțimea $xH=\{xh\mid h\in H\}$ se numește \emph{clasă la stânga} a lui $x$ în raport cu $H$.
\end{definitie}

\begin{lema}\label{lema:clase}
\emph{(i)} $x\in xH$;\quad \emph{(ii)} $xH=yH\iff x^{-1}y\in H$;\quad \emph{(iii)} două clase la stânga sunt fie egale, fie disjuncte;\quad \emph{(iv)} $|xH|=|H|$ pentru orice $x$.
\end{lema}

\begin{proof}
(i) $x=xe$. (ii) Dacă $x^{-1}y=h\in H$, atunci $y=xh$, deci $yH=x(hH)=xH$. Reciproc, $y\in yH=xH$ dă $y=xh$, deci $x^{-1}y=h\in H$. (iii) Dacă $z\in xH\cap yH$, atunci $x^{-1}z,y^{-1}z\in H$, deci $x^{-1}y=(x^{-1}z)(y^{-1}z)^{-1}\in H$ și aplicăm (ii). (iv) Funcția $h\mapsto xh$ este o bijecție de la $H$ la $xH$ (injectivă prin simplificare, surjectivă prin definiție).
\end{proof}

\begin{teorema}[Lagrange]\label{teo:lagrange}
Dacă $G$ este un grup finit și $H\le G$, atunci $|H|$ divide $|G|$. Numărul claselor la stânga, notat $[G:H]$ (\emph{indicele} lui $H$), satisface $|G|=[G:H]\cdot|H|$.
\end{teorema}

\begin{proof}
Din Lema~\ref{lema:clase}, clasele la stânga formează o partiție a lui $G$ în mulțimi de câte $|H|$ elemente.
\end{proof}

\begin{corolar}\label{cor:lagrange}
Fie $G$ finit cu $|G|=n$. Atunci: \emph{(i)} $\ord(x)\mid n$ pentru orice $x\in G$; \emph{(ii)} $x^{n}=e$ pentru orice $x\in G$; \emph{(iii)} dacă $n$ este prim, $G$ este ciclic și orice element $x\neq e$ îl generează.
\end{corolar}

\begin{proof}
(i) $\ord(x)=|\langle x\rangle|$ divide $n$ (Teorema~\ref{teo:ordin}(iii) și Lagrange). (ii) $x^n=(x^{\ord x})^{n/\ord x}=e$. (iii) Pentru $x\neq e$, $\ord(x)\in\{1,n\}$ și $\ord(x)\neq 1$, deci $\langle x\rangle$ are $n$ elemente, adică este $G$.
\end{proof}

\begin{observatie}
\textbf{Reciproca teoremei lui Lagrange este falsă}: nu pentru orice divizor $d$ al lui $|G|$ există un subgrup de ordin $d$. Exemplul clasic: grupul $A_4$ al permutărilor pare din $S_4$ ($12$ elemente) nu are subgrup de ordin $6$. Într-adevăr, un subgrup de ordin $6$ ar avea indice $2$, deci ar fi normal și ar conține toate pătratele elementelor lui $A_4$; dar pătratul oricărui $3$-ciclu este tot un $3$-ciclu, iar $A_4$ are $8$ ciclii de lungime $3$ --- prea mulți pentru un subgrup cu $6$ elemente. Reciproca este însă adevărată pentru grupurile ciclice (Teorema~\ref{teo:subgrupciclic}).
\end{observatie}

\begin{propozitie}[subgrupurile finite ale lui $\CC^*$]\label{prop:subgrupC}
Orice subgrup finit cu $n$ elemente al grupului $(\CC^*,\cdot)$ este egal cu $U_n$.
\end{propozitie}

\begin{proof}
Fie $G\le\CC^*$ cu $|G|=n$. Din Corolarul~\ref{cor:lagrange}(ii), $z^n=1$ pentru orice $z\in G$, deci $G\subseteq U_n$. Cum $|G|=n=|U_n|$, avem egalitate.
\end{proof}

\begin{observatie}
Tehnica de \emph{împerechere} $x\leftrightarrow x^{-1}$ (dintr-o demonstrație clasică a faptului că un grup de ordin par conține un element de ordin $2$) este un instrument frecvent la ONM; ea reapare la Teorema~\ref{teo:produs} (produsul tuturor elementelor unui grup abelian finit).
\end{observatie}

\subsection*{Teorema lui Cauchy}

\begin{teorema}[Cauchy]\label{teo:cauchy}
Fie $G$ un grup finit și $p$ un număr prim cu $p\mid|G|$. Atunci $G$ conține un element de ordin $p$.
\end{teorema}

\begin{proof}[Demonstrație (McKay)]
Considerăm mulțimea $p$-uplurilor
\[
S=\bigl\{(x_1,x_2,\dots,x_p)\in G^p\ \big|\ x_1x_2\cdots x_p=e\bigr\}.
\]
Primele $p-1$ componente se aleg liber, iar ultima este determinată, $x_p=(x_1\cdots x_{p-1})^{-1}$; deci $|S|=|G|^{p-1}$, număr divizibil cu $p$ (căci $p\mid|G|$ și $p-1\ge1$).

Fie $\sigma:S\to S$ ,,rotația`` $\sigma(x_1,x_2,\dots,x_p)=(x_2,\dots,x_p,x_1)$. Ea este bine definită: dacă $x_1x_2\cdots x_p=e$, atunci
\[
x_2\cdots x_p\,x_1=x_1^{-1}(x_1x_2\cdots x_p)\,x_1=e
\]
(același truc de conjugare ca în Propoziția~\ref{prop:ordinv}(iii)); este evident bijectivă și $\sigma^p=\mathrm{id}$.

Grupăm elementele lui $S$ în \emph{orbite} $O_t=\{t,\sigma(t),\dots,\sigma^{p-1}(t)\}$; relația $s\in O_t$ este de echivalență (simetria: $s=\sigma^k(t)\Rightarrow t=\sigma^{p-k}(s)$), deci orbitele partiționează $S$. Susținem că fiecare orbită are exact $p$ elemente sau exact $1$ element, ultimul caz însemnând tupluri \emph{constante} $(x,x,\dots,x)$. Într-adevăr, dacă $\sigma^k(t)=t$ pentru un $1\le k\le p-1$, alegem $u,v\in\ZZ$ cu $uk+vp=1$ (posibil, căci $(k,p)=1$) și obținem
\[
\sigma(t)=\sigma^{uk+vp}(t)=(\sigma^k)^u(t)=t,
\]
deci toate componentele lui $t$ coincid.

Fie $N$ numărul tuplurilor constante din $S$, adică numărul elementelor $x\in G$ cu $x^p=e$. Numărând pe orbite:
\[
|G|^{p-1}=|S|=N+p\cdot(\text{numărul orbitelor cu }p\text{ elemente}),
\]
deci $p\mid N$. Cum $(e,\dots,e)\in S$, avem $N\ge1$, deci $N\ge p\ge2$: există $x\neq e$ cu $x^p=e$, iar $p$ fiind prim, $\ord(x)\in\{1,p\}$, adică $\ord(x)=p$.
\end{proof}

\begin{observatie}
Demonstrația este un exemplu-tipar de \emph{numărare cu simetrie de rotație}: aceeași mulțime se numără în două moduri, folosind că orbitele nebanale au exact $p$ elemente --- aici intervine esențial primalitatea lui $p$. Cazul $p=2$ dă din nou existența unui element de ordin $2$ într-un grup de ordin par. Împreună cu Lagrange, Cauchy dă echivalența: $G$ finit are un element de ordin $p$ prim $\iff p\mid|G|$.
\end{observatie}

\begin{corolar}[ridicarea la puterea $p$]\label{lema:puterep}
Fie $G$ un grup finit cu $n$ elemente, $p\ge 2$ un întreg, și $H=\{x^p\mid x\in G\}$ mulțimea puterilor de ordin $p$. Atunci
\[
(p,n)=1 \iff H=G.
\]
(Aici $H=G$ înseamnă că aplicația $x\mapsto x^p$ este surjectivă; pe un grup finit aceasta echivalează cu injectivitatea.)
\end{corolar}

\begin{proof}
$(\Rightarrow)$ Fie $(p,n)=1$; există $a,b\in\ZZ$ cu $ap+bn=1$ (Bézout). Pentru orice $g\in G$, ordinul lui $g$ divide $n$ (Lagrange), deci $g^{n}=e$ și astfel $g^{bn}=e$. Atunci
\[
g=g^{ap+bn}=\bigl(g^{a}\bigr)^{p}\,g^{bn}=\bigl(g^{a}\bigr)^{p}\in H,
\]
deci $H=G$.

$(\Leftarrow)$ Prin contrapoziție: presupunem $(p,n)=d>1$ și fie $q$ un divizor prim al lui $d$. Cum $q\mid n$, teorema lui Cauchy (Teorema~\ref{teo:cauchy}) dă un element $x$ de ordin $q$. Cum $q\mid p$, avem $x^{p}=e=e^{p}$; deci $x\neq e$ și $e$ au aceeași imagine, aplicația $x\mapsto x^{p}$ nu e injectivă, deci (pe mulțime finită) nici surjectivă: $H\neq G$.
\end{proof}

\begin{observatie}
Enunțul nu cere ca $G$ să fie abelian și nu afirmă că $H$ este subgrup --- în grupuri neabeliene, $x\mapsto x^p$ nu e neapărat morfism, iar $H$ poate să nu fie stabil. Când $(p,n)=1$, ridicarea la puterea $p$ este chiar o \emph{bijecție} a lui $G$ (permutarea inversă e $x\mapsto x^{a}$ cu $ap\equiv 1\pmod n$), fapt des folosit la ONM pentru a extrage rădăcini de ordin $p$ în grupuri finite.
\end{observatie}

\subsection*{Teorema lui Cayley}

\begin{teorema}[Cayley]\label{teo:cayley}
Orice grup $G$ este izomorf cu un subgrup al grupului $\bigl(S(G),\circ\bigr)$ al bijecțiilor mulțimii $G$. În particular, orice grup finit cu $n$ elemente este izomorf cu un subgrup al lui $S_n$.
\end{teorema}

\begin{proof}
Pentru $g\in G$, \emph{translația la stânga} $\lambda_g:G\to G$, $\lambda_g(x)=gx$, este bijectivă (ecuația $gx=b$ are soluție unică --- Propoziția~\ref{prop:simplificare}), deci $\lambda_g\in S(G)$. Funcția $\lambda:G\to S(G)$, $\lambda(g)=\lambda_g$, este morfism:
\[
\lambda_{gh}(x)=(gh)x=g(hx)=\bigl(\lambda_g\circ\lambda_h\bigr)(x),
\]
și este injectivă: $\lambda_g=1_G$ înseamnă $gx=x$ pentru toți $x$, deci $g=e$; prin urmare $\lambda_g=\lambda_h\Rightarrow g=h$. Imaginea $\lambda(G)=\{\lambda_g\mid g\in G\}$ este parte stabilă a lui $S(G)$ (căci $\lambda_g\circ\lambda_h=\lambda_{gh}$), conține $1_G=\lambda_e$ și inversele ($\lambda_g^{-1}=\lambda_{g^{-1}}$), deci este subgrup; iar $\lambda:G\to\lambda(G)$ este morfism bijectiv, deci izomorfism. Așadar $G\simeq\lambda(G)\le S(G)$. Pentru $|G|=n$, numerotarea elementelor lui $G$ dă $S(G)\simeq S_n$.
\end{proof}

\begin{observatie}
Teorema explică rolul special al grupurilor de permutări: \emph{orice} grup ,,trăiește`` într-un $S(X)$. Scufundarea nu este economică ($S_n$ are $n!$ elemente), dar este conceptual fundamentală.
\end{observatie}

\subsection*{Teoremele lui Sylow}

Fie $|G|=p^a m$ cu $p$ prim, $a\ge1$ și $p\nmid m$. Un subgrup de ordin $p^a$ (puterea maximă a lui $p$ permisă de Lagrange) se numește \emph{$p$-subgrup Sylow}; notăm $n_p$ numărul lor.

\begin{teorema}[Sylow]\label{teo:sylow}
Fie $G$ un grup finit cu $|G|=p^am$, $p\nmid m$. Atunci:
\begin{enumerate}[label=\emph{(\Roman*)},nosep,leftmargin=2.2em]
\item \emph{(existența)} $G$ are cel puțin un $p$-subgrup Sylow;
\item \emph{(conjugarea)} orice subgrup al lui $G$ de ordin o putere a lui $p$ este conținut într-un conjugat al unui $p$-subgrup Sylow; în particular, oricare două $p$-subgrupuri Sylow sunt conjugate;
\item \emph{(numărarea)} $n_p=[G:N_G(P)]$, $\;n_p\mid m\;$ și $\;n_p\equiv1\pmod p$.
\end{enumerate}
\end{teorema}

Demonstrația --- prin acțiuni de grupuri --- se află în secțiunea \emph{Teorie avansată} (pagina~\pageref{dem:sylow}). Reținem consecința folosită cel mai des:

\begin{corolar}\label{cor:sylownormal}
$n_p=1$ dacă și numai dacă $p$-subgrupul Sylow este \textbf{normal} în $G$.
\end{corolar}

\begin{proof}
Conjugatele unui $p$-subgrup Sylow $P$ epuizează mulțimea celor $n_p$ subgrupuri Sylow (Sylow~(II)), iar numărul lor este $[G:N_G(P)]$; $P$ este normal $\iff$ nu are alte conjugate $\iff n_p=1$ (vezi și Teorema~\ref{teo:conjsub}).
\end{proof}

\subsection*{Teorema lui Frobenius}

\begin{teorema}[Frobenius, 1895]\label{teo:frobgrup}
Fie $G$ un grup finit și $n$ un divizor al lui $|G|$. Atunci numărul soluțiilor ecuației $x^n=e$ în $G$ este un \textbf{multiplu de $n$}:
\[
n\ \Big|\ \#\{x\in G\mid x^n=e\}.
\]
\end{teorema}

Pentru grupuri ciclice regăsim numărul exact $(n,|G|)=n$ (Propoziția~\ref{prop:solutii}); teorema afirmă că în orice grup acest număr rămâne divizibil cu $n$. Cazul $n=p$ prim redă concluzia demonstrației McKay a teoremei lui Cauchy. Demonstrația (numărare cu acțiuni, lema lui Burnside și coliere aperiodice) se află în secțiunea \emph{Teorie avansată} (pagina~\pageref{dem:frobenius}).

\begin{corolar}[Frobenius, forma generală cu c.m.m.d.c.]\label{cor:frobgcd}
Fie $G$ un grup finit și $n\in\NN^*$ \textbf{arbitrar} (nu neapărat divizor al lui $|G|$). Atunci
\[
\bigl(n,\,|G|\bigr)\ \Big|\ \#\{x\in G\mid x^n=e\}.
\]
\end{corolar}

\begin{proof}
Fie $d=(n,|G|)$. Susținem că \emph{cele două ecuații au aceleași soluții}: $\{x\mid x^n=e\}=\{x\mid x^d=e\}$. Într-adevăr, dacă $x^d=e$, atunci $x^n=(x^d)^{n/d}=e$, căci $d\mid n$. Reciproc, dacă $x^n=e$, atunci $\ord(x)\mid n$; dar $\ord(x)\mid|G|$ (Lagrange), deci $\ord(x)$ divide $(n,|G|)=d$, adică $x^d=e$. Cum $d\mid|G|$, Teorema~\ref{teo:frobgrup} aplicată lui $d$ dă $d\mid\#\{x\mid x^d=e\}=\#\{x\mid x^n=e\}$.
\end{proof}

\begin{obs}
Cele două formulări sunt deci echivalente (pentru $n\mid|G|$ avem $d=n$ și corolarul redă teorema), iar ipoteza $n\mid|G|$ din enunțul clasic nu se poate pur și simplu omite: în $\ZZ_2$, ecuația $x^4=e$ are $2$ soluții --- nu e multiplu de $4$, dar e multiplu de $(4,2)=2$. Forma cu c.m.m.d.c.\ este cea ,,corectă`` pentru $n$ arbitrar.
\end{obs}

\section{Teorie avansată}\label{sec:avansat}

\emph{Notă}: întreaga secțiune \textbf{depășește programa clasei a XII-a}; este însă standard în pregătirea pentru baraje, loturi și olimpiade studențești și se sprijină exclusiv pe rezultatele deja demonstrate în fișă. Cele trei subsecțiuni tratează, pe rând, structurile factor, acțiunile de grupuri (cu demonstrațiile teoremelor lui Sylow și Frobenius) și clasificările structurale (grupuri abeliene, simple, rezolubile).

\subsection{Relații de echivalență, congruențe și grupuri factor}\label{sec:avansat1}

\begin{definitie}[relații binare; relații de echivalență]\label{def:echiv}
O \emph{relație binară} pe mulțimea nevidă $A$ este o submulțime $\rho\subseteq A\times A$; scriem $x\,\rho\,y$ pentru $(x,y)\in\rho$. Relația $\rho$ este: \emph{reflexivă} dacă $x\,\rho\,x$ pentru orice $x$; \emph{simetrică} dacă $x\,\rho\,y\Rightarrow y\,\rho\,x$; \emph{tranzitivă} dacă $x\,\rho\,y$ și $y\,\rho\,z\Rightarrow x\,\rho\,z$. O relație reflexivă, simetrică și tranzitivă se numește \emph{relație de echivalență}.
\end{definitie}

\begin{definitie}[clase, mulțimea cât, sistem de reprezentanți]
Fie $\rho$ o relație de echivalență pe $A$. \emph{Clasa de echivalență} a lui $x$ este $\widehat x=\{y\in A\mid x\,\rho\,y\}$. Mulțimea claselor,
\[
A/\rho=\{\widehat x\mid x\in A\},
\]
se numește \emph{mulțimea cât} (mulțimea factor) a lui $A$ prin $\rho$. Un \emph{sistem complet de reprezentanți} este o submulțime $S\subseteq A$ care conține \textbf{exact câte un} element din fiecare clasă.
\end{definitie}

\begin{exemplu}
Congruența modulo $n$ pe $\ZZ$: $x\,\rho\,y\iff n\mid x-y$ este relație de echivalență; clasele sunt $\cl0,\cl1,\dots,\cl{n-1}$, deci $\ZZ/\rho=\ZZ_n$, iar $S=\{0,1,\dots,n-1\}$ este un sistem complet de reprezentanți (la fel $\{1,2,\dots,n\}$ sau orice alegere de câte un element din fiecare clasă).
\end{exemplu}

\begin{definitie}[partiție]
O \emph{partiție} a lui $A$ este o familie de submulțimi nevide ale lui $A$, două câte două disjuncte, a căror reuniune este $A$.
\end{definitie}

\begin{teorema}[echivalențe $\leftrightarrow$ partiții]\label{teo:partitie}
\emph{(i)} Clasele unei relații de echivalență pe $A$ formează o partiție a lui $A$ (notată $A/\rho$). \emph{(ii)} Reciproc, orice partiție a lui $A$ provine dintr-o unică relație de echivalență: $x\,\rho\,y\iff x$ și $y$ aparțin aceluiași bloc al partiției.
\end{teorema}

\begin{proof}
(i) $x\in\widehat x$ (reflexivitate), deci clasele sunt nevide și acoperă $A$. Dacă $z\in\widehat x\cap\widehat y$, atunci $x\,\rho\,z$ și $y\,\rho\,z$, deci (simetrie și tranzitivitate) $x\,\rho\,y$; de aici $\widehat x=\widehat y$, prin dublă incluziune cu tranzitivitatea. Așadar două clase sunt egale sau disjuncte. (ii) Verificare directă a celor trei proprietăți; clasele sunt exact blocurile.
\end{proof}

\begin{definitie}[relațiile canonice pe un grup]
Fie $G$ un grup și $H\subseteq G$ o submulțime. Definim pe $G$ relațiile
\[
x\,\rho_H\,y\iff x^{-1}y\in H\qquad\text{și}\qquad x\,\rho'_H\,y\iff xy^{-1}\in H .
\]
\end{definitie}

\begin{propozitie}[relațiile induse de un subgrup]\label{prop:rhoH}
$\rho_H$ este relație de echivalență pe $G$ \textbf{dacă și numai dacă} $H$ este subgrup al lui $G$ (analog pentru $\rho'_H$). În acest caz, clasa lui $x$ față de $\rho_H$ este $xH$, iar față de $\rho'_H$ este $Hx$:
\[
G/\rho_H=\{xH\mid x\in G\},\qquad G/\rho'_H=\{Hx\mid x\in G\}.
\]
$\rho_H$ și $\rho'_H$ se numesc \emph{relațiile de echivalență induse de subgrupul $H$} (la stânga, respectiv la dreapta).
\end{propozitie}

\begin{proof}
Presupunem că $\rho_H$ este echivalență. Reflexivitatea dă $e=x^{-1}x\in H$. Pentru $h\in H$ avem $e\,\rho_H\,h$ (căci $e^{-1}h=h\in H$), deci prin simetrie $h\,\rho_H\,e$, adică $h^{-1}=h^{-1}e\in H$. Pentru $h,k\in H$: $e\,\rho_H\,h$ și $h\,\rho_H\,hk$ (căci $h^{-1}(hk)=k\in H$), deci prin tranzitivitate $e\,\rho_H\,hk$, adică $hk\in H$. Așadar $H\le G$.
Reciproc, pentru $H\le G$, verificările sunt cele din Lema~\ref{lema:clase}: reflexivitate din $e\in H$; simetrie din $y^{-1}x=(x^{-1}y)^{-1}\in H$; tranzitivitate din $x^{-1}z=(x^{-1}y)(y^{-1}z)\in H$. Clasa lui $x$: $\{y\mid x^{-1}y\in H\}=xH$. Pentru $\rho'_H$, totul este simetric, cu clasele $Hx$.
\end{proof}

\begin{observatie}
În acest limbaj, teorema lui Lagrange (Teorema~\ref{teo:lagrange}) spune că partiția $G/\rho_H$ are toate blocurile de cardinal $|H|$. Numărul claselor la stânga este egal cu numărul claselor la dreapta: $xH\mapsto Hx^{-1}$ este o bijecție bine definită între $G/\rho_H$ și $G/\rho'_H$ (din $xH=yH\iff x^{-1}y\in H\iff Hx^{-1}=Hy^{-1}$), deci indicele $[G:H]$ nu depinde de parte.
\end{observatie}

\begin{teorema}[normalitatea în limbajul claselor]\label{teo:normalrho}
Pentru $H\le G$ sunt echivalente: \emph{(i)} $H\trianglelefteq G$; \emph{(ii)} $\rho_H=\rho'_H$; \emph{(iii)} partițiile $G/\rho_H$ și $G/\rho'_H$ coincid.
\end{teorema}

\begin{proof}
(i)$\iff$(ii): relațiile coincid $\iff$ clasa fiecărui $x$ este aceeași în ambele, adică $xH=Hx$ pentru orice $x$ --- exact Definiția~\ref{def:normal}. (ii)$\Rightarrow$(iii) evident. (iii)$\Rightarrow$(ii): blocul care îl conține pe $x$ este unic în fiecare partiție, anume $xH$, respectiv $Hx$; dacă familiile de blocuri coincid, atunci $xH=Hx$ pentru fiecare $x$.
\end{proof}

\begin{definitie}[congruență]\label{def:congruenta}
O relație de echivalență $\rho$ pe grupul $G$ se numește \emph{congruență} dacă este \emph{compatibilă cu operația}:
\[
x\,\rho\,y\ \text{ și }\ x'\,\rho\,y'\ \Longrightarrow\ xx'\,\rho\,yy' .
\]
\end{definitie}

\begin{observatie}[compatibilitatea cu inversarea]\label{obs:invcompat}
Într-un \textbf{grup}, dintr-o congruență rezultă automat: $x\,\rho\,y\Rightarrow x^{-1}\,\rho\,y^{-1}$. Într-adevăr, din $x^{-1}\,\rho\,x^{-1}$ și $x\,\rho\,y$ obținem $e=x^{-1}x\,\rho\,x^{-1}y$; apoi din aceasta și $y^{-1}\,\rho\,y^{-1}$ obținem $y^{-1}\,\rho\,x^{-1}yy^{-1}=x^{-1}$, iar simetria încheie. (Într-un monoid, condiția de compatibilitate cu inversarea s-ar impune separat, acolo unde are sens.)
\end{observatie}

\begin{teorema}[congruențele provin din subgrupuri normale]\label{teo:congruenta}
Fie $G$ un grup și $\rho\subseteq G\times G$ o relație de echivalență. Atunci $\rho$ este \textbf{congruență} dacă și numai dacă există un subgrup normal $N\trianglelefteq G$ astfel încât $\rho=\rho_N$. În acest caz $N=\widehat e$ (clasa elementului neutru) și $\rho_N=\rho'_N$.
\end{teorema}

\begin{proof}
($\Leftarrow$) Fie $N\trianglelefteq G$; $\rho_N$ este echivalență (Propoziția~\ref{prop:rhoH}). Compatibilitatea: dacă $x^{-1}y=n\in N$ și $x'^{-1}y'=n'\in N$, atunci
\[
(xx')^{-1}(yy')=x'^{-1}(x^{-1}y)\,y'=\bigl(x'^{-1}n\,x'\bigr)\bigl(x'^{-1}y'\bigr)=\underbrace{x'^{-1}nx'}_{\in N}\cdot\,n'\in N,
\]
unde $x'^{-1}nx'\in N$ prin normalitate. Egalitatea $\rho_N=\rho'_N$: Teorema~\ref{teo:normalrho}.

($\Rightarrow$) Fie $\rho$ congruență și $N=\widehat e=\{x\in G\mid x\,\rho\,e\}$. \emph{$N$ este subgrup}: $e\in N$; pentru $x,y\in N$, compatibilitatea dă $xy\,\rho\,e\cdot e=e$; iar $x\,\rho\,e$ implică $x^{-1}\,\rho\,e^{-1}=e$ (Observația~\ref{obs:invcompat}). \emph{$N$ este normal}: pentru $x\in N$ și $g\in G$, din $g\,\rho\,g$, $x\,\rho\,e$ și $g^{-1}\,\rho\,g^{-1}$ rezultă, aplicând compatibilitatea de două ori, $gxg^{-1}\,\rho\,geg^{-1}=e$, deci $gxg^{-1}\in N$; Propoziția~\ref{prop:normalcrit} încheie. \emph{$\rho=\rho_N$}: din $x\,\rho\,y$, înmulțind cu $x^{-1}\,\rho\,x^{-1}$, obținem $e\,\rho\,x^{-1}y$, adică $x^{-1}y\in N$; reciproc, din $e\,\rho\,x^{-1}y$, înmulțind cu $x\,\rho\,x$, obținem $x\,\rho\,y$.
\end{proof}

\begin{propozitie}[congruența asociată unui morfism]\label{prop:congmorf}
Pentru orice morfism $f:G\to G'$, relația
\[
x\sim y\iff f(x)=f(y)
\]
este o congruență pe $G$ și coincide cu $\rho_{\Ker f}=\rho'_{\Ker f}$.
\end{propozitie}

\begin{proof}
$f(x)=f(y)\iff f(x^{-1}y)=e'\iff x^{-1}y\in\Ker f$ și, la fel, $\iff f(xy^{-1})=e'\iff xy^{-1}\in\Ker f$. Deci $\sim$ este simultan $\rho_{\Ker f}$ și $\rho'_{\Ker f}$; compatibilitatea cu operația este imediată ($f(xx')=f(x)f(x')=f(y)f(y')=f(yy')$) --- în acord cu Teorema~\ref{teo:congruenta}, întrucât $\Ker f\trianglelefteq G$ (Teorema~\ref{teo:nucleu}).
\end{proof}

\begin{teorema}[grupul factor / grupul cât]\label{teo:factor}
Fie $N\trianglelefteq G$. Pe mulțimea claselor
\[
G/N:=G/\rho_N=\{xN\mid x\in G\}
\]
operația $(xN)\cdot(yN)=(xy)N$ este \textbf{bine definită} (independentă de reprezentanți) și dă o structură de grup, numit \emph{grupul factor} (sinonim: \emph{grupul cât}) al lui $G$ prin $N$: neutrul este $N=eN$, iar $(xN)^{-1}=x^{-1}N$. Funcția
\[
\pi:G\to G/N,\qquad\pi(x)=xN\qquad(\text{\emph{proiecția canonică}})
\]
este morfism surjectiv cu $\Ker\pi=N$; pentru $G$ finit, $|G/N|=[G:N]=|G|/|N|$.
\textbf{Reciproc}: dacă pentru un subgrup $H\le G$ formula $(xH)(yH)=(xy)H$ este bine definită pe clasele la stânga, atunci $H\trianglelefteq G$.
\end{teorema}

\begin{proof}
\emph{Buna definire}: dacă $x'=xn$ și $y'=ym$ cu $n,m\in N$, atunci $x'y'=xnym=xy\,(y^{-1}ny)\,m\in xyN$, folosind $y^{-1}ny\in N$ (normalitate). Asociativitatea, neutrul și inversele se verifică pe reprezentanți, moștenite din $G$. $\pi$ este morfism chiar prin definiția operației, evident surjectiv, iar $\pi(x)=eN\iff xN=N\iff x\in N$.
\emph{Reciproca}: fie $x\in G$ și $h\in H$. Perechile de reprezentanți $(x,\,x^{-1})$ și $(xh,\,x^{-1})$ definesc aceleași clase, deci buna definire cere $(x\cdot x^{-1})H=(xh\cdot x^{-1})H$, adică $xhx^{-1}\in H$ pentru orice $x,h$; criteriul de normalitate (Propoziția~\ref{prop:normalcrit}) încheie.
\end{proof}

\begin{exemplu}
$\ZZ/n\ZZ$: pe $(\ZZ,+)$, subgrupul $n\ZZ$ este normal (grup abelian), clasele lui $\rho_{n\ZZ}$ sunt clasele de resturi, iar grupul factor este exact $(\ZZ_n,+)$. Construcția lui $\ZZ_n$ din clasa a X-a este deci cazul fondator al întregii teorii.
\end{exemplu}

\begin{teorema}[teorema fundamentală de izomorfism]\label{teo:fundamental}
Fie $f:G\to G'$ un morfism de grupuri. Atunci funcția
\[
\bar f:G/\Ker f\longrightarrow\Ima f,\qquad\bar f\bigl(x\,\Ker f\bigr)=f(x),
\]
este bine definită și este \textbf{izomorfism} de grupuri:
\[
G/\Ker f\ \simeq\ \Ima f .
\]
\textbf{Consecință}: dacă $f:G\to G'$ este morfism \textbf{surjectiv}, atunci $G/\Ker f\simeq G'$.
\end{teorema}

\begin{proof}
Notăm $K=\Ker f$. Buna definire și injectivitatea, dintr-o singură echivalență: $xK=yK\iff x^{-1}y\in K\iff f(x)=f(y)$ (Propoziția~\ref{prop:congmorf}). Morfism: $\bar f\bigl((xK)(yK)\bigr)=\bar f(xyK)=f(xy)=f(x)f(y)=\bar f(xK)\,\bar f(yK)$. Surjectivitatea pe $\Ima f$ este evidentă. Consecința: $\Ima f=G'$.
\end{proof}

\begin{exemplu}[aplicații ale teoremei fundamentale]\label{ex:fundamental}
\begin{enumerate}[label=(\alph*),nosep]
\item $f:\ZZ\to\ZZ_n$, $f(k)=\cl k$: surjectiv, $\Ker f=n\ZZ$, deci $\ZZ/n\ZZ\simeq\ZZ_n$.
\item $f:(\RR,+)\to(\CC^*,\cdot)$, $f(t)=\cos2\pi t+i\sin2\pi t$: morfism (formulele de adunare), cu imaginea $U=\{z\in\CC\mid|z|=1\}$ (cercul unitate) și nucleul $\ZZ$; deci $\RR/\ZZ\simeq U$.
\item $\Phi:G\to\Aut(G)$, $\Phi(g)=i_g$, are $\Ima\Phi=\Inn(G)$ și $\Ker\Phi=Z(G)$ (Teorema~\ref{teo:inn}); deci $G/Z(G)\simeq\Inn(G)$.
\item $\mathrm{sgn}:S_n\to(\{\pm1\},\cdot)$ ($n\ge2$; Teorema~\ref{teo:sgn}): surjectiv (o transpoziție are semnul $-1$), cu nucleul $A_n$ (grupul altern al permutărilor pare); deci $S_n/A_n\simeq\ZZ_2$ și $|A_n|=n!/2$.
\end{enumerate}
\end{exemplu}

\begin{teorema}[teorema de corespondență]\label{teo:corespondenta}
Fie $N\trianglelefteq G$ și $\pi:G\to G/N$ proiecția canonică. Atunci
\[
H\ \longmapsto\ \pi(H)=H/N=\{hN\mid h\in H\}
\]
este o \textbf{bijecție} între subgrupurile lui $G$ care conțin $N$ și subgrupurile lui $G/N$, cu inversa $\overline H\mapsto\pi^{-1}(\overline H)$. Bijecția păstrează incluziunile și normalitatea: pentru $N\le H\le G$,
\[
H\trianglelefteq G\iff H/N\trianglelefteq G/N .
\]
Pentru $G$ finit, $\bigl|\pi^{-1}(\overline H)\bigr|=|\overline H|\cdot|N|$.
\end{teorema}

\begin{proof}
$\pi(H)$ și $\pi^{-1}(\overline H)$ sunt subgrupuri (Propoziția~\ref{prop:imagini}), iar $\pi^{-1}(\overline H)\supseteq\pi^{-1}(\{eN\})=N$. Compunerile: $\pi\bigl(\pi^{-1}(\overline H)\bigr)=\overline H$ din surjectivitatea lui $\pi$; iar pentru $N\le H\le G$: $x\in\pi^{-1}(\pi(H))\iff xN=hN$ pentru un $h\in H\iff x\in hN\subseteq H$ (căci $N\subseteq H$), deci $\pi^{-1}(\pi(H))=H$. Păstrarea incluziunilor este evidentă. Normalitatea: dacă $H\trianglelefteq G$, atunci $(gN)(hN)(gN)^{-1}=(ghg^{-1})N\in H/N$; reciproc, din $(ghg^{-1})N\in H/N$ rezultă $ghg^{-1}\in HN=H$. Cardinalul: $\pi^{-1}(\overline H)$ este reuniunea disjunctă a celor $|\overline H|$ clase, fiecare de cardinal $|N|$.
\end{proof}

\begin{lema}\label{lema:HNsubgrup}
Fie $G$ un grup, $H\le G$ și $N\trianglelefteq G$. Atunci $HN=\{hn\mid h\in H,\ n\in N\}$ este subgrup al lui $G$, $N\trianglelefteq HN$ și $H\cap N\trianglelefteq H$.
\end{lema}
\begin{proof}
$e=ee\in HN$. Pentru $h_1n_1,h_2n_2\in HN$, folosind normalitatea ($Nh=hN$):
$(h_1n_1)(h_2n_2)^{-1}=h_1n_1n_2^{-1}h_2^{-1}=h_1h_2^{-1}\,(h_2 n_1n_2^{-1}h_2^{-1})\in HN$,
căci $h_2n_1n_2^{-1}h_2^{-1}\in N$. Deci $HN\le G$. $N\trianglelefteq HN$ pentru că $N\trianglelefteq G$; iar $H\cap N\trianglelefteq H$ deoarece pentru $h\in H$, $h(H\cap N)h^{-1}\subseteq H$ ($H$ subgrup) și $\subseteq N$ ($N$ normal), deci $\subseteq H\cap N$.
\end{proof}

\begin{teorema}[a doua teoremă de izomorfism]\label{teo:izo2}
Fie $G$ un grup, $H\le G$ și $N\trianglelefteq G$. Atunci
\[
\frac{HN}{N}\ \cong\ \frac{H}{H\cap N}.
\]
\end{teorema}
\begin{proof}
Definim $\varphi:H\to HN/N$ prin $\varphi(h)=hN$; este morfism, ca restricție la $H$ a proiecției canonice $HN\to HN/N$. Este \emph{surjectiv}: orice element al lui $HN/N$ este $hnN=hN=\varphi(h)$ (căci $nN=N$). Nucleul: $\varphi(h)=N\iff hN=N\iff h\in N$, deci $\Ker\varphi=H\cap N$. Teorema fundamentală de izomorfism dă $H/(H\cap N)\cong\Ima\varphi=HN/N$.
\end{proof}

\begin{teorema}[a treia teoremă de izomorfism]\label{teo:izo3}
Fie $G$ un grup și $M,N\trianglelefteq G$ cu $M\subseteq N$. Atunci $M\trianglelefteq N$, $N/M\trianglelefteq G/M$ și
\[
\frac{G/M}{N/M}\ \cong\ \frac{G}{N}.
\]
\end{teorema}
\begin{proof}
$M\trianglelefteq N$ e clar ($M$ normal în $G$, deci și în subgrupul $N$). Definim $\psi:G/M\to G/N$ prin $\psi(gM)=gN$; este \emph{bine definită}, căci $M\subseteq N$ (dacă $g_1M=g_2M$ atunci $g_1^{-1}g_2\in M\subseteq N$, deci $g_1N=g_2N$). Este morfism surjectiv, iar
\[
\Ker\psi=\{gM\mid gN=N\}=\{gM\mid g\in N\}=N/M.
\]
Teorema fundamentală de izomorfism dă $(G/M)\big/(N/M)\cong G/N$ (în particular $N/M\trianglelefteq G/M$, fiind nucleu).
\end{proof}

\subsection{Acțiuni de grupuri; teoremele lui Sylow și Frobenius}\label{sec:avansat2}

Această subsecțiune dezvoltă limbajul acțiunilor de grupuri și îl folosește pentru a demonstra complet teoremele lui Sylow (Teorema~\ref{teo:sylow}) și a lui Frobenius (Teorema~\ref{teo:frobgrup}), enunțate în §\ref{sec:mari}.

\subsubsection*{Acțiuni de grupuri și ecuația claselor}

Acțiunile sunt instrumentul unificator: ele traduc probleme de grupuri în numărări de orbite, iar teoremele lui Sylow se demonstrează integral prin acest limbaj.

\begin{definitie}[acțiune de grup]\label{def:actiune}
O \emph{acțiune} a grupului $G$ pe mulțimea nevidă $X$ este o funcție $G\times X\to X$, $(g,x)\mapsto g\cdot x$, cu proprietățile
\[
e\cdot x=x\qquad\text{și}\qquad g\cdot(h\cdot x)=(gh)\cdot x\qquad(\forall\,g,h\in G,\ \forall\,x\in X).
\]
\end{definitie}

\begin{propozitie}[acțiuni $=$ morfisme în $S(X)$]\label{prop:actmorf}
A da o acțiune a lui $G$ pe $X$ revine la a da un morfism $\varphi:G\to S(X)$, prin corespondența $\varphi(g)(x)=g\cdot x$. Acțiunea se numește \emph{fidelă} dacă $\varphi$ este injectiv.
\end{propozitie}

\begin{proof}
Dintr-o acțiune: $\varphi(g)\circ\varphi(g^{-1})=\varphi(g^{-1})\circ\varphi(g)=\varphi(e)=1_X$ (axiomele), deci fiecare $\varphi(g)$ este bijecție, iar a doua axiomă spune exact $\varphi(gh)=\varphi(g)\circ\varphi(h)$. Reciproc, orice morfism $\varphi:G\to S(X)$ definește o acțiune prin $g\cdot x=\varphi(g)(x)$.
\end{proof}

\begin{exemplu}[acțiuni fundamentale]\label{ex:actiuni}
\begin{enumerate}[label=(\alph*),nosep]
\item \emph{Translațiile}: $G$ acționează pe el însuși prin $g\cdot x=gx$; acțiunea este fidelă (demonstrația teoremei lui Cayley).
\item \emph{Conjugarea pe elemente}: $G$ acționează pe $G$ prin $g\cdot x=gxg^{-1}$ (axiomele: $e\cdot x=x$, iar $g\cdot(h\cdot x)=g(hxh^{-1})g^{-1}=(gh)x(gh)^{-1}=(gh)\cdot x$).
\item \emph{Conjugarea pe submulțimi}: $G$ acționează pe mulțimea submulțimilor lui $G$ prin $g\cdot M=gMg^{-1}$; în particular pe mulțimea \emph{subgrupurilor} (conjugatul $gHg^{-1}$ al unui subgrup este tot subgrup, căci $x\mapsto gxg^{-1}$ este bijecție ce păstrează operația).
\item \emph{Acțiunea pe clase}: pentru $H\le G$, grupul $G$ acționează pe mulțimea claselor la stânga $\{xH\mid x\in G\}$ prin $g\cdot(xH)=(gx)H$; formula este bine definită ($xH=yH\Rightarrow(gx)H=(gy)H$) și \textbf{nu cere normalitatea} lui $H$.
\end{enumerate}
\end{exemplu}

\begin{definitie}[orbită; stabilizator]
Pentru o acțiune a lui $G$ pe $X$ și $x\in X$ definim
\[
\Orb(x)=G\cdot x=\{g\cdot x\mid g\in G\}\subseteq X,
\qquad
\Stab(x)=\{g\in G\mid g\cdot x=x\}\subseteq G .
\]
\end{definitie}

\begin{teorema}[orbită--stabilizator]\label{teo:orbstab}
\emph{(i)} Relația $x\sim y\iff\exists\,g\in G:\ y=g\cdot x$ este o relație de echivalență pe $X$, ale cărei clase sunt orbitele; în particular, \textbf{orbitele partiționează} $X$.
\emph{(ii)} $\Stab(x)\le G$.
\emph{(iii)} Corespondența $g\,\Stab(x)\mapsto g\cdot x$ este o bijecție bine definită de la mulțimea claselor la stânga $\{g\,\Stab(x)\mid g\in G\}$ la $\Orb(x)$; în particular, pentru $G$ finit,
\[
|\Orb(x)|=[\,G:\Stab(x)\,]\qquad\text{și}\qquad|G|=|\Orb(x)|\cdot|\Stab(x)| .
\]
\end{teorema}

\begin{proof}
(i) Reflexivă ($e\cdot x=x$), simetrică ($y=g\cdot x\Rightarrow x=g^{-1}\cdot y$), tranzitivă (compunerea acțiunilor); clasa de echivalență a lui $x$ este chiar $\Orb(x)$, iar clasele unei relații de echivalență formează o partiție. (ii) $e\in\Stab(x)$; dacă $g,h\in\Stab(x)$, atunci $h^{-1}\cdot x=h^{-1}\cdot(h\cdot x)=x$, deci $(gh^{-1})\cdot x=g\cdot x=x$; criteriul de subgrup. (iii)
\[
\begin{aligned}
g\cdot x=h\cdot x&\iff(h^{-1}g)\cdot x=x\iff h^{-1}g\in\Stab(x)\\
&\iff g\,\Stab(x)=h\,\Stab(x),
\end{aligned}
\]
deci corespondența este bine definită și injectivă; surjectivitatea pe orbită este definiția. Egalitățile finale: teorema lui Lagrange.
\end{proof}

\begin{teorema}[ecuația claselor pentru o acțiune]\label{teo:ecclase}
Fie $G$ un grup finit care acționează pe o mulțime finită $X$, și fie $x_1,\dots,x_r$ un sistem de reprezentanți ai orbitelor. Atunci
\[
|X|=\sum_{i=1}^{r}|\Orb(x_i)|=\sum_{i=1}^{r}[\,G:\Stab(x_i)\,].
\]
Separând orbitele cu un singur element (adică punctele fixe $X^{G}=\{x\mid g\cdot x=x\ \forall g\in G\}$) de cele cu cel puțin două, se obține forma
\[
|X|=|X^{G}|+\sum_{\substack{i\\ |\Orb(x_i)|\ge 2}}[\,G:\Stab(x_i)\,],
\]
în care fiecare termen al sumei este un divizor $>1$ al lui $|G|$.
\end{teorema}

\begin{proof}
Orbitele formează o partiție a lui $X$ (Teorema~\ref{teo:orbstab}(i)), deci $|X|$ este suma cardinalelor orbitelor distincte; prin orbită--stabilizator, $|\Orb(x_i)|=[\,G:\Stab(x_i)\,]$, divizor al lui $|G|$ prin Lagrange. O orbită are un singur element exact când $x_i\in X^{G}$, ceea ce dă a doua formă.
\end{proof}

\begin{teorema}[conjugarea pe elemente; ecuația claselor]\label{teo:clase}
Pentru acțiunea prin conjugare a lui $G$ pe el însuși:
\emph{(i)} orbita lui $x$ este \emph{clasa de conjugare} $\mathcal C(x)=\{gxg^{-1}\mid g\in G\}$, iar $\Stab(x)=C_G(x)$; deci $|\mathcal C(x)|=[G:C_G(x)]$ pentru $G$ finit;
\emph{(ii)} $\mathcal C(x)=\{x\}\iff x\in Z(G)$;
\emph{(iii)} \emph{(ecuația claselor)} dacă $G$ este finit și $x_1,\dots,x_r$ sunt reprezentanți ai claselor de conjugare cu cel puțin două elemente, atunci
\[
|G|=|Z(G)|+\sum_{i=1}^{r}\,[\,G:C_G(x_i)\,],
\]
fiecare termen al sumei fiind $>1$ și divizor al lui $|G|$.
\end{teorema}

\begin{proof}
(i) $g\cdot x=x\iff gxg^{-1}=x\iff gx=xg\iff g\in C_G(x)$; apoi orbită--stabilizator. (ii) Direct din (i). (iii) Orbitele partiționează $G$; clasele cu un singur element sunt, prin (ii), exact elementele centrului, iar celelalte au cardinalele $[G:C_G(x_i)]>1$, divizori ai lui $|G|$.
\end{proof}

\begin{teorema}[submulțimi conjugate]\label{teo:conjsub}
Pentru acțiunea prin conjugare pe submulțimile lui $G$: $\Stab(M)=N_G(M)$ (normalizatorul!), deci, pentru $G$ finit, numărul conjugatelor distincte ale lui $M$ este
\[
\bigl|\{gMg^{-1}\mid g\in G\}\bigr|=[\,G:N_G(M)\,].
\]
În particular, pentru un subgrup $H\le G$: numărul conjugatelor lui $H$ este $[G:N_G(H)]$, iar $H\trianglelefteq G$ dacă și numai dacă $H$ este singurul său conjugat.
\end{teorema}

\begin{proof}
$g\cdot M=M\iff gMg^{-1}=M\iff g\in N_G(M)$; apoi Teorema~\ref{teo:orbstab} și Definiția~\ref{def:normal}.
\end{proof}

\begin{propozitie}[$p^k$ și centrul]\label{prop:pkcentru}
Fie $G$ un grup finit și $p$ prim. Presupunem că $p^k\mid|G|$ pentru un $k\ge1$. Atunci fie $p^k$ divide ordinul unui subgrup \emph{propriu} al lui $G$, fie $p\mid|Z(G)|$.
\end{propozitie}
\begin{proof}
Folosim ecuația claselor $|G|=|Z(G)|+\sum_{i=1}^r[G:C_G(g_i)]$, suma fiind luată peste clasele de conjugare cu mai mult de un element (deci $g_i\notin Z(G)$, adică $C_G(g_i)\neq G$: centralizatoarele sunt subgrupuri \emph{proprii}). Dacă pentru vreun $i$ avem $p^k\mid|C_G(g_i)|$, concluzia e atinsă (subgrup propriu). Altfel, pentru fiecare $i$, din $p^k\mid|G|=[G:C_G(g_i)]\cdot|C_G(g_i)|$ și $p^k\nmid|C_G(g_i)|$ rezultă $p\mid[G:C_G(g_i)]$. Atunci $p$ divide suma $\sum_i[G:C_G(g_i)]$ și divide $|G|$, deci divide diferența $|Z(G)|=|G|-\sum_i[G:C_G(g_i)]$.
\end{proof}

\begin{corolar}[centrul $p$-grupurilor]\label{cor:pcentru}
Dacă $|G|=p^k$ cu $p$ prim și $k\ge1$ ($G$ este un \emph{$p$-grup}), atunci $Z(G)\neq\{e\}$.
\end{corolar}

\begin{proof}
În ecuația claselor, fiecare termen $[G:C_G(x_i)]$ este un divizor $>1$ al lui $p^k$, deci multiplu de $p$. Rezultă $|Z(G)|\equiv|G|\equiv0\pmod p$, deci $|Z(G)|\ge p$.
\end{proof}

\begin{lema}\label{lema:GZ}
Dacă $G/Z(G)$ este ciclic, atunci $G$ este abelian.
\end{lema}

\begin{proof}
Fie $G/Z(G)=\langle g\,Z(G)\rangle$. Orice element al lui $G$ aparține unei clase $g^i\,Z(G)$, deci se scrie $g^iz$ cu $z\in Z(G)$. Atunci, pentru $x=g^iz$ și $y=g^jw$ ($z,w\in Z(G)$):
\[
xy=g^iz\,g^jw=g^{i+j}zw=g^jw\,g^iz=yx,
\]
folosind centralitatea lui $z,w$ și comutarea puterilor lui $g$ între ele.
\end{proof}

\begin{corolar}\label{cor:p2}
Orice grup de ordin $p^2$ ($p$ prim) este abelian.
\end{corolar}

\begin{proof}
Din Corolarul~\ref{cor:pcentru} și Lagrange, $|Z(G)|\in\{p,\,p^2\}$. Dacă $|Z(G)|=p$, atunci $|G/Z(G)|=p$, deci $G/Z(G)$ este ciclic (Corolarul~\ref{cor:lagrange}(iii)), iar Lema~\ref{lema:GZ} dă $G$ abelian, adică $Z(G)=G$ --- contradicție. Deci $Z(G)=G$.
\end{proof}

\begin{propozitie}[grupuri de ordin $p^2$]\label{prop:ordp2struct}
Orice grup de ordin $p^2$ ($p$ prim) este izomorf fie cu $\ZZ_{p^2}$, fie cu $\ZZ_p\times\ZZ_p$ (caz în care orice element $\neq e$ are ordinul $p$).
\end{propozitie}

\begin{proof}
$G$ este abelian (Corolarul~\ref{cor:p2}). Dacă are un element de ordin $p^2$, este ciclic, $\simeq\ZZ_{p^2}$. Altfel orice $x\neq e$ are ordinul $p$; alegem $a\neq e$ și $b\notin\langle a\rangle$. Atunci $\langle a\rangle\cap\langle b\rangle=\{e\}$ (ordinele fiind prime), iar $(\cl i,\cl j)\mapsto a^ib^j$ este morfism (căci $G$ e abelian) de la $\ZZ_p\times\ZZ_p$ la $G$, injectiv ($a^ib^j=e\Rightarrow a^i=b^{-j}\in\langle a\rangle\cap\langle b\rangle$) și deci bijectiv, $|G|=p^2$. Așadar $G\simeq\ZZ_p\times\ZZ_p$.
\end{proof}

\begin{lema}[normalizatorul crește]\label{lema:normcreste}
Fie $G$ un $p$-grup finit și $H<G$ un subgrup propriu. Atunci $H\subsetneq N_G(H)$.
\end{lema}
\begin{proof}
Inducție după $|G|$. Centrul $Z=Z(G)$ este netrivial (Corolarul~\ref{cor:pcentru}). Dacă $Z\not\subseteq H$, alegem $z\in Z\setminus H$: fiind central, $z$ normalizează $H$, deci $z\in N_G(H)\setminus H$. Dacă $Z\subseteq H$, trecem la $G/Z$, un $p$-grup mai mic, în care $H/Z<G/Z$ este propriu; prin ipoteza de inducție $H/Z\subsetneq N_{G/Z}(H/Z)$. Prin teorema de corespondență $N_{G/Z}(H/Z)=N_G(H)/Z$, deci $H\subsetneq N_G(H)$.
\end{proof}

\begin{propozitie}[lanț normal în $p$-grupuri]\label{prop:pchain}
Fie $G$ un $p$-grup și $H<G$ un subgrup propriu. Atunci există un subgrup $K$ al lui $G$ cu $H\trianglelefteq K$ și $K/H$ ciclic de ordin $p$.
\end{propozitie}
\begin{proof}
Prin Lema~\ref{lema:normcreste}, $H\subsetneq N_G(H)$, iar $H\trianglelefteq N_G(H)$ ($H$ e normal în normalizatorul său). Grupul factor $N_G(H)/H$ este un $p$-grup netrivial, deci (teorema lui Cauchy) conține un element de ordin $p$, care generează un subgrup de ordin $p$ al lui $N_G(H)/H$. Prin teorema de corespondență, acesta este $K/H$ pentru un subgrup $K$ cu $H\trianglelefteq K\le N_G(H)$, iar $K/H$ este ciclic de ordin $p$.
\end{proof}

\begin{lema}[punctele fixe ale $p$-grupurilor]\label{lema:fix}
Fie $P$ un $p$-grup finit care acționează pe mulțimea finită $X$ și
\[
X^{P}=\{x\in X\mid g\cdot x=x\ \text{pentru orice }g\in P\}
\]
mulțimea punctelor fixe. Atunci $|X|\equiv|X^{P}|\pmod p$.
\end{lema}

\begin{proof}
Orbitele partiționează $X$ și au cardinalele $[P:\Stab(x)]$, divizori ai lui $|P|$, deci puteri ale lui $p$. Orbitele cu un singur element sunt exact punctele fixe; toate celelalte au cardinal divizibil cu $p$. Sumând pe orbite, $|X|\equiv|X^P|\pmod p$.
\end{proof}

\subsubsection*{Lema lui Burnside}

\begin{lema}[Burnside; numărarea orbitelor]\label{lema:burnside}
Dacă grupul finit $G$ acționează pe mulțimea finită $X$, atunci
\[
\sum_{g\in G}\bigl|\mathrm{Fix}(g)\bigr|=|G|\cdot(\text{numărul de orbite}),\qquad \mathrm{Fix}(g)=\{x\in X\mid g\cdot x=x\}.
\]
În particular, $\sum_{g\in G}|\mathrm{Fix}(g)|\equiv0\pmod{|G|}$.
\end{lema}

\begin{proof}
Numărăm în două moduri perechile $(g,x)$ cu $g\cdot x=x$. Sumând după $g$ obținem $\sum_g|\mathrm{Fix}(g)|$. Sumând după $x$ obținem $\sum_x|\Stab(x)|=\sum_x\frac{|G|}{|\Orb(x)|}=|G|\sum_x\frac1{|\Orb(x)|}$, iar pe fiecare orbită $O$ suma $\sum_{x\in O}\frac1{|O|}=1$, deci totalul este $|G|\cdot(\text{numărul de orbite})$.
\end{proof}

\subsubsection*{Demonstrația teoremelor lui Sylow}\label{dem:sylow}

Demonstrăm Teorema~\ref{teo:sylow} punct cu punct, folosind exclusiv acțiuni de grupuri. Fie $|G|=p^am$, $p\nmid m$.

\paragraph{(I) Existența (Wielandt).}
Fie $\Omega$ mulțimea submulțimilor lui $G$ cu exact $p^a$ elemente; atunci $|\Omega|=\binom{p^am}{p^a}$. \emph{Arătăm întâi $p\nmid\binom{p^am}{p^a}$}: scriem $\binom{p^am}{p^a}=\prod_{j=0}^{p^a-1}\frac{p^am-j}{p^a-j}$; factorul $j=0$ este $m$, nedivizibil cu $p$, iar pentru $1\le j\le p^a-1$, cu $c$ exponentul lui $p$ în $j$ (deci $c<a$), atât $p^am-j$ cât și $p^a-j$ au exponentul $p$-adic $c$, deci fiecare factor are exponentul $0$. Așadar $p\nmid|\Omega|$. Grupul $G$ acționează pe $\Omega$ prin translație la stânga, $g\cdot X=gX$; orbitele partiționează $\Omega$, deci există o orbită $O$ cu $p\nmid|O|$. Fie $X\in O$ și $H=\Stab(X)=\{g\mid gX=X\}$; prin orbită--stabilizator (Teorema~\ref{teo:orbstab}), $|O|=[G:H]$, deci $p^a\mid|H|$, adică $|H|\ge p^a$. Pe de altă parte, fixând $x_0\in X$, aplicația $h\mapsto hx_0$ trimite $H$ injectiv în $X$ (căci $hX=X$), deci $|H|\le|X|=p^a$. Prin urmare $|H|=p^a$: $H$ este $p$-subgrup Sylow. \emph{(Aceasta regenerează teorema lui Cauchy}: un element $x\neq e$ al lui $H$ are ordinul $p^j$, iar $x^{p^{j-1}}$ are ordinul $p$.\emph{)}

\paragraph{(II) Conjugarea.}
Fie $P$ un $p$-subgrup Sylow și $Q\le G$ cu $|Q|$ o putere a lui $p$. $Q$ acționează pe mulțimea claselor $X=\{xP\mid x\in G\}$ prin $q\cdot(xP)=(qx)P$ (Exemplul~\ref{ex:actiuni}(d)); $|X|=[G:P]=m$, nedivizibil cu $p$. Prin lema punctelor fixe (Lema~\ref{lema:fix}), $|X^Q|\equiv m\not\equiv0\pmod p$, deci există o clasă fixă $gP$: $(qg)P=gP$ pentru orice $q\in Q$, adică $g^{-1}Qg\subseteq P$, deci $Q\subseteq gPg^{-1}$. Dacă $Q$ este el însuși Sylow, egalitatea $|Q|=|gPg^{-1}|=p^a$ forțează $Q=gPg^{-1}$: oricare două Sylow-uri sunt conjugate.

\paragraph{(III) Numărarea.}
Fie $\Sigma$ mulțimea $p$-subgrupurilor Sylow; $G$ acționează pe $\Sigma$ prin conjugare, tranzitiv (punctul (II)), cu $\Stab(P)=N_G(P)$ (Teorema~\ref{teo:conjsub}), deci $n_p=|\Sigma|=[G:N_G(P)]$. Din $P\le N_G(P)\le G$ rezultă $[G:P]=[G:N_G(P)][N_G(P):P]$, deci $n_p\mid[G:P]=m$. Pentru congruență, restrângem acțiunea la $P$: dacă $Q\in\Sigma$ este fixat de tot $P$, atunci $P\subseteq N_G(Q)$, deci $P,Q$ sunt Sylow-uri ale lui $N_G(Q)$; prin (II) aplicat în $N_G(Q)$ ele sunt conjugate acolo, dar $Q\trianglelefteq N_G(Q)$, deci $P=Q$. Astfel $\Sigma^P=\{P\}$, iar Lema~\ref{lema:fix} dă $n_p=|\Sigma|\equiv|\Sigma^P|=1\pmod p$.

\subsubsection*{Demonstrația teoremei lui Frobenius}\label{dem:frobenius}

\begin{definitie}[funcția lui Möbius]\label{def:mobius}
Pentru $n\in\NN^*$ definim $\mu(1)=1$, $\mu(n)=(-1)^k$ dacă $n$ este produsul a $k$ numere prime distincte și $\mu(n)=0$ dacă $n$ se divide cu pătratul unui număr prim.
\end{definitie}

\begin{propozitie}[proprietatea de bază; inversiunea Möbius]\label{prop:mobius}
\emph{(i)} $\displaystyle\sum_{d\mid n}\mu(d)$ este egală cu $1$ pentru $n=1$ și cu $0$ pentru $n>1$.
\emph{(ii)} (\emph{inversiunea Möbius}) Dacă $(M_d)_{d\ge1}$ și $(N_d)_{d\ge1}$ sunt numere (sau, mai general, elemente ale unui grup abelian) cu $N_d=\sum_{e\mid d}M_e$ pentru orice $d$, atunci $M_d=\sum_{e\mid d}\mu(d/e)\,N_e$ pentru orice $d$.
\end{propozitie}

\begin{proof}
(i) Fie $n>1$ cu divizorii primi distincți $q_1,\dots,q_k$, $k\ge1$. Divizorii $d$ cu $\mu(d)\neq0$ sunt produsele a câte $j$ dintre $q_1,\dots,q_k$, pentru care $\mu(d)=(-1)^j$; suma este $\sum_{j=0}^{k}\binom kj(-1)^j=(1-1)^k=0$.
(ii) Înlocuim $N_e$ și schimbăm ordinea sumării: pentru $c\mid d$ fixat, $e$ parcurge multiplii lui $c$ care divid $d$, iar $f=d/e$ parcurge divizorii lui $d/c$. Obținem
\[
\sum_{e\mid d}\mu(d/e)\,N_e=\sum_{e\mid d}\mu(d/e)\sum_{c\mid e}M_c=\sum_{c\mid d}M_c\sum_{f\mid d/c}\mu(f)=M_d,
\]
căci, prin (i), suma interioară este $1$ pentru $c=d$ și $0$ în rest.
\end{proof}

Folosim funcția lui Möbius și inversiunea Möbius (Propoziția~\ref{prop:mobius}(ii)).

\begin{lema}[coliere aperiodice]\label{lema:coliere}
Pentru întregi pozitivi $a\mid b$,\quad $\displaystyle\sum_{f\mid a}\mu(f)\binom{b/f}{a/f}\equiv0\pmod b$.
\end{lema}

\begin{proof}
$\binom{b/f}{a/f}$ numără șirurile binare de lungime $b$ cu $a$ cifre de $1$ care sunt invariante la rotația de perioadă $f$ (un astfel de șir este determinat de primele $b/f$ poziții, care conțin $a/f$ cifre de $1$). Prin inversiune Möbius, suma numără șirurile \emph{aperiodice} (fără nicio simetrie de rotație netrivială). Cele $b$ rotații ale unui șir aperiodic sunt distincte, deci șirurile aperiodice se grupează în orbite de câte $b$, iar numărul lor este multiplu de $b$.
\end{proof}

\emph{Demonstrația Teoremei~\ref{teo:frobgrup}} (D.~Speyer). Fie $n\mid|G|$. Pentru $d\mid|G|$ notăm $M_d$ numărul elementelor de ordin exact $d$ și $N_d=\#\{x:x^d=e\}=\sum_{e\mid d}M_e$. Grupul $G$ acționează prin translație la stânga pe mulțimea $\Omega$ a submulțimilor cu $n$ elemente. Un element $g$ de ordin $d$ fixează $X\in\Omega$ dacă și numai dacă $X$ este reuniune de clase la stânga ale lui $\langle g\rangle$; aceasta cere $d\mid n$, iar atunci se aleg $n/d$ din cele $|G|/d$ clase, deci $|\mathrm{Fix}(g)|=\binom{|G|/d}{n/d}$ (și $0$ dacă $d\nmid n$). Lema lui Burnside (Lema~\ref{lema:burnside}) dă
\[
\sum_{d\mid n}\binom{|G|/d}{n/d}\,M_d\equiv0\pmod{|G|}.\tag{$\dagger$}
\]
Înlocuind $M_d=\sum_{e\mid d}\mu(d/e)N_e$ și schimbând ordinea sumării (cu $f=d/e$):
\[
\sum_{e\mid n}N_e\,S_e\equiv0\pmod{|G|},\qquad S_e:=\sum_{f\mid n/e}\mu(f)\binom{|G|/(ef)}{n/(ef)}.\tag{$\ddagger$}
\]
Prin Lema~\ref{lema:coliere} cu $a=n/e$, $b=|G|/e$ (și $a\mid b$, căci $n\mid|G|$), fiecare $S_e$ este divizibil cu $|G|/e$. Procedăm prin inducție după $n$ (baza $n=1$: $N_1=1$). Pentru orice divizor $e<n$ al lui $n$, ipoteza de inducție dă $e\mid N_e$, deci $N_eS_e$ este divizibil cu $e\cdot\frac{|G|}{e}=|G|$. Din $(\ddagger)$ rezultă că și termenul $e=n$ este divizibil cu $|G|$; dar $S_n=\binom{|G|/n}{1}=|G|/n$, deci
\[
N_n\cdot\frac{|G|}{n}\equiv0\pmod{|G|}\quad\Longrightarrow\quad n\mid N_n.\qquad\blacksquare
\]

\subsubsection*{Aplicații}

\begin{aplicatie}[grupurile de ordin 15 sunt ciclice]\label{apl:15}
Fie $|G|=15=3\cdot5$. Din Sylow~III: $n_3\equiv1\pmod3$ și $n_3\mid5$, deci $n_3=1$; analog $n_5\equiv1\pmod5$ și $n_5\mid3$, deci $n_5=1$. Fie $P\trianglelefteq G$ Sylow de ordin $3$ și $Q\trianglelefteq G$ Sylow de ordin $5$ (Corolarul~\ref{cor:sylownormal}); $P\cap Q=\{e\}$, căci ordinul intersecției divide și $3$, și $5$. Fie $x\in P$ de ordin $3$ și $y\in Q$ de ordin $5$. Comutatorul lor stă în intersecție:
\[
xyx^{-1}y^{-1}=\underbrace{x\,\bigl(yx^{-1}y^{-1}\bigr)}_{\in P,\ \text{căci }P\trianglelefteq G}=\underbrace{\bigl(xyx^{-1}\bigr)\,y^{-1}}_{\in Q,\ \text{căci }Q\trianglelefteq G}\ \in\ P\cap Q=\{e\},
\]
deci $xy=yx$. Atunci $\ord(xy)=3\cdot5=15$ (Teorema~\ref{teo:ordprodus}, cazul coprim), deci $G=\langle xy\rangle$ este ciclic.
\end{aplicatie}

\begin{propozitie}[produs de două numere prime]\label{prop:pqciclic}
Fie $p<q$ prime cu $p\nmid q-1$. Atunci orice grup de ordin $pq$ este ciclic. (Teoremele lui Sylow sunt astfel punctul de intrare în clasificarea grupurilor finite mici.)
\end{propozitie}

\begin{proof}
Din Sylow (Teorema~\ref{teo:sylow}): $n_q\equiv1\pmod q$ și $n_q\mid p$ dau $n_q=1$; iar $n_p\equiv1\pmod p$, $n_p\mid q$ și $p\nmid q-1$ dau $n_p=1$. Deci Sylow-urile $P$ (ordin $p$) și $Q$ (ordin $q$) sunt normale, cu $P\cap Q=\{e\}$. Exact ca în Aplicația~\ref{apl:15}, generatorii lor comută, iar produsul are ordinul $pq$ (Teorema~\ref{teo:ordprodus}), deci $G$ este ciclic.
\end{proof}


\subsubsection*{Numere ciclice}

\begin{definitie}
Un număr natural $n\ge1$ se numește \emph{număr ciclic} dacă \textbf{orice} grup de ordin $n$ este ciclic.
\end{definitie}

\begin{teorema}[caracterizarea numerelor ciclice]\label{teo:numereciclice}
Un număr $n$ este ciclic dacă și numai dacă
\[
\bigl(n,\varphi(n)\bigr)=1 .
\]
\end{teorema}

\begin{proof}
Vom folosi în mod repetat că, pentru $q$ prim, $\Aut(C_q)\simeq(\ZZ_q^*,\cdot)$ este ciclic de ordin $q-1$ (Teorema~\ref{teo:ciclicitate}).

\smallskip
\emph{$(n,\varphi(n))=1\Rightarrow n$ este liber de pătrate.} Dacă $p^2\mid n$ pentru un prim $p$, atunci $p\mid n$ și, cum $\varphi(p^2)=p(p-1)$ divide $\varphi(n)$, avem $p\mid\varphi(n)$; deci $p\mid(n,\varphi(n))$, contradicție. Așadar $n=p_1p_2\cdots p_k$ cu primi distincți, iar condiția $(n,\varphi(n))=1$ (cu $\varphi(n)=\prod_i(p_i-1)$) se rescrie:
\[
p_i\nmid p_j-1\quad\text{pentru orice }i,j.\tag{$\ast$}
\]

\smallskip
\emph{Necesitatea (}$n$ ciclic $\Rightarrow(n,\varphi(n))=1$\emph{).} Prin contrapunere: presupunem $(n,\varphi(n))>1$ și construim un grup \textbf{ne}ciclic de ordin $n$. Fie $p$ prim cu $p\mid n$ și $p\mid\varphi(n)$.
\begin{itemize}[nosep,leftmargin=1.4em]
\item Dacă $p^2\mid n$: grupul abelian $C_p\times C_{n/p}$ are ordinul $n$, dar exponentul $[p,\,n/p]=n/p<n$ (căci $p\mid n/p$), deci nu are element de ordin $n$: neciclic.
\item Dacă $p\,\|\,n$ (exact o dată), din $p\mid\varphi(n)=\prod_{q}\varphi(q^{a_q})$ rezultă $p\mid\varphi(q^{a})=q^{a-1}(q-1)$ pentru un prim $q\neq p$; cum $p\neq q$, obligatoriu $p\mid q-1$, adică $q\equiv1\pmod p$. Fie $H\le\ZZ_q^*$ subgrupul (ciclic) de ordin $p$ (există, căci $p\mid q-1=|\ZZ_q^*|$). Mulțimea aplicațiilor afine
\[
G_0=\bigl\{\,x\mapsto ax+b\ \big|\ a\in H,\ b\in\ZZ_q\,\bigr\}
\]
este grup față de compunere, de ordin $|H|\cdot q=pq$, și este \textbf{neabelian} (compunerea $x\mapsto a(x+b)$ vs.\ $x\mapsto ax+b$ diferă pentru $a\neq1$, $b\neq0$). Atunci $G_0\times C_{n/(pq)}$ este neabelian, deci neciclic, de ordin $n$.
\end{itemize}

\smallskip
\emph{Suficiența (}$(n,\varphi(n))=1\Rightarrow n$ ciclic\emph{).} Fie $G$ de ordin $n$ liber de pătrate, cu condiția $(\ast)$; arătăm prin inducție după $k$ că $G$ este ciclic. Pentru $k=1$, $|G|=p$ este prim, deci ciclic. Pentru pasul de inducție, fie $p$ cel mai mic prim care divide $n$ și $P$ un $p$-subgrup Sylow (de ordin $p$, deci $P\simeq C_p$). Grupul $N_G(P)/C_G(P)$ se scufundă în $\Aut(P)\simeq C_{p-1}$, deci ordinul său divide $p-1$; dar divide și $n$. Orice divizor prim al lui $p-1$ este $<p$, deci (prin minimalitatea lui $p$) nu divide $n$; așadar $\bigl(p-1,n\bigr)=1$ și $N_G(P)=C_G(P)$, adică $P\le Z\bigl(N_G(P)\bigr)$. Prin \emph{teorema complementului normal a lui Burnside} (rezultat clasic, a cărui demonstrație folosește transferul și depășește cadrul fișei), $G$ admite un \emph{$p$-complement normal}: un subgrup normal $M\trianglelefteq G$ de ordin $n/p$, cu $G=M\rtimes P$. Ordinul $n/p$ satisface aceeași condiție, deci prin inducție $M\simeq C_{n/p}$ este ciclic. Acțiunea prin conjugare $P\to\Aut(M)$ are imaginea de ordin divizor al lui $(p,\varphi(n/p))$; dar $\varphi(n/p)\mid\varphi(n)$ și $(p,\varphi(n))=1$, deci acțiunea este trivială: $G=M\times P\simeq C_{n/p}\times C_p$. Cum $(n/p,p)=1$, produsul este ciclic (Teorema~\ref{teo:ordprodus}, cazul coprim), deci $G\simeq C_n$.
\end{proof}

\begin{observatie}[ciclic vs.\ abelian --- o precizare importantă]\label{obs:ciclicabelian}
Condiția $(n,\varphi(n))=1$ caracterizează numerele pentru care orice grup de ordin $n$ este \textbf{ciclic}, \emph{nu} doar abelian. Cele două noțiuni chiar diferă: pentru $n=4$, orice grup de ordin $4$ este abelian ($C_4$ sau grupul lui Klein), dar $(4,\varphi(4))=(4,2)=2\neq1$; într-adevăr, grupul lui Klein nu este ciclic. Mai general, orice $n=p^2$ este ,,număr abelian`` (grupurile de ordin $p^2$ sunt abeliene --- Corolarul~\ref{cor:p2}) fără a fi număr ciclic. Mulțimea numerelor $n$ pentru care orice grup de ordin $n$ este abelian este deci strict mai mare. Caracterizarea ei (teoremă clasică, pe care o cităm fără demonstrație) este: $n=p_1^{a_1}\cdots p_k^{a_k}$ este ,,număr abelian`` dacă și numai dacă toți exponenții satisfac $a_i\le2$ (adică $n$ este liber de cuburi) și $p_i\nmid p_j^{a_j}-1$ pentru orice $i,j$. Numerele abeliene mai mici decât $200$ care nu sunt ciclice sunt $4,9,25,45,49,99,121,153,169,175$. De pildă $55=5\cdot11$ \textbf{nu} este abelian: cum $5\mid11-1$, construcția din demonstrația Teoremei~\ref{teo:numereciclice} dă un grup neabelian de ordin $55$.
\end{observatie}

\begin{corolar}
Orice număr liber de pătrate $n$ cu proprietatea $(\ast)$ --- niciun factor prim nu divide (vreun alt factor prim) $-1$ --- este ciclic. În particular, orice număr prim, orice produs $pq$ de primi cu $p\nmid q-1$ și $q\nmid p-1$ (de exemplu $15=3\cdot5$, $33=3\cdot11$, $35=5\cdot7$) este ciclic: orice grup de acel ordin este $C_n$.
\end{corolar}

\begin{proof}
Direct din Teorema~\ref{teo:numereciclice} și rescrierea $(\ast)$ a condiției $(n,\varphi(n))=1$.
\end{proof}

\subsection{Produse directe de grupuri}\label{sec:produsedirecte}

\begin{definitie}
Fie $G_1,\dots,G_n$ grupuri. \emph{Produsul direct} $G_1\times\cdots\times G_n$ este produsul cartezian al mulțimilor, înzestrat cu operația pe componente
\[
(x_1,\dots,x_n)(y_1,\dots,y_n)=(x_1y_1,\dots,x_ny_n).
\]
\end{definitie}

\begin{propozitie}
$G_1\times\cdots\times G_n$ este grup, cu neutrul $(e_1,\dots,e_n)$ și inversul $(x_1,\dots,x_n)^{-1}=(x_1^{-1},\dots,x_n^{-1})$; este abelian dacă și numai dacă toate $G_i$ sunt abeliene, iar $|G_1\times\cdots\times G_n|=\prod_i|G_i|$. Dacă fiecare $x_i$ are ordin finit, atunci
\[
\ord\big((x_1,\dots,x_n)\big)=\big[\ord(x_1),\dots,\ord(x_n)\big]\quad(\text{c.m.m.m.c.}).
\]
\end{propozitie}
\begin{proof}
Axiomele de grup se verifică pe componente. Pentru ordin: $(x_1,\dots,x_n)^k=(x_1^k,\dots,x_n^k)=(e_1,\dots,e_n)$ dacă și numai dacă $\ord(x_i)\mid k$ pentru fiecare $i$, adică dacă și numai dacă $[\ord(x_1),\dots,\ord(x_n)]\mid k$; cel mai mic astfel de $k$ este c.m.m.m.c.
\end{proof}

\begin{corolar}[criteriul de ciclicitate a produsului]\label{cor:produsciclic}
Pentru grupuri ciclice, $C_m\times C_n\cong C_{mn}$ dacă și numai dacă $(m,n)=1$. În particular $C_2\times C_3\cong C_6$ este ciclic, pe când $C_2\times C_2$ (grupul lui Klein) nu este ciclic.
\end{corolar}
\begin{proof}
Fie $x,y$ generatori ai lui $C_m$, respectiv $C_n$. Ordinul lui $(x,y)$ este $[m,n]$. Dacă $(m,n)=1$, atunci $[m,n]=mn=|C_m\times C_n|$, deci $(x,y)$ generează, iar produsul e ciclic. Dacă $d=(m,n)>1$, atunci $[m,n]=\tfrac{mn}{d}<mn$: orice element are ordinul divizor al lui $[m,n]<mn$, deci nu există generator (Teorema despre ordinul unui produs de elemente care comută).
\end{proof}

\subsection{Grupuri abeliene finit generate; grupuri simple și rezolubile}\label{sec:avansat3}

Această subsecțiune încununează teoria grupurilor cu trei rezultate structurale majore, construite pe Sylow, teorema de corespondență și grupurile factor.

\subsubsection*{Grupuri abeliene finit generate}

Lucrăm în grupul aditiv $\ZZ^n$ al $n$-uplurilor de întregi. O listă $b_1,\dots,b_r\in\ZZ^n$ este \emph{bază integrală} (sau $\ZZ$-bază) a unui subgrup $H\le\ZZ^n$ dacă orice element al lui $H$ se scrie \textbf{unic} sub forma $m_1b_1+\dots+m_rb_r$ cu $m_i\in\ZZ$ (și fiecare astfel de combinație aparține lui $H$).

\begin{observatie}
Elementele $b_1,\dots,b_n$ formează o $\ZZ$-bază a lui $\ZZ^n$ însuși dacă și numai dacă matricea (întreagă) cu liniile $b_i$ are determinantul $\pm1$ (matrice \emph{unimodulară}); inversa unei astfel de matrice este tot întreagă. Schimbările de bază admisibile în $\ZZ^n$ sunt deci exact transformările unimodulare.
\end{observatie}

\begin{teorema}[bază adaptată; ,,stacked basis``]\label{teo:adaptata}
Fie $H\le\ZZ^n$ un subgrup netrivial. Atunci există o bază integrală $(b_1,\dots,b_n)$ a lui $\ZZ^n$, un întreg $s$ cu $1\le s\le n$ și întregi pozitivi $k_1,\dots,k_s$ astfel încât $(k_1b_1,\dots,k_sb_s)$ este bază integrală a lui $H$.
\end{teorema}

\begin{proof}
Inducție după $n$. Pentru $n=1$: orice subgrup netrivial al lui $\ZZ$ este $k\ZZ$ (Teorema~\ref{teo:subgrupZ}), deci $b_1=1$, $k_1=k$.

Fie $n>1$. Printre toate perechile (bază $(u_1,\dots,u_n)$ a lui $\ZZ^n$; element $h\in H$), scriem $h=\sum_i m_iu_i$ și alegem $k_1$ = \emph{cea mai mică valoare strict pozitivă} a unui coeficient $m_i$ care apare. Putem presupune că indicele este $1$: există o bază $(u_1,\dots,u_n)$ și $h_0=k_1u_1+m_2u_2+\dots+m_nu_n\in H$.

\emph{Toți $m_i$ se divid cu $k_1$}: scriem $m_i=q_ik_1+r_i$, $0\le r_i<k_1$, și punem $b_1=u_1+\sum_{i\ge2}q_iu_i$. Atunci $(b_1,u_2,\dots,u_n)$ este tot bază (schimbare unimodulară) și $h_0=k_1b_1+\sum_{i\ge2}r_iu_i$. Minimalitatea lui $k_1$ forțează $r_i=0$, deci $h_0=k_1b_1$.

Fie $\iota:\ZZ^{n-1}\to\ZZ^n$ morfismul injectiv $(m_2,\dots,m_n)\mapsto\sum_{i\ge2}m_iu_i$ și $H'=\iota^{-1}(H)$. Orice $h\in H$ se scrie unic $h=m\,b_1+\iota(h'')$; scriind $m=qk_1+r$ ($0\le r<k_1$), avem $h-qh_0=r\,b_1+\iota(\dots)$, iar minimalitatea lui $k_1$ dă $r=0$. Așadar $k_1\mid m$ și $\iota(h'')=h-qk_1b_1\in H$, deci $h''\in H'$. Prin ipoteza de inducție, $H'$ are o bază adaptată: există o bază $(b'_2,\dots,b'_n)$ a lui $\ZZ^{n-1}$ și $k_2,\dots,k_s>0$ cu $(k_2b'_2,\dots,k_sb'_s)$ bază a lui $H'$. Punând $b_i=\iota(b'_i)$, lista $(b_1,\dots,b_n)$ este bază a lui $\ZZ^n$, iar $(k_1b_1,\dots,k_sb_s)$ este bază a lui $H$.
\end{proof}

\begin{lema}\label{lema:fg}
Un grup abelian netrivial $G$ este \emph{finit generat} dacă și numai dacă există $n\ge1$ și un morfism surjectiv $\pi:\ZZ^n\to G$.
\end{lema}

\begin{proof}
Dacă $\pi:\ZZ^n\to G$ este surjectiv, atunci $G$ este generat de $\pi(e_1),\dots,\pi(e_n)$ (baza canonică). Reciproc, dacă $G=\langle g_1,\dots,g_n\rangle$ (aditiv), morfismul $\pi(m_1,\dots,m_n)=m_1g_1+\dots+m_ng_n$ este bine definit (comutativitate) și surjectiv.
\end{proof}

\begin{teorema}[structura grupurilor abeliene finit generate]\label{teo:structura}
Fie $G$ un grup abelian finit generat, netrivial. Atunci există $n\ge1$ și $s$ cu $0\le s\le n$ astfel încât
\[
G\ \simeq\ C_{k_1}\times C_{k_2}\times\dots\times C_{k_s}\times\ZZ^{\,n-s},
\]
unde $C_{k_i}$ este grup ciclic de ordin $k_i>0$ (dacă $s=0$, $G\simeq\ZZ^n$). Mai mult, se pot alege $k_1,\dots,k_s$ cu $k_1>1$ și $k_1\mid k_2\mid\dots\mid k_s$ (\emph{factorii invarianți}), caz în care descompunerea este unic determinată.
\end{teorema}

\begin{proof}
Prin Lema~\ref{lema:fg}, fie $\pi:\ZZ^n\to G$ surjectiv și $H=\Ker\pi$, deci $G\simeq\ZZ^n/H$ (teorema fundamentală, Teorema~\ref{teo:fundamental}). Dacă $H=\{0\}$, $G\simeq\ZZ^n$. Altfel, alegem o bază adaptată $(b_1,\dots,b_n)$ cu $(k_1b_1,\dots,k_sb_s)$ bază a lui $H$ (Teorema~\ref{teo:adaptata}). Morfismul care duce $\sum m_ib_i$ în \[ (\cl{m_1},\dots,\cl{m_s},m_{s+1},\dots,m_n)\in C_{k_1}\times\dots\times C_{k_s}\times\ZZ^{n-s} \] este surjectiv și are nucleul exact $H$; teorema fundamentală încheie. Reducerea la forma cu $k_1\mid\dots\mid k_s$ și unicitatea sunt un pas suplimentar standard (nu-l detaliem).
\end{proof}

\begin{corolar}[grupuri abeliene finite]\label{cor:abelianfinit}
Orice grup abelian finit netrivial este produs direct de grupuri ciclice:
$G\simeq C_{k_1}\times\dots\times C_{k_n}$.
\end{corolar}

\begin{proof}
În Teorema~\ref{teo:structura}, finitudinea forțează $s=n$ (nu poate apărea factor $\ZZ$).
\end{proof}

\begin{observatie}[descompunerea primară]
Combinând cu lema chineză (Teorema~\ref{teo:crt}), fiecare $C_{k}$ cu $k=p_1^{a_1}\cdots p_r^{a_r}$ se descompune $C_k\simeq C_{p_1^{a_1}}\times\dots\times C_{p_r^{a_r}}$. Astfel orice grup abelian finit este produs de grupuri ciclice de ordin putere de prim (forma cu \emph{divizori elementari}). De exemplu, grupurile abeliene de ordin $8$ sunt $C_8$, $C_4\times C_2$, $C_2\times C_2\times C_2$.
\end{observatie}

\begin{teorema}[produsul tuturor elementelor; generalizarea teoremei lui Wilson]\label{teo:produs}
Fie $G$ un grup \textbf{abelian} finit și $H=\{x\in G\mid x^2=e\}$. Atunci:
\emph{(i)} $H\le G$;\quad
\emph{(ii)} $\displaystyle\prod_{x\in G}x=\prod_{x\in H}x$, adică produsul tuturor elementelor se reduce la produsul elementelor de ordin cel mult $2$;\quad
\emph{(iii)} dacă $G$ are \textbf{exact un} element $t$ de ordin $2$, produsul tuturor elementelor este $t$; în toate celelalte cazuri produsul este $e$.
\end{teorema}

\begin{proof}
(i) $e\in H$, iar pentru $x,y\in H$ avem $(xy^{-1})^2=x^2(y^{-1})^2=e$ (comutativitate), deci se aplică criteriul de subgrup.

(ii) Elementele din $G\setminus H$ satisfac $x\neq x^{-1}$ și se grupează în perechi disjuncte $\{x,x^{-1}\}$ cu produsul $e$; grupul fiind abelian, putem rearanja factorii, iar aceste perechi se simplifică.

(iii) Dacă $H=\{e\}$, produsul este $e$; dacă $H=\{e,t\}$, produsul este $t$. Rămâne cazul în care $H$ conține două elemente distincte $a,b\neq e$. Atunci $K=\{e,a,b,ab\}$ este un subgrup cu $4$ elemente al lui $H$: $ab\notin\{e,a,b\}$ (căci $ab=e\Rightarrow b=a$, $ab=a\Rightarrow b=e$, $ab=b\Rightarrow a=e$), iar închiderea se verifică direct ($a\cdot ab=b$, $b\cdot ab=a$, $(ab)^2=a^2b^2=e$). Clasele $xK$ partiționează $H$ (Lema~\ref{lema:clase}), iar produsul pe fiecare clasă este
\[
x\cdot xa\cdot xb\cdot xab \;=\; x^4\,(ab)^2 \;=\; e,
\]
folosind comutativitatea și $x^4=(x^2)^2=e$. Înmulțind pe clase, $\prod_{x\in H}x=e$.
\end{proof}

\begin{observatie}
Pentru $G=(\ZZ_p^*,\cdot)$ cu $p$ prim impar, singurul element de ordin $2$ este $-\cl1$ (soluțiile ecuației $x^2=\cl1$ într-un corp sunt $\pm\cl1$), deci teorema dă $\cl1\cdot\cl2\cdots\cl{p-1}=-\cl1$ --- exact teorema lui Wilson (Teorema~\ref{teo:wilson}), văzută acum ca fapt de teoria grupurilor.
\end{observatie}

\subsubsection*{Grupuri simple}

\begin{definitie}
Un grup netrivial $G$ este \emph{simplu} dacă singurele sale subgrupuri normale sunt $\{e\}$ și $G$.
\end{definitie}

\begin{lema}\label{lema:abeliansimplu}
Un grup abelian este simplu dacă și numai dacă este ciclic de ordin prim.
\end{lema}

\begin{proof}
Într-un grup abelian orice subgrup este normal (Exemplul~\ref{ex:normal}(a)); deci simplitatea revine la a nu avea subgrupuri proprii netriviale, ceea ce, prin Propoziția~\ref{prop:farasub}, caracterizează grupurile ciclice de ordin prim. Reciproc, $C_p$ este simplu (Corolarul~\ref{cor:lagrange}(iii)).
\end{proof}

\begin{lema}[formula produsului]\label{lema:produsHK}
Fie $H,K$ subgrupuri finite ale unui grup $G$. Atunci mulțimea $HK=\{hk\mid h\in H,\ k\in K\}$ satisface
\[
|HK|\cdot|H\cap K|=|H|\cdot|K|.
\]
\end{lema}

\begin{proof}
Considerăm $\psi:H\times K\to HK$, $\psi(h,k)=hk$ (notația $\psi$, pentru a nu o confunda cu funcția lui Möbius). Avem $h_1k_1=h_2k_2\iff h_2^{-1}h_1=k_2k_1^{-1}\in H\cap K$; notând $x=h_2^{-1}h_1\in H\cap K$, perechile cu aceeași imagine sunt exact $(h_1x^{-1},xk_1)$, în număr de $|H\cap K|$. Deci fiecare element din $HK$ are exact $|H\cap K|$ preimagini, de unde $|H\times K|=|HK|\cdot|H\cap K|$.
\end{proof}

\begin{lema}\label{lema:sylownormalsimplu}
Dacă un grup finit $G$ are un $p$-subgrup Sylow normal propriu netrivial (echivalent, $n_p=1$ cu $p^a<|G|$), atunci $G$ nu este simplu. În particular, dacă din Sylow~III rezultă $n_p=1$ pentru un prim $p\mid|G|$ cu $|G|\neq p^a$, grupul nu e simplu.
\end{lema}

\begin{proof}
Un $p$-subgrup Sylow normal netrivial și diferit de $G$ este un subgrup normal propriu netrivial (Corolarul~\ref{cor:sylownormal}).
\end{proof}

\begin{teorema}[grupuri de ordine mici]\label{teo:simplemici}
Următoarele grupuri finite nu sunt simple (deci singurele grupuri simple de ordin $<60$ sunt cele ciclice de ordin prim):
\emph{(i)} orice grup de ordin $p^k$ cu $k\ge2$ (are subgrup normal de ordin $p$);
\emph{(ii)} orice grup de ordin $pq$ cu $p<q$ prime (are Sylow $q$-subgrup normal);
\emph{(iii)} grupurile de ordin $18,50,54$ (Sylow-ul de la primul impar este unic), $20,28,40,44,45,52$ (Sylow-ul de la cel mai mare prim este unic), $42$ (Sylow $7$ unic), $30$ și $56$ (numărare de elemente), $12,24,36,48$ (intersecții de Sylow-uri).
Orice număr compus mai mic decât $60$ se află într-unul dintre cazurile (i), (ii), (iii) (verificare directă pe lista numerelor compuse), iar un grup de ordin prim este ciclic, deci simplu.
\end{teorema}

\begin{proof}
(i) Un grup de ordin $p^k$ are centru netrivial (Corolarul~\ref{cor:pcentru}); prin Cauchy, $Z(G)$ conține un element de ordin $p$, generând un subgrup normal de ordin $p$ (subgrup al centrului --- Exemplul~\ref{ex:normal}(b)).

(ii) Sylow~III: $n_q\equiv1\pmod q$ și $n_q\mid p$; cum $1<p<q$, singura posibilitate este $n_q=1$, deci Sylow-ul de ordin $q$ este normal (Lema~\ref{lema:sylownormalsimplu}).

(iii) \emph{Ordin $18=2\cdot3^2$}: $n_3\equiv1\pmod3$, $n_3\mid2\Rightarrow n_3=1$: Sylow $3$ (ordin $9$) normal. Analog $50=2\cdot5^2$ (Sylow $25$) și $54=2\cdot3^3$ (Sylow $27$). \emph{Ordin $42=2\cdot3\cdot7$}: $n_7\equiv1\pmod7$, $n_7\mid6\Rightarrow n_7=1$. \emph{Ordin $30$}: dacă $n_3>1$ atunci $n_3=10$, dând $20$ elemente de ordin $3$; dacă $n_5>1$ atunci $n_5=6$, dând $24$ elemente de ordin $5$; nu pot coexista în $30$ elemente, deci un Sylow (de ordin $3$ sau $5$) este unic, deci normal.

\emph{Ordinele $20=2^2\cdot5$, $28=2^2\cdot7$, $40=2^3\cdot5$, $44=2^2\cdot11$, $45=3^2\cdot5$, $52=2^2\cdot13$}: pentru cel mai mare prim $q$ ($5,7,5,11,5,13$), Sylow~III dă $n_q\equiv1\pmod q$ și $n_q\mid4,4,8,4,9,4$; singurul divizor $\equiv1\pmod q$ este $1$ (divizorii sunt $1,2,4$, respectiv $1,2,4,8$, respectiv $1,3,9$). Deci $n_q=1$ și Sylow-ul de ordin $q$ este normal.

\emph{Ordin $56=2^3\cdot7$}: $n_7\equiv1\pmod7$ și $n_7\mid8$ dau $n_7\in\{1,8\}$. Dacă $n_7=8$, cele $8$ subgrupuri de ordin $7$ se intersectează două câte două trivial (ordin prim), deci conțin $8\cdot6=48$ elemente de ordin $7$. Rămân exact $8$ elemente de ordin diferit de $7$; orice $2$-subgrup Sylow (de ordin $8$) nu conține elemente de ordin $7$, deci coincide cu mulțimea acestor $8$ elemente. Așadar $n_2=1$. În ambele cazuri un Sylow este normal.

\emph{Ordinele $12,24,36,48$}: luăm $p=2$ pentru $12,24,48$ și $p=3$ pentru $36$; $p$-subgrupurile Sylow au ordinul $4,8,16$, respectiv $9$. Dacă $n_p=1$, am terminat. Altfel fie $H\neq K$ două $p$-subgrupuri Sylow și $D=H\cap K$. Din formula produsului (Lema~\ref{lema:produsHK}), $|H|^2/|D|=|HK|\le|G|$, deci $|D|\ge|H|^2/|G|$; cum $|D|$ este un divizor propriu al lui $|H|$, obținem $|D|=2$ (ordin $12$), $|D|=4$ (ordin $24$), $|D|=8$ (ordin $48$), respectiv $|D|=3$ (ordin $36$). În fiecare caz $[H:D]=[K:D]=p$, iar un subgrup de indice $p$ al unui $p$-grup este normal în acesta (Lema~\ref{lema:normcreste}: $D\subsetneq N_H(D)$ forțează $N_H(D)=H$). Prin urmare $N=N_G(D)$ conține $H$ și $K$, deci și mulțimea $HK$, de unde $|N|\ge|HK|=|H|^2/|D|$, adică $|N|\ge8,16,32$, respectiv $27$. Cum $|N|$ este multiplu de $|H|$ și divide $|G|$, rezultă $|N|=|G|$ în toate cele patru cazuri: $D$ este subgrup normal, netrivial și propriu al lui $G$.
\end{proof}

\begin{propozitie}[$A_5$ este simplu]\label{prop:A5}
Grupul altern $A_5$ (de ordin $60$) este simplu și neabelian; el este cel mai mic grup simplu neabelian. (Aceasta explică de ce clasificarea de mai sus se oprește la $60$.)
\end{propozitie}

\begin{proof}
Determinăm clasele de conjugare din $|\mathcal C(x)|=[A_5:C_{A_5}(x)]$ (Teorema~\ref{teo:clase}). Un ciclu de lungime $5$ are $C_{S_5}(x)=\langle x\rangle$ (ordin $5$), inclus în $A_5$; deci clasa sa din $S_5$ (de $24$ elemente) se sparge în $A_5$ în două clase de câte $60/5=12$. Analog: dublele transpoziții dau o clasă de $15$, iar ciclii de lungime $3$ o clasă de $20$. Cu identitatea, cardinalele claselor sunt $1,15,20,12,12$. Un subgrup normal $N\trianglelefteq A_5$ este o reuniune de clase care conține identitatea, deci $|N|=1+(\text{sumă de termeni din }\{15,20,12,12\})$, iar $|N|$ trebuie să dividă $60$. Verificând, singurele astfel de sume care divid $60$ sunt $|N|=1$ și $|N|=1+15+20+12+12=60$ (celelalte, $16,21,13,25,36,28,33,40,45,48$, nu divid $60$). Deci $N\in\{\{e\},A_5\}$: $A_5$ este simplu. Neabelian: $(1\,2\,3)(3\,4\,5)\neq(3\,4\,5)(1\,2\,3)$.
\end{proof}

\subsubsection*{Grupuri rezolubile}

\begin{definitie}
Un grup $G$ este \emph{rezolubil} (solubil) dacă există un lanț finit de subgrupuri
\[
\{e\}=G_0\trianglelefteq G_1\trianglelefteq\dots\trianglelefteq G_m=G
\]
în care fiecare $G_{i-1}$ este normal în $G_i$, iar toate câturile $G_i/G_{i-1}$ sunt abeliene.
\end{definitie}

\begin{exemplu}
$S_4$ este rezolubil: cu $V=\{e,(1\,2)(3\,4),(1\,3)(2\,4),(1\,4)(2\,3)\}$ (grupul lui Klein) și $A_4$ (permutările pare), avem
\[
\{e\}\trianglelefteq V\trianglelefteq A_4\trianglelefteq S_4,
\]
cu câturile $V\simeq C_2\times C_2$ (abelian), $A_4/V\simeq C_3$ și $S_4/A_4\simeq C_2$ (Exemplul~\ref{ex:fundamental}(d)) --- toate abeliene.
\end{exemplu}

\begin{propozitie}[proprietăți de închidere]\label{prop:rezolubil}
Fie $G$ un grup și $N\trianglelefteq G$. Atunci:
\emph{(i)} dacă $G$ este rezolubil, orice subgrup $H\le G$ este rezolubil;
\emph{(ii)} dacă $G$ este rezolubil, $G/N$ este rezolubil;
\emph{(iii)} dacă $N$ și $G/N$ sunt rezolubile, atunci $G$ este rezolubil.
\end{propozitie}

\begin{proof}[Schiță]
(i) Intersectăm lanțul lui $G$ cu $H$: $H_i=H\cap G_i$; atunci $H_{i-1}\trianglelefteq H_i$ și $H_i/H_{i-1}$ se scufundă în $G_i/G_{i-1}$ (abelian), deci este abelian. (ii) Proiectăm lanțul prin $\pi:G\to G/N$: $\pi(G_i)$ formează un lanț cu câturi factori ai unor grupuri abeliene, deci abeliene. (iii) Concatenăm un lanț al lui $N$ cu preimaginile unui lanț al lui $G/N$ (teorema de corespondență, Teorema~\ref{teo:corespondenta}, asigură normalitatea și, prin teorema fundamentală, câturile corecte).
\end{proof}

\begin{observatie}
\label{obs:simplurezol}Un grup simplu rezolubil este obligatoriu ciclic de ordin prim: lanțul din definiție trebuie să aibă un cât $G/G_{m-1}$ abelian netrivial, forțând $G_{m-1}=\{e\}$ și $G=G_m$ abelian, deci (fiind simplu) $\simeq C_p$ (Lema~\ref{lema:abeliansimplu}).
\end{observatie}

\begin{teorema}[$S_n$ nu este rezolubil pentru $n\ge5$]\label{teo:Snrezol}
Pentru $n\ge5$, grupurile $A_n$ și $S_n$ nu sunt rezolubile.
\end{teorema}

\begin{proof}
$A_5$ este simplu și neabelian (Propoziția~\ref{prop:A5}), deci nu este $\simeq C_p$; prin Observația~\ref{obs:simplurezol}, un grup simplu rezolubil ar fi ciclic de ordin prim --- contradicție, deci $A_5$ nu este rezolubil. Pentru $n\ge5$, $A_5$ se scufundă în $A_n$ și în $S_n$ (permutările pare ale lui $\{1,\dots,5\}$, fixând $\{6,\dots,n\}$), iar un subgrup al unui grup rezolubil este rezolubil (Propoziția~\ref{prop:rezolubil}(i)). Prin urmare nici $A_n$, nici $S_n$ nu sunt rezolubile pentru $n\ge5$.
\end{proof}

\begin{teorema}[Abel--Ruffini]\label{teo:abelruffini}
Ecuația polinomială generală de grad $\ge5$ nu se rezolvă prin radicali.
\end{teorema}

Legătura cu cele de mai sus, prin teoria lui Galois: grupul Galois al polinomului general de grad $n$ este $S_n$, iar rezolvarea prin radicali a ecuației corespunde exact rezolubilității acestui grup. Nerezolubilitatea lui $S_n$ pentru $n\ge5$ (Teorema~\ref{teo:Snrezol}) dă deci concluzia.

\chapter{Inele și corpuri}

\section{Inele: definiții și reguli de calcul}

\begin{definitie}
$(A,+,\cdot)$ este \emph{inel} dacă $(A,+)$ este grup abelian (cu neutrul $0$), înmulțirea este asociativă și distributivă față de adunare: $a(b+c)=ab+ac$, $(a+b)c=ac+bc$. Inelul este \emph{unitar} dacă înmulțirea are element neutru $1$ și \emph{comutativ} dacă înmulțirea este comutativă.
\end{definitie}

\begin{observatie}[cele două convenții; convenția noastră]\label{obs:conventieinel}
În literatură circulă \textbf{două definiții} ale noțiunii de inel:
\begin{itemize}[nosep,leftmargin=1.4em]
\item \emph{fără unitate}: se cer doar axiomele de mai sus (grup abelian, asociativitate, distributivitate). Un astfel de obiect se numește uneori \emph{pseudo-inel} sau, în engleză, \emph{rng} (de la „ring'' fără „i'' de la „identity''). Exemplu: $2\ZZ$ cu operațiile uzuale --- e închis la $+$ și $\cdot$, dar nu conține $1$.
\item \emph{cu unitate}: se cere în plus existența lui $1$ cu $1\cdot a=a\cdot 1=a$.
\end{itemize}
\textbf{În programa românească --- și în tot acest volum --- „inel'' înseamnă întotdeauna \emph{inel unitar}.} Așadar $2\ZZ$ \emph{nu} este inel pentru noi (este doar un ideal al lui $\ZZ$), iar când spunem „subinel'' cerem ca acesta să conțină același $1$. Este o convenție, nu o teoremă: la olimpiadă enunțul o presupune tacit, dar merită știut că altundeva „inel'' poate însemna altceva.
\end{observatie}

\begin{definitie}[morfism de inele; morfism unitar]\label{def:morfinel}
Fie $A,B$ inele. O aplicație $f:A\to B$ se numește \emph{morfism de inele} dacă păstrează ambele operații:
\[
f(x+y)=f(x)+f(y),\qquad f(xy)=f(x)f(y)\qquad\text{pentru orice }x,y\in A.
\]
Morfismul se numește \emph{unitar} dacă, în plus,
\[
f(1_A)=1_B .
\]
Un morfism bijectiv se numește \emph{izomorfism} (inversul lui este automat morfism); un morfism $f:A\to A$ se numește \emph{endomorfism}.
\end{definitie}

\begin{observatie}[de ce $f(1)=1$ trebuie cerut separat]
Din $f(x+y)=f(x)+f(y)$ rezultă întotdeauna $f(0)=0$ și $f(-x)=-f(x)$ (proprietăți de morfism de grupuri aditive), \textbf{dar} $f(1_A)=1_B$ \emph{nu} rezultă din axiome. Contraexemplu: $f:\ZZ\to\ZZ\times\ZZ$, $f(x)=(x,0)$, păstrează adunarea și înmulțirea, însă $f(1)=(1,0)\neq(1,1)=1_{\ZZ\times\ZZ}$. De aceea, \textbf{în acest volum, prin „morfism de inele'' înțelegem implicit morfism unitar} ($f(1)=1$) --- convenția standard în programa românească; ea este esențială, de pildă, la subinelul prim (Propoziția~\ref{prop:subinelprim}) și la faptul că un morfism duce elementele inversabile în elemente inversabile: $f(u^{-1})=f(u)^{-1}$.
\end{observatie}

\begin{exemplu}
$(\ZZ,+,\cdot)$, $(\QQ,+,\cdot)$, $(\RR,+,\cdot)$, $(\CC,+,\cdot)$, $(\ZZ_n,+,\cdot)$ --- comutative; $(M_2(\RR),+,\cdot)$ --- necomutativ; $\ZZ[i]=\{a+bi\mid a,b\in\ZZ\}$ (\emph{întregii lui Gauss}) și $\ZZ[\sqrt d\,]=\{a+b\sqrt d\mid a,b\in\ZZ\}$ --- comutative; inelele de polinoame $A[X]$ (§\ref{sec:polinoame}).
\end{exemplu}

\begin{propozitie}[reguli de calcul]\label{prop:reguli}
Într-un inel $A$: \emph{(i)} $a\cdot 0=0\cdot a=0$; \emph{(ii)} $a(-b)=(-a)b=-(ab)$ și $(-a)(-b)=ab$; \emph{(iii)} $a(b-c)=ab-ac$; \emph{(iv)} distributivitatea se extinde la sume finite.
\end{propozitie}

\begin{proof}
(i) $a\cdot0=a(0+0)=a\cdot0+a\cdot0$; simplificăm în grupul $(A,+)$. (ii) $ab+a(-b)=a(b+(-b))=a\cdot0=0$, deci $a(-b)$ este opusul lui $ab$; restul analog.
\end{proof}

\begin{observatie}
Dacă într-un inel $1=0$, atunci $x=x\cdot 1=x\cdot 0=0$ pentru orice $x$, deci $A=\{0\}$. De aceea la corpuri se cere explicit $1\neq 0$.
\end{observatie}

\begin{definitie}
$U(A)=\{a\in A\mid a\text{ inversabil față de }\cdot\}$ este mulțimea \emph{unităților} inelului $A$.
\end{definitie}

\begin{propozitie}
$(U(A),\cdot)$ este grup, numit \emph{grupul unităților}.
\end{propozitie}

\begin{proof}
$1\in U(A)$; dacă $a,b\in U(A)$, atunci $(ab)(b^{-1}a^{-1})=1=(b^{-1}a^{-1})(ab)$, deci $ab\in U(A)$, iar $a^{-1}\in U(A)$ evident.
\end{proof}

\begin{exemplu}[{$U(\ZZ[i])$}]\label{ex:gauss}
$U(\ZZ)=\{\pm1\}$. În $\ZZ[i]$ definim \emph{norma} $N(a+bi)=a^2+b^2=|a+bi|^2$; ea este multiplicativă: $N(zw)=|zw|^2=N(z)N(w)$. Dacă $uv=1$ în $\ZZ[i]$, atunci $N(u)N(v)=1$ în $\NN$, deci $N(u)=1$, adică $u\in\{\pm1,\pm i\}$; reciproc, acestea sunt inversabile. Așadar $U(\ZZ[i])=\{\pm1,\pm i\}$. (Norma multiplicativă --- unealtă standard pentru inele de tip $\ZZ[\sqrt d\,]$.)
\end{exemplu}

\begin{definitie}[centrul unui inel]\label{def:centruinel}
\emph{Centrul} inelului $A$ este
\[
Z(A)=\{a\in A\mid ax=xa\ \text{pentru orice }x\in A\}.
\]
$A$ este comutativ dacă și numai dacă $Z(A)=A$.
\end{definitie}

\begin{definitie}[subinel]\label{def:subinel}
O submulțime $B\subseteq A$ se numește \emph{subinel} al lui $A$ dacă este stabilă la operațiile lui $A$ și devine ea însăși inel cu operațiile induse; concret (în convenția noastră, cu unitate):
\[
1_A\in B,\qquad x-y\in B\ \text{ și }\ xy\in B\qquad\text{pentru orice }x,y\in B.
\]
(Mulțimea $B$ este nevidă, căci $1_A\in B$; condiția $x-y\in B$ dă $0=1_A-1_A\in B$ și apoi $-y=0-y\in B$, deci $(B,+)$ este subgrup al lui $(A,+)$, prin criteriul de subgrup.) Exemple: $\ZZ\subseteq\QQ\subseteq\RR\subseteq\CC$; $\ZZ[i]\subseteq\CC$; matricele scalare în $M_n(K)$; iar $2\ZZ$ \textbf{nu} este subinel al lui $\ZZ$ (nu conține $1$).
\end{definitie}

\begin{propozitie}\label{prop:centruinel}
$Z(A)$ este un subinel al lui $A$ care conține $1$; în plus, $(Z(A),+)$ este subgrup al lui $(A,+)$. Orice element inversabil din $Z(A)$ are inversul tot în $Z(A)$.
\end{propozitie}

\begin{proof}
$0,1\in Z(A)$. Dacă $a,b\in Z(A)$, atunci pentru orice $x$: $(a-b)x=ax-bx=xa-xb=x(a-b)$ și $(ab)x=a(bx)=a(xb)=(ax)b=(xa)b=x(ab)$, deci $a-b,ab\in Z(A)$. Dacă $a\in Z(A)$ este inversabil, din $ax=xa$ rezultă (înmulțind cu $a^{-1}$ la stânga și la dreapta) $xa^{-1}=a^{-1}x$.
\end{proof}

\begin{lema}[Bézout aditiv pentru centru]\label{lema:bezoutcentru}
Fie $a\in A$ și $n,p\in\ZZ$. Dacă $na\in Z(A)$ și $pa\in Z(A)$ (multipli întregi), atunci $(n,p)\,a\in Z(A)$. În particular, dacă $(n,p)=1$, atunci $a\in Z(A)$.
\end{lema}

\begin{proof}
$(Z(A),+)$ este subgrup al lui $(A,+)$ (Propoziția~\ref{prop:centruinel}). Cu $d=(n,p)$ și $un+vp=d$ (Bézout), avem $da=u(na)+v(pa)\in Z(A)$, căci un subgrup este închis la combinații întregi ale elementelor sale. Este analogul aditiv al Lemei~\ref{lema:bezoutexp}.
\end{proof}

\begin{definitie}
Un element $a\neq0$ este \emph{divizor al lui zero} dacă există $b\neq0$ cu $ab=0$ sau $ba=0$. Un inel comutativ, unitar, cu $1\neq0$ și \textbf{fără} divizori ai lui zero se numește \emph{inel integru} (domeniu de integritate).
\end{definitie}

\begin{propozitie}\label{prop:integru}
\emph{(i)} Într-un inel integru funcționează simplificarea: $ab=ac$ și $a\neq0$ implică $b=c$. \emph{(ii)} Un element inversabil nu este divizor al lui zero (în orice inel).
\end{propozitie}

\begin{proof}
(i) $a(b-c)=0$ și $a\neq 0$ dau $b-c=0$. (ii) Dacă $ab=0$ și $a$ inversabil, atunci $b=a^{-1}(ab)=0$; analog pe partea cealaltă.
\end{proof}

\begin{propozitie}[binomul lui Newton în inele]\label{prop:binom}
Dacă $x,y\in A$ \textbf{comută} ($xy=yx$), atunci pentru orice $n\in\NN$:
\[
(x+y)^n=\sum_{k=0}^{n}\binom{n}{k}x^k y^{n-k},
\]
unde $\binom nk=C_n^k$. (Demonstrație: inducție după $n$, exact ca în $\RR$; comutativitatea este esențială pentru gruparea termenilor.)
\end{propozitie}

\subsection*{Centrul și morfismele inelelor de matrice}

\begin{propozitie}[centrul inelului de matrice]\label{prop:centrumatrice}
Fie $A$ un inel comutativ (unitar) și $n\ge1$. Atunci centrul inelului de matrice este format din matricele scalare:
\[
Z\big(M_n(A)\big)=\{a I_n\mid a\in A\},
\]
iar $a\mapsto aI_n$ este un izomorfism de inele $A\xrightarrow{\ \sim\ }Z\big(M_n(A)\big)$.
\end{propozitie}
\begin{proof}
Matricele scalare comută cu orice matrice (folosind comutativitatea lui $A$), deci $\{aI_n\}\subseteq Z(M_n(A))$. Reciproc, fie $B=(b_{k\ell})\in Z(M_n(A))$ și $E_{ij}$ matricea unitate (cu $1$ pe poziția $(i,j)$, $0$ în rest). Din $BE_{ij}=E_{ij}B$, comparând intrarea $(k,\ell)$:
\[
(BE_{ij})_{k\ell}=b_{ki}\,\delta_{j\ell},\qquad (E_{ij}B)_{k\ell}=\delta_{ik}\,b_{j\ell}.
\]
Pentru $k\neq i$, $\ell=j$: $b_{ki}=0$ --- deci $B$ e diagonală. Pentru $k=i$, $\ell=j$: $b_{ii}=b_{jj}$ --- deci toate intrările diagonale sunt egale, $B=aI_n$ cu $a=b_{11}$. Aplicația $a\mapsto aI_n$ e morfism injectiv de inele (căci $(aI_n)(bI_n)=abI_n$, $I_n\mapsto$ unitatea) și, prin cele de mai sus, surjectiv pe centru.
\end{proof}

\begin{definitie}[ideal la stânga, la dreapta, bilateral]\label{def:ideal}
Fie $A$ un inel. O submulțime $I\subseteq A$ se numește:
\begin{itemize}[nosep,leftmargin=1.4em]
\item \emph{ideal la stânga} dacă $(I,+)$ este subgrup al lui $(A,+)$ și $a x\in I$ pentru orice $a\in A$, $x\in I$;
\item \emph{ideal la dreapta} dacă $(I,+)$ este subgrup al lui $(A,+)$ și $x a\in I$ pentru orice $a\in A$, $x\in I$;
\item \emph{ideal bilateral} (pe scurt: \emph{ideal}) dacă este simultan ideal la stânga și la dreapta, adică $(I,+)\le (A,+)$ și
\[
a x\in I \quad\text{și}\quad x a\in I \qquad \text{pentru orice } a\in A,\ x\in I .
\]
\end{itemize}
Într-un inel \emph{comutativ} cele trei noțiuni coincid. Idealele $\{0\}$ și $A$ se numesc \emph{improprii}; un inel nenul fără alte ideale bilaterale se numește \emph{simplu}.
\end{definitie}

\begin{observatie}
Un ideal este mai mult decât un subinel: cere absorbție la înmulțirea cu \emph{orice} element al inelului, nu doar stabilitate internă. De pildă $\ZZ\subseteq\QQ$ este subinel, dar nu ideal ($\frac12\cdot 1\notin\ZZ$); în schimb $n\ZZ\subseteq\ZZ$ este ideal. Reciproc, un ideal propriu nu conține $1$ (altfel $a=a\cdot 1\in I$ pentru orice $a$, deci $I=A$), deci nu e subinel unitar. Atenție și la centru: $Z(A)$ este subinel (Propoziția~\ref{prop:centruinel}), dar în general \textbf{nu} este ideal.
\end{observatie}

\begin{observatie}[nucleul, câtul, idealul maximal]\label{obs:idealcat}
Idealele bilaterale sunt exact nucleele morfismelor de inele: dacă $f:A\to B$ e morfism, atunci $\Ker f$ e ideal bilateral (din $f(x)=0$ rezultă $f(ax)=f(a)f(x)=0$ și $f(xa)=0$); reciproc, pentru orice ideal bilateral $I$ mulțimea claselor $A/I=\{a+I\}$ devine inel cu operațiile
\[
(a+I)+(b+I)=(a+b)+I,\qquad (a+I)(b+I)=ab+I
\]
(bine definite tocmai datorită absorbției), numit \emph{inelul cât}, iar $\pi:A\to A/I$, $\pi(a)=a+I$, e morfism surjectiv cu $\Ker\pi=I$. Un ideal propriu $M\subsetneq A$ se numește \emph{maximal} dacă nu există ideal $J$ cu $M\subsetneq J\subsetneq A$; pentru $A$ comutativ, $M$ este maximal dacă și numai dacă $A/M$ este corp --- faptul folosit la construcția lui Kronecker ($K[X]/(p)$ cu $p$ ireductibil).
\end{observatie}

\begin{lema}[$M_n(K)$ este inel simplu]\label{lema:mnsimplu}
Fie $K$ un corp și $n\ge1$. Singurele ideale bilaterale ale lui $M_n(K)$ sunt $\{O\}$ și $M_n(K)$.
\end{lema}
\begin{proof}
Fie $I\neq\{O\}$ un ideal bilateral și $A\in I$ cu o intrare $a_{pq}\neq0$. Atunci $E_{ip}\,A\,E_{qj}=a_{pq}E_{ij}\in I$, deci $E_{ij}=a_{pq}^{-1}(a_{pq}E_{ij})\in I$ pentru orice $i,j$; sumând, $I_n\in I$, deci $I=M_n(K)$.
\end{proof}

\begin{propozitie}[morfisme între inele de matrice]\label{prop:morfmatrice}
Fie $K,L$ corpuri comutative și $m,n\in\NN^*$ cu $n>m$. Atunci nu există niciun morfism de inele (unitar, conform convenției noastre) $f:M_n(K)\to M_m(L)$. Mai mult, chiar renunțând la condiția $f(I_n)=I_m$, singura aplicație $f:M_n(K)\to M_m(L)$ care păstrează adunarea și înmulțirea este aplicația nulă.
\end{propozitie}
\begin{proof}
Demonstrăm a doua afirmație, care o implică pe prima (un morfism unitar este nenul, căci $f(I_n)=I_m\neq O$). Fie $f$ o aplicație nenulă care păstrează adunarea și înmulțirea. $\Ker f$ e ideal bilateral al lui $M_n(K)$, iar prin Lema~\ref{lema:mnsimplu} $\Ker f\in\{O,\,M_n(K)\}$; cum $f\neq0$, $\Ker f=\{O\}$, deci $f$ e \emph{injectiv}. Fie $e_{ij}=f(E_{ij})$. Atunci $e_{11},\dots,e_{nn}$ sunt idempotente ($e_{ii}^2=f(E_{ii}^2)=e_{ii}$), \emph{nenule} (injectivitate) și \emph{ortogonale} ($e_{ii}e_{jj}=f(E_{ii}E_{jj})=0$ pentru $i\neq j$). Într-un spațiu vectorial peste $L$, imaginea sumei de idempotente ortogonale este suma directă a imaginilor, deci
\[
\rang\Big(\sum_{i=1}^n e_{ii}\Big)=\sum_{i=1}^n\rang(e_{ii})\ge n
\]
(fiecare $e_{ii}\neq0$ are rang $\ge1$). Dar rangul oricărei matrice din $M_m(L)$ este $\le m$, deci $n\le m$ --- contradicție cu $n>m$.
\end{proof}

\begin{corolar}\label{cor:izomatrice}
Pentru corpuri comutative $K,L$ și $m,n\in\NN^*$: inelele $M_n(K)$ și $M_m(L)$ sunt izomorfe dacă și numai dacă $n=m$ și $K\cong L$.
\end{corolar}
\begin{proof}
$(\Leftarrow)$ e clar. $(\Rightarrow)$ Un izomorfism $M_n(K)\to M_m(L)$ e nenul, deci $n\le m$ (Propoziția~\ref{prop:morfmatrice}); inversul dă $m\le n$, deci $n=m$. Un izomorfism de inele duce centrul în centru, deci $Z(M_n(K))\cong Z(M_m(L))$; prin Propoziția~\ref{prop:centrumatrice}, $K\cong Z(M_n(K))\cong Z(M_m(L))\cong L$.
\end{proof}

\section{Inelul $\ZZ_n$. Euler, Fermat, Wilson}\label{sec:Zn}

Reamintim: $\ZZ_n=\{\cl0,\cl1,\dots,\cl{n-1}\}$, cu $\cl a+\cl b=\cl{a+b}$ și $\cl a\cdot\cl b=\cl{ab}$, operații bine definite (independente de reprezentanți, deoarece $n\mid a-a'$ și $n\mid b-b'$ implică $n\mid ab-a'b'$). $(\ZZ_n,+,\cdot)$ este inel comutativ unitar.

\begin{teorema}[structura lui $\ZZ_n$]\label{teo:Zn}
Fie $\cl a\in\ZZ_n$, $\cl a\neq\cl0$. Atunci:
\emph{(i)} $\cl a$ este inversabil $\iff(a,n)=1$;\quad
\emph{(ii)} $\cl a$ este divizor al lui zero $\iff(a,n)>1$.\\
Prin urmare, în $\ZZ_n$ orice element nenul este fie inversabil, fie divizor al lui zero.
\end{teorema}

\begin{proof}
(i) Dacă $(a,n)=1$, relația lui Bézout dă $ax+ny=1$, deci $\cl a\cdot\cl x=\cl1$. Reciproc, $\cl a\cl b=\cl1$ înseamnă $ab-1=kn$, deci orice divizor comun al lui $a$ și $n$ divide $1$.
(ii) Dacă $d=(a,n)>1$, atunci $\cl a\cdot\cl{\tfrac nd}=\cl{\tfrac ad\cdot n}=\cl 0$, cu $\cl{\tfrac nd}\neq\cl0$. Reciproc, dacă $(a,n)=1$ și $\cl a\cl b=\cl0$, atunci $n\mid ab$, deci $n\mid b$ (lema lui Gauss), adică $\cl b=\cl 0$.
\end{proof}

\begin{corolar}\label{cor:Zpcorp}
$\ZZ_n$ este corp $\iff$ $\ZZ_n$ este inel integru $\iff$ $n$ este prim.
\end{corolar}

\begin{proof}
Dacă $n=ab$ cu $1<a,b<n$, atunci $\cl a\cl b=\cl0$ dă divizori ai lui zero. Dacă $n$ este prim, orice $\cl a\neq\cl0$ are $(a,n)=1$, deci este inversabil.
\end{proof}

\begin{definitie}
$\varphi(n)$ (\emph{indicatorul lui Euler}) este numărul întregilor $k\in\{1,\dots,n\}$ cu $(k,n)=1$. Din Teorema~\ref{teo:Zn}, $|U(\ZZ_n)|=\varphi(n)$.
\end{definitie}

\begin{teorema}[Euler; Fermat]\label{teo:euler}
\emph{(i)} Dacă $(a,n)=1$, atunci $a^{\varphi(n)}\equiv1\pmod n$.
\emph{(ii)} \emph{(Mica teoremă a lui Fermat)} Dacă $p$ este prim și $p\nmid a$, atunci $a^{p-1}\equiv1\pmod p$; pentru orice $a\in\ZZ$, $a^p\equiv a\pmod p$.
\end{teorema}

\begin{proof}
(i) $\cl a$ aparține grupului $U(\ZZ_n)$, de ordin $\varphi(n)$; Corolarul~\ref{cor:lagrange}(ii) dă $\cl a^{\,\varphi(n)}=\cl1$. (ii) este cazul $n=p$, cu $\varphi(p)=p-1$; a doua formă rezultă înmulțind cu $a$, respectiv direct când $p\mid a$.
\end{proof}

\begin{teorema}[lema chineză a resturilor]\label{teo:crt}
Fie $(m,n)=1$ și $a,b\in\ZZ$. Sistemul $x\equiv a\pmod m$, $x\equiv b\pmod n$ are soluție, unică modulo $mn$.
\end{teorema}

\begin{proof}
Căutăm $x=a+mt$; condiția a doua devine $mt\equiv b-a\pmod n$, rezolvabilă deoarece $\cl m\in U(\ZZ_n)$. Unicitatea: dacă $x,y$ sunt soluții, atunci $m\mid x-y$ și $n\mid x-y$, deci $mn\mid x-y$ (căci $(m,n)=1$).
\end{proof}

\begin{corolar}[formula lui Euler]\label{cor:phi}
$\varphi$ este multiplicativă: $(m,n)=1\Rightarrow\varphi(mn)=\varphi(m)\varphi(n)$. Cum $\varphi(p^k)=p^k-p^{k-1}$ (dintre $1,\dots,p^k$ eliminăm cei $p^{k-1}$ multipli de $p$), rezultă
$\varphi(n)=n\prod_{p\mid n}\Bigl(1-\frac1p\Bigr)$.
\end{corolar}

\begin{proof}
Prin lema chineză, $x\mapsto(x\bmod m,\;x\bmod n)$ este o bijecție de la resturile modulo $mn$ la perechile de resturi; iar $(x,mn)=1\iff(x,m)=1$ și $(x,n)=1$. Numărând perechile ,,coprime`` obținem $\varphi(mn)=\varphi(m)\varphi(n)$.
\end{proof}

\begin{propozitie}[ecuația liniară în $\ZZ_n$]\label{prop:liniara}
Ecuația $\cl a\,x=\cl b$ în $\ZZ_n$ are soluții dacă și numai dacă $d=(a,n)$ divide $b$; în acest caz are exact $d$ soluții.
\end{propozitie}

\begin{proof}
Ecuația revine la $ax\equiv b\pmod n$, adică $ax+ny=b$: membrul stâng e mereu divizibil cu $d$, deci condiția e necesară. Dacă $d\mid b$, împărțim prin $d$: $a'x\equiv b'\pmod{n'}$ cu $(a',n')=1$, deci $x\equiv x_0\pmod{n'}$; modulo $n$ soluțiile sunt $x_0,\,x_0+n',\dots,x_0+(d-1)n'$, în număr de $d$.
\end{proof}

\begin{teorema}[Wilson]\label{teo:wilson}
Pentru $n\ge2$: $n$ este prim $\iff(n-1)!\equiv-1\pmod n$.
\end{teorema}

\begin{proof}
($\Rightarrow$) În corpul $\ZZ_p$, ecuația $x^2=\cl1$, adică $(x-\cl1)(x+\cl1)=\cl0$, are exact soluțiile $\pm\cl1$ (nu există divizori ai lui zero). Deci în produsul $\cl1\cdot\cl2\cdots\cl{p-1}$ elementele diferite de $\pm\cl1$ se grupează în perechi $\{x,x^{-1}\}$ cu produs $\cl1$, iar produsul total este $\cl1\cdot(-\cl1)=-\cl1$. (Pentru $p=2,3$ verificare directă.)
($\Leftarrow$) Dacă $n=de$ cu $1<d<n$, atunci $d\mid(n-1)!$ și $d\mid(n-1)!+1$, deci $d\mid1$, contradicție.
\end{proof}

\begin{observatie}
Complementar: dacă $n>4$ este compus, atunci $(n-1)!\equiv0\pmod n$. Într-adevăr, $n=ab$ cu $1<a<b<n$ face ca $a$ și $b$ să apară ca factori distincți în $(n-1)!$; iar dacă $n=p^2$ cu $p>2$, factorii $p$ și $2p$ (ambii $<n$) dau $p^2\mid(n-1)!$.
\end{observatie}

\section{Corpuri. Caracteristică. Morfismul lui Frobenius}

\begin{definitie}
Un inel unitar $K$ cu $1\neq0$ în care orice element nenul este inversabil se numește \emph{corp}; dacă înmulțirea este comutativă, \emph{corp comutativ}. Exemple: $\QQ$, $\RR$, $\CC$, $\ZZ_p$ ($p$ prim), precum și
\[
\QQ(\sqrt d\,)=\{a+b\sqrt d\mid a,b\in\QQ\}\subset\RR\quad(d\in\NN\text{ nu este pătrat perfect}),
\]
unde inversul se obține prin raționalizare: $\dfrac1{a+b\sqrt d}=\dfrac{a-b\sqrt d}{a^2-db^2}$, numitorul fiind nenul pentru $(a,b)\neq(0,0)$ deoarece $\sqrt d\notin\QQ$.
\end{definitie}

\begin{propozitie}\label{prop:corpintegru}
Un corp nu are divizori ai lui zero; în particular, orice corp comutativ este inel integru.
\end{propozitie}

\begin{proof}
Elementele nenule sunt inversabile, deci nu pot fi divizori ai lui zero (Propoziția~\ref{prop:integru}(ii)).
\end{proof}

\begin{teorema}[inel integru finit]\label{teo:integrufinit}
Orice inel integru \textbf{finit} este corp. Mai general, într-un inel finit orice element nenul care nu este divizor al lui zero este inversabil.
\end{teorema}

\begin{proof}
Fie $a\neq0$ nedivizor al lui zero. Funcțiile $x\mapsto ax$ și $x\mapsto xa$ sunt injective ($ax=ay\Rightarrow a(x-y)=0\Rightarrow x=y$), deci, $A$ fiind finit, surjective: există $b,c$ cu $ab=1$ și $ca=1$. Atunci $c=c(ab)=(ca)b=b$, deci $a$ este inversabil.
\end{proof}

\begin{propozitie}[când un inel finit este corp]\label{prop:finitcorp}
Un inel finit $A$ (cu $1\neq0$) este corp dacă și numai dacă nu are elemente nilpotente nenule și ecuația $x^2=x$ are numai soluțiile $0$ și $1$.
\end{propozitie}
\begin{proof}
$(\Rightarrow)$ Într-un corp, $a^k=0$ cu $a\neq0$ ar da $1=a^{-k}a^k=0$, exclus; iar $x^2=x\Rightarrow x(x-1)=0\Rightarrow x\in\{0,1\}$ (fără divizori ai lui zero).
$(\Leftarrow)$ Fie $a\neq0$. Cum $A$ e finit, șirul $a,a^2,a^3,\dots$ se repetă, deci există $N\ge1$, multiplu al perioadei și suficient de mare, cu $a^{2N}=a^N$; atunci $e:=a^N$ e idempotent, deci $e\in\{0,1\}$. Dacă $e=0$, $a$ ar fi nilpotent nenul --- exclus; deci $a^N=1$, adică $a\cdot a^{N-1}=a^{N-1}\cdot a=1$ și $a$ e inversabil. Orice element nenul fiind inversabil, $A$ e corp (finit, deci comutativ prin Wedderburn).
\end{proof}

\begin{propozitie}[endomorfisme ale grupului aditiv]\label{prop:endoaditiv}
Fie $A$ un inel cu $1\neq0$.
\emph{(i)} Dacă orice endomorfism nenul al grupului $(A,+)$ este injectiv, atunci $A$ nu are divizori ai lui zero.
\emph{(ii)} Dacă orice endomorfism nenul al grupului $(A,+)$ este surjectiv, atunci $A$ este corp.
\end{propozitie}
\begin{proof}
Pentru $a\in A$, translațiile $L_a(x)=ax$ și $R_a(x)=xa$ sunt endomorfisme ale lui $(A,+)$ (distributivitate).
(i) Fie $a\neq0$; $L_a$ e nenul ($L_a(1)=a\neq0$), deci injectiv. Din $ab=0=L_a(0)$ rezultă $b=0$: $a$ nu e divizor al lui zero la stânga; analog cu $R_a$ la dreapta.
(ii) Fie $a\neq0$; $L_a$ nenul, deci surjectiv: există $b$ cu $ab=1$. Cum $b\neq0$, $L_b$ e nenul, deci surjectiv: există $c$ cu $bc=1$. Atunci $a=a(bc)=(ab)c=c$, deci $ba=bc=1$: $a$ e inversabil (bilateral). Orice element nenul fiind inversabil, $A$ e corp.
\end{proof}

\begin{propozitie}\label{prop:morfcorp}
\emph{(i)} Orice morfism (unitar) de corpuri este injectiv. \emph{(ii)} Unicul morfism unitar de inele $f:\ZZ\to A$ este $f(k)=k\cdot1_A$. \emph{(iii)} Unicul endomorfism (unitar) al corpului $\QQ$ este identitatea.
\end{propozitie}

\begin{proof}
(i) Dacă $f(a)=0$ cu $a\neq0$, atunci $1=f(1)=f(a)f(a^{-1})=0$, contradicție cu $1\neq0$ în codomeniu. (ii) Din $f(1)=1_A$ și aditivitate, $f(k)=k\cdot1_A$ pentru $k\in\ZZ$; se verifică imediat că aceasta este morfism. (iii) Ca la Observația~\ref{obs:endo}(b): $f(1)=1$ forțează $f(k)=k$, apoi $f(p/q)=p/q$.
\end{proof}

\begin{definitie}
\emph{Caracteristica} $\car A$ a unui inel este cel mai mic $n\ge1$ cu $\underbrace{1+\dots+1}_{n}=n\cdot1=0$, dacă există; altfel $\car A=0$. De exemplu $\car\ZZ_n=n$, $\car\QQ=\car\RR=\car\CC=0$.
\end{definitie}

\begin{propozitie}\label{prop:caract}
\emph{(i)} Într-un inel integru (în particular, într-un corp comutativ), caracteristica este $0$ sau un număr prim. \emph{(ii)} Orice inel finit are caracteristică nenulă; deci un corp finit are caracteristică un prim $p$, și atunci $p\cdot x=0$ pentru orice $x$.
\end{propozitie}

\begin{proof}
(i) Fie $\car A=n\ge1$ și $n=ab$ cu $1\le a,b\le n$. Atunci $(a\cdot1)(b\cdot1)=ab\cdot1=0$ (distributivitate), deci $a\cdot1=0$ sau $b\cdot1=0$; minimalitatea lui $n$ dă $a=n$ sau $b=n$, adică $n$ prim (sau $n=1$, imposibil căci $1\neq0$). (ii) Într-un inel finit, șirul $1,2\cdot1,3\cdot1,\dots$ ia de două ori aceeași valoare: $m\cdot1=k\cdot1$ cu $m<k$ dă $(k-m)\cdot1=0$. Ultima afirmație: $p\cdot x=(p\cdot1)x=0$.
\end{proof}

\begin{propozitie}[subinelul prim]\label{prop:subinelprim}
Fie $A$ un inel unitar și $P(A)=\{k\cdot 1_A\mid k\in\ZZ\}$. Atunci $P(A)$ este un subinel al lui $A$, inclus în orice subinel al lui $A$ (deci este cel mai mic subinel --- \emph{subinelul prim}). Mai mult:
\[
\car(A)=0\iff P(A)\cong\ZZ,\qquad \car(A)=n\in\NN^*\iff P(A)\cong\ZZ_n.
\]
\end{propozitie}
\begin{proof}
Aplicația $\varphi:\ZZ\to A$, $\varphi(k)=k\cdot 1_A$, este morfism unitar de inele; imaginea sa este $P(A)$, care e deci subinel. Orice subinel conține $1_A$, deci conține toți $k\cdot 1_A$, adică $P(A)$: minimalitatea. Nucleul este $\Ker\varphi=\{k\in\ZZ\mid k\cdot 1_A=0\}$, subgrup al lui $(\ZZ,+)$, deci $=d\ZZ$ pentru un $d\ge0$; prin definiția caracteristicii, $d=\car(A)$ (cu $d=0$ pentru caracteristică $0$). Teorema fundamentală de izomorfism dă $P(A)=\Ima\varphi\cong\ZZ/\Ker\varphi$, adică $\ZZ$ dacă $\car(A)=0$ și $\ZZ_n$ dacă $\car(A)=n$.
\end{proof}

\begin{propozitie}[caracteristica și morfismele]\label{prop:carmorf}
\emph{(i)} Fie $A,B$ inele cu $\car(A)=n$, $\car(B)=m$. Dacă există un morfism (unitar) de inele $f:A\to B$, atunci $m\mid n$.
\emph{(ii)} Dacă $K,L$ sunt corpuri și există un morfism $f:K\to L$, atunci $\car(K)=\car(L)$.
\end{propozitie}
\begin{proof}
(i) $f(1_A)=1_B$. Din $n\cdot 1_A=0$: $0=f(n\cdot 1_A)=n\cdot 1_B$, deci $\ord(1_B)\mid n$ în $(B,+)$, adică $m\mid n$ (pentru $n=0$ afirmația e vidă). (ii) Un morfism de corpuri e injectiv (Propoziția~\ref{prop:morfcorp}(i)). Dacă $\car(K)=0$, subinelul prim $\ZZ\hookrightarrow K\hookrightarrow L$ arată $\car(L)=0$. Dacă $\car(K)=p$ prim, atunci $p\cdot 1_L=f(p\cdot 1_K)=0$, deci $\car(L)\mid p$; cum $1_L\neq0$, $\car(L)=p$.
\end{proof}

\begin{teorema}[Wedderburn, 1905]\label{teo:wedderburn}
Orice corp finit este comutativ.
\end{teorema}

Demonstrația uzuală (prin polinoame ciclotomice și ecuația claselor) depășește cadrul fișei; teorema se citează. Structural, pentru fiecare $p^{\,n}$ ($p$ prim) există, până la izomorfism, exact un corp cu $p^{\,n}$ elemente, notat $\mathbb F_{p^n}$: \emph{existența} este demonstrată în §\ref{sec:polcorpfinite} (Observația~\ref{obs:existafpn}), iar \emph{unicitatea} până la izomorfism o admitem fără demonstrație.

\begin{teorema}[cardinalul unui inel de caracteristică primă]\label{teo:cardpn}
Fie $A$ un inel finit cu $\car A=p$ prim. Atunci $|A|=p^{\,n}$ pentru un $n\ge1$.
\end{teorema}

\begin{proof}
Din $\car A=p$ avem $p\cdot x=(p\cdot1)x=0$ pentru orice $x\in A$, deci în grupul abelian $(A,+)$ orice element nenul are ordinul $p$. Definind înmulțirea cu scalari $\cl k\cdot x:=k\cdot x$ (bine definită, căci $p\cdot x=0$), $(A,+)$ devine spațiu vectorial peste corpul $\ZZ_p$. Fiind finit, are dimensiune finită $n$, deci $|A|=|\ZZ_p|^{\,n}=p^{\,n}$.
\end{proof}

\begin{corolar}
Orice corp finit $K$ are $|K|=p^{\,n}$, unde $p=\car K$ (combinând Propoziția~\ref{prop:caract}(ii) cu Teorema~\ref{teo:cardpn}). Notăm un astfel de corp $\mathbb F_{p^n}$.
\end{corolar}

\begin{teorema}[funcții polinomiale peste un corp finit]\label{teo:functiipol}
Fie $K$ un corp finit cu $|K|=q$. Atunci \textbf{orice} funcție $f:K\to K$ este polinomială: există $P\in K[X]$ cu $f(a)=P(a)$ pentru orice $a\in K$. Mai mult, există exact un astfel de $P$ de grad $\le q-1$.
\end{teorema}

\begin{proof}
Existența și unicitatea polinomului de grad $\le q-1$ sunt exact interpolarea Lagrange (Teorema~\ref{teo:lagrangeint}) în cele $q$ puncte distincte ale lui $K$, cu valorile impuse $f(a)$. (Numărătoare: funcțiile $K\to K$ sunt în număr de $q^{q}$, la fel ca polinoamele de grad $\le q-1$; corespondența polinom $\mapsto$ funcție este deci o bijecție.)
\end{proof}

\begin{propozitie}[reciproca: funcțiile polinomiale forțează corp finit]\label{prop:functiipolrecip}
Fie $A$ un inel cu $0\neq1$ în care orice funcție $f:A\to A$ este polinomială (adică $f(x)=P(x)$ pentru un $P\in A[X]$). Atunci:
\emph{(i)} $A$ este finit;
\emph{(ii)} $A$ este corp; dacă în plus $A$ este comutativ, $A$ este un corp finit.
\end{propozitie}
\begin{proof}
(i) Dacă $A$ ar fi infinit, atunci $|A[X]|=|A|$, deci mulțimea \emph{funcțiilor polinomiale} (imaginea aplicației $A[X]\to A^A$, $P\mapsto\widetilde P$) are cardinalul cel mult $|A|$. Dar mulțimea \emph{tuturor} funcțiilor $A^A$ are cardinalul $|A|^{|A|}>|A|$, deci nu orice funcție ar putea fi polinomială --- contradicție. Așadar $A$ este finit.

(ii) Considerăm funcția $e:A\to A$, $e(0)=1$ și $e(x)=0$ pentru $x\neq0$. Prin ipoteză există $P=c_0+c_1X+\dots+c_mX^m\in A[X]$ cu $e(x)=P(x)=c_0+c_1x+\dots+c_mx^m$ pentru orice $x\in A$ (coeficienții se scriu la stânga). Pentru $x=0$ obținem $c_0=e(0)=1$. Fie acum $a\neq0$. Atunci
\[
0=e(a)=P(a)=1+t\,a,\qquad\text{unde } t=c_1+c_2a+\dots+c_ma^{m-1},
\]
deci $(-t)\,a=1$: orice element nenul are un invers la stânga. De aici rezultă că orice element nenul este inversabil: dacă $a\neq0$ și $ba=1$, atunci $b\neq0$ (altfel $1=0$), deci există $c$ cu $cb=1$; atunci $c=c(ba)=(cb)a=a$, adică $ab=cb=1$, iar $b$ este inversul (bilateral) al lui $a$. Așadar $A$ este corp; fiind finit (punctul (i)), este comutativ prin teorema lui Wedderburn. Enunțul este reciproca teoremei directe de mai sus.
\end{proof}

\begin{teorema}[corpuri cu puține automorfisme]\label{teo:autcorp}
Un corp (inel cu diviziune) cu un \textbf{număr finit} de automorfisme este comutativ.
\end{teorema}

\begin{proof}
Fie $D$ un corp cu centrul $Z=Z(D)$; notăm $D^*=D\setminus\{0\}$ și $Z^*=Z\setminus\{0\}$. Presupunem, prin absurd, că $D$ nu este comutativ. Pentru $g\in D^*$, aplicația $i_g(x)=gxg^{-1}$ păstrează suma și produsul, $i_g(1)=1$, iar inversa ei este $i_{g^{-1}}$; deci $i_g$ este automorfism al inelului $D$. Ca în Teorema~\ref{teo:inn}, $g\mapsto i_g$ este morfism de grupuri de la $(D^*,\cdot)$ la $(\Aut(D),\circ)$, cu nucleul $\{g\in D^*\mid gx=xg\ \forall x\}=Z^*$; prin teorema fundamentală de izomorfism, $\Aut(D)$ conține un subgrup izomorf cu $D^*/Z^*$. Rămâne să arătăm că $D^*/Z^*$ este infinit. Presupunem contrariul.

\emph{Dacă $Z$ este finit}, atunci $D^*$ este reuniunea unui număr finit de clase $gZ^*$, fiecare cu $|Z^*|$ elemente, deci $D$ este finit și, prin teorema lui Wedderburn, comutativ: contradicție.

\emph{Dacă $Z$ este infinit}, fixăm $x\in D\setminus Z$ (există, căci $D$ nu este comutativ). Pentru orice $z\in Z$ avem $x+z\neq0$ (altfel $x=-z\in Z$), deci $x+z\in D^*$. Elementele $x+z$, $z\in Z$, sunt în număr infinit, iar clasele $gZ^*$ sunt în număr finit; deci există $z_1\neq z_2$ în $Z$ cu $x+z_1$ și $x+z_2$ în aceeași clasă, adică $x+z_1=\lambda(x+z_2)$ pentru un $\lambda\in Z^*$. Rezultă $(1-\lambda)x=\lambda z_2-z_1$. Dacă $\lambda=1$, atunci $z_1=z_2$, exclus. Altfel $1-\lambda\in Z^*$, iar inversul său este tot central (Propoziția~\ref{prop:centruinel}), deci $x=(1-\lambda)^{-1}(\lambda z_2-z_1)\in Z$: contradicție.

Așadar $D^*/Z^*$ este infinit, deci un corp necomutativ are o infinitate de automorfisme (interioare). Prin contrapunere, un corp cu un număr finit de automorfisme este comutativ.
\end{proof}

\begin{teorema}[morfismul lui Frobenius; ,,visul bobocului``]\label{teo:frobenius}
Fie $A$ un inel \textbf{comutativ} de caracteristică $p$ prim. Atunci pentru orice $x,y\in A$:
\[
(x+y)^p=x^p+y^p,\qquad (x-y)^p=x^p-y^p,
\]
iar $F:A\to A$, $F(x)=x^p$, este morfism de inele.
\end{teorema}

\begin{proof}
Cheia: $p\mid\binom pk$ pentru $1\le k\le p-1$. Într-adevăr, din identitatea $k\binom pk=p\binom{p-1}{k-1}$ rezultă $p\mid k\binom pk$, iar $(p,k)=1$ dă $p\mid\binom pk$. În binomul lui Newton (Propoziția~\ref{prop:binom}) toți termenii intermediari au coeficienți multipli de $p$, deci se anulează ($p\cdot z=0$). A doua identitate: pentru $p$ impar, $(-y)^p=-y^p$; pentru $p=2$, $-1=1$. Multiplicativitatea lui $F$ este imediată din comutativitate.
\end{proof}

\begin{aplicatie}[Fermat, a doua demonstrație]
În $\ZZ_p$: $(\cl a+\cl 1)^p=\cl a^p+\cl 1$, deci prin inducție după $a$ obținem $\cl a^p=\cl a$ pentru orice $\cl a$, adică $a^p\equiv a\pmod p$.
\end{aplicatie}

\section{Elemente nilpotente și idempotente}

\begin{definitie}[nilpotent, idempotent, inel redus, idempotent central]
Fie $A$ un inel. Un element $x\in A$ este:
\begin{itemize}[nosep,leftmargin=1.4em]
\item \emph{nilpotent} dacă $x^n=0$ pentru un $n\ge1$;
\item \emph{idempotent} dacă $x^2=x$.
\end{itemize}
Inelul $A$ se numește \emph{redus} dacă nu are elemente nilpotente nenule --- echivalent, dacă $x^2=0\Rightarrow x=0$. Un idempotent $e$ se numește \emph{central} dacă aparține centrului inelului, adică $e\in Z(A)$ (Definiția~\ref{def:centruinel}): $ex=xe$ pentru orice $x\in A$.
\end{definitie}

\begin{propozitie}\label{prop:nilpotent}
Dacă $x^n=0$, atunci $1-x$ și $1+x$ sunt inversabile, cu
\[
\begin{gathered}
(1-x)^{-1}=1+x+x^2+\dots+x^{n-1},\\
(1+x)^{-1}=1-x+x^2-\dots+(-1)^{n-1}x^{n-1}.
\end{gathered}
\]
\end{propozitie}

\begin{proof}
$(1-x)(1+x+\dots+x^{n-1})=1-x^n=1$, la fel în cealaltă ordine (puterile lui $x$ comută între ele și cu $1$). A doua formulă: înlocuim $x$ cu $-x$.
\end{proof}

\begin{propozitie}\label{prop:idempotent}
\emph{(i)} Dacă $x$ este idempotent, atunci $1-x$ este idempotent. \emph{(ii)} Într-un inel integru, singurii idempotenți sunt $0$ și $1$.
\end{propozitie}

\begin{proof}
(i) $(1-x)^2=1-2x+x^2=1-2x+x=1-x$. (ii) $x^2=x\iff x(x-1)=0$, deci $x=0$ sau $x=1$.
\end{proof}

\begin{lema}[idempotenți centrali în inele reduse]\label{teo:machale}
Într-un inel \textbf{redus} $A$, orice idempotent este central: $e^2=e\Rightarrow e\in Z(A)$.
\end{lema}

\begin{proof}
Fie $e^2=e$ și $x\in A$ arbitrar. Punem $u=ex-exe$. Folosind $e^2=e$:
\[
u=e(x-xe),\qquad (x-xe)e=xe-xe^2=xe-xe=0,
\]
deci $u^2=e(x-xe)\cdot e(x-xe)=e\bigl[(x-xe)e\bigr](x-xe)=0$. Inelul fiind redus, $u=0$, adică $ex=exe$. Simetric, pentru $v=xe-exe=(x-ex)e$ avem $e(x-ex)=ex-e^2x=0$, deci $v^2=(x-ex)\bigl[e(x-ex)\bigr]e=0$, de unde $v=0$ și $xe=exe$. Așadar $ex=exe=xe$, adică $e\in Z(A)$.
\end{proof}

\begin{observatie}
Rezultatul rămâne valabil sub ipoteza mai slabă $x^2=0\Rightarrow x\in Z(A)$ (nu neapărat $x=0$): forma generală Buckley--MacHale. Ideea ,,$(ex-exe)^2=0$`` (respectiv identitatea $xy+yx\in Z(A)$) este motorul multor teoreme de comutativitate --- vezi Teorema~\ref{teo:machalecomut} (MacHale) și ierarhia de la finalul secțiunii.
\end{observatie}

\begin{lema}[puteri într-un $n$-inel]\label{lema:ninel}
Fie $n\ge2$ fixat și $A$ un \emph{$n$-inel}, adică $x^n=x$ pentru orice $x\in A$. Atunci pentru orice $a\in A$ și orice $m\ge1$,
\[
a^{\,m(n-1)+1}=a.
\]
\end{lema}

\begin{proof}
Inducție după $m$. Pentru $m=1$: $a^{(n-1)+1}=a^n=a$. Pasul: presupunând $a^{m(n-1)+1}=a$,
\[
a^{(m+1)(n-1)+1}=a^{m(n-1)+1}\cdot a^{n-1}=a\cdot a^{n-1}=a^n=a.\qedhere
\]
\end{proof}

\begin{propozitie}[inele reduse: coborârea exponentului]\label{prop:redus}
Fie $A$ un inel \textbf{redus} și $t\ge1$ fixat. Dacă $x\in A$ satisface $x^{n+t}=x^n$ pentru un $n\ge1$, atunci
\[
x^{t+1}=x.
\]
\end{propozitie}

\begin{proof}
Fie $m\ge1$ cel mai mic exponent cu $x^{m+t}=x^m$. Presupunem, prin absurd, $m\ge2$ și considerăm
\[
y=x^{m+t-1}-x^{m-1}=x^{m-1}\bigl(x^{t}-1\bigr).
\]
Puterile lui $x$ comută între ele, deci $y^2=x^{2m-2}(x^t-1)^2=x^{2m+2t-2}-2x^{2m+t-2}+x^{2m-2}$. Înmulțind relația $x^{m+t}=x^m$ cu $x^{m-2}$ (posibil, $m\ge2$) obținem $x^{2m+t-2}=x^{2m-2}$; înmulțind-o cu $x^{m+t-2}$ obținem $x^{2m+2t-2}=x^{2m+t-2}$. Prin urmare toți cei trei termeni sunt egali cu $x^{2m-2}$ și
\[
y^2=x^{2m-2}-2x^{2m-2}+x^{2m-2}=0.
\]
Inelul fiind redus, $y=0$, adică $x^{(m-1)+t}=x^{m-1}$ --- contrazice minimalitatea lui $m$. Deci $m=1$, adică $x^{1+t}=x^1$.
\end{proof}

\begin{teorema}[condiția $x^{k+1}=x^k$ dă inel boolean]\label{teo:kk1}
Fie $A$ un inel unitar cu proprietatea: pentru orice $x\in A$ există $k=k(x)\ge1$ cu $x^{k+1}=x^k$. Atunci $A$ este redus, $x^2=x$ pentru orice $x$ (deci $A$ este boolean) și, în consecință, $A$ este comutativ.
\end{teorema}

\begin{proof}
\emph{Redus}: fie $t$ nilpotent. Atunci $1+t$ este inversabil (Propoziția~\ref{prop:nilpotent}); dar orice element inversabil $u$ satisface $u^{k+1}=u^k$, de unde, înmulțind cu $(u^{-1})^k$, $u=1$. Deci $1+t=1$, adică $t=0$: $A$ nu are nilpotenți nenuli.
\emph{Boolean}: aplicăm Propoziția~\ref{prop:redus} cu $t=1$; ipoteza $x^{k+1}=x^k$ este exact $x^{n+t}=x^n$ (cu $n=k$), deci $x^{2}=x$ pentru orice $x$.
\emph{Comutativ}: Teorema~\ref{teo:boole}.
\end{proof}



\begin{teorema}[inele booleene]\label{teo:boole}
Dacă $x^2=x$ pentru \textbf{orice} $x\in A$, atunci $1+1=0$ și $A$ este comutativ.
\end{teorema}

\begin{proof}
Pentru orice $x$: $(x+x)=(x+x)^2=x^2+x^2+x^2+x^2=4x$, deci $2x=0$, adică fiecare element este propriul opus. Apoi
$x+y=(x+y)^2=x^2+xy+yx+y^2=x+xy+yx+y$, de unde $xy+yx=0$, deci $xy=-yx=yx$.
\end{proof}

\begin{corolar}\label{cor:boolefarad}
Un inel boolean fără divizori ai lui zero are cel mult două elemente (este $\{0\}$ sau $\simeq\ZZ_2$).
\end{corolar}

\begin{proof}
$x^2=x\iff x(x-1)=0$; fără divizori ai lui zero, $x=0$ sau $x=1$, deci $A\subseteq\{0,1\}$.
\end{proof}

\begin{propozitie}[cardinalul unui inel boolean]\label{prop:boolcard}
Pe o mulțime finită cu $m$ elemente se poate defini o structură de inel boolean (unitar, $x^2=x$ pentru orice $x$) dacă și numai dacă $m=2^k$ pentru un $k\in\NN^*$.
\end{propozitie}
\begin{proof}
$(\Rightarrow)$ Într-un inel boolean $1+1=0$, deci $\car(A)=2$ și $2\cdot x=0$ pentru orice $x$; atunci $(A,+)$ devine spațiu vectorial peste $\ZZ_2$ (cu $\cl k\cdot x:=k\cdot x$, bine definit). Fiind finit de dimensiune $k$, $|A|=2^k$.
$(\Leftarrow)$ Inelul produs $\ZZ_2^{\,k}$ (operații pe componente) este boolean și are $2^k$ elemente.
\end{proof}

\begin{observatie}
Se mai poate arăta (rezultat mai fin, care se citează) că un \emph{$n$-inel} ($x^n=x$ pentru orice $x$) este boolean dacă este îndeplinită una dintre condițiile: \emph{(i)} există $k>m$ naturale și $p$ prim cu $n=2^k+2^m=p+1$; \emph{(ii)} $x+x=0$ pentru orice $x$ și există $k$ cu $n=2^k+1$.
\end{observatie}

\subsection*{Teoreme de comutativitate}

Următoarele sunt \textbf{teoreme clasice de comutativitate}. Demonstrațiile lor generale depășesc programa și le enunțăm fără demonstrație; cazurile mici sunt însă elementare: pentru $x^2=x$ vezi Teorema~\ref{teo:boole}, iar pentru $x^3=x$ vezi demonstrația de după Teorema~\ref{teo:jacobson}. Fac excepție teorema lui MacHale și varianta ei cubică (Propoziția~\ref{prop:machalegen}), demonstrate complet aici, și cazurile deja tratate. Fiecare enunț îl generalizează, în esență, pe cel dinaintea sa.

\begin{teorema}[MacHale]\label{teo:machalecomut}
Fie $A$ un inel și $f:A\to A$ un endomorfism \textbf{surjectiv} cu proprietatea
\[
x^2-f(x)\in Z(A)\qquad\text{pentru orice }x\in A.
\]
Atunci $A$ este comutativ.
\end{teorema}

\begin{proof}
Notăm $g(x)=x^2-f(x)\in Z(A)$. Cum $f$ este aditiv, $f(x+y)=f(x)+f(y)$, iar $(x+y)^2-x^2-y^2=xy+yx$; prin urmare
\[
g(x+y)-g(x)-g(y)=(x+y)^2-x^2-y^2=xy+yx.
\]
Membrul stâng aparține lui $Z(A)$ (grup aditiv), deci
\[
xy+yx\in Z(A)\qquad\text{pentru orice }x,y\in A.\tag{$\star$}
\]
Fie $s=xy+yx\in Z(A)$; fiind central, $s$ comută cu $x$: din $xs=sx$ obținem
\[
x^2y+xyx=xyx+yx^2\ \Longrightarrow\ x^2y=yx^2\quad(\forall y),
\]
deci $x^2\in Z(A)$ pentru orice $x$ (toate pătratele sunt centrale). Atunci
\[
f(x)=x^2-g(x)\in Z(A)\qquad(\text{diferență de elemente centrale}).
\]
În fine, $f$ fiind \textbf{surjectiv}, orice element al lui $A$ este de forma $f(x)$, deci aparține lui $Z(A)$: rezultă $A=Z(A)$, adică $A$ este comutativ.
\end{proof}

\begin{observatie}
Demonstrația folosește doar \emph{aditivitatea} și \emph{surjectivitatea} lui $f$ (nu și multiplicativitatea). Pasul-cheie $(\star)$ --- ,,$xy+yx$ este central`` --- forțează deja \emph{toate pătratele să fie centrale}; de aici, surjectivitatea transferă centralitatea pe tot inelul. Un caz particular des întâlnit la concurs: $f=1_A$ (identitatea, evident surjectivă), adică ipoteza $x^2-x\in Z(A)$ pentru orice $x$ implică $A$ comutativ (Herstein cu $n=2$).
\end{observatie}

\begin{propozitie}[o variantă cubică a teoremei lui MacHale]\label{prop:machalegen}
Fie $A$ un inel cu proprietatea că $x^3+x^2\in Z(A)$ pentru orice $x\in A$. Atunci $A$ este comutativ. Aceeași concluzie are loc dacă $x^3-x^2\in Z(A)$ pentru orice $x\in A$: cele două ipoteze sunt echivalente, căci înlocuind $x$ cu $-x$ obținem $(-x)^3+(-x)^2=-(x^3-x^2)$, iar $Z(A)$ este subgrup aditiv.
\end{propozitie}
\begin{proof}
Notăm $h(x)=x^3+x^2$. Folosim că $Z(A)$ este un subinel care conține $1$ (Propoziția~\ref{prop:centruinel}), deci este închis la sume și diferențe și conține toți multiplii întregi $k\cdot1$. Calculele de mai jos implică doar puteri ale unui singur element și constante întregi, care comută între ele.

\emph{Pasul 1.} $h(-x)=-x^3+x^2\in Z(A)$, deci $h(x)+h(-x)=2x^2\in Z(A)$.

\emph{Pasul 2.} $h(x+1)=(x+1)^3+(x+1)^2=x^3+4x^2+5x+2$, deci $h(x+1)-h(x)-2=3x^2+5x\in Z(A)$. Scăzând $2x^2$ (Pasul 1), obținem
\[
u(x):=x^2+5x\in Z(A)\qquad\text{pentru orice }x\in A.
\]

\emph{Pasul 3.} $u(x+1)-u(x)=(x^2+7x+6)-(x^2+5x)=2x+6\in Z(A)$, deci $2x\in Z(A)$; apoi $x^2+x=u(x)-2\cdot(2x)\in Z(A)$ pentru orice $x\in A$.

\emph{Pasul 4 (mecanismul lui MacHale).} Aplicând Pasul 3 lui $x$, $y$ și $x+y$,
\[
(x+y)^2+(x+y)-(x^2+x)-(y^2+y)=xy+yx\in Z(A).
\]
În particular $xy+yx$ comută cu $x$: $x^2y+xyx=xyx+yx^2$, adică $x^2y=yx^2$ pentru orice $y$, deci $x^2\in Z(A)$. În fine, $x=(x^2+x)-x^2\in Z(A)$ pentru orice $x$, deci $A$ este comutativ.
\end{proof}

\begin{observatie}
O formulare ,,mai generală``, cu $f(x^3)-g(x^2)\in Z(A)$ (sau $f(x)-g(x^3)\in Z(A)$) pentru endomorfisme surjective \emph{arbitrare} $f,g$ ale grupului $(A,+)$, este \textbf{falsă}. În inelul necomutativ $A$ al matricelor $2\times2$ superior triunghiulare peste $\ZZ_2$ (cu centrul $\{O,I_2\}$), un calcul direct pe cele $8$ elemente arată că $x^3-x^2$ și $x-x^3$ aparțin mulțimii $\{O,E_{12}\}$ pentru orice $x$. Luând o bijecție aditivă $f$ a lui $A$ cu $f(E_{12})=I_2$ (există, $(A,+)$ fiind spațiu vectorial de dimensiune $3$ peste $\ZZ_2$) și $g=f$, obținem $f(x^3)-g(x^2)=f(x^3-x^2)\in\{O,I_2\}\subseteq Z(A)$ și, analog, $f(x)-g(x^3)\in Z(A)$, deși $A$ nu este comutativ.
\end{observatie}

\begin{teorema}[Jacobson]\label{teo:jacobson}
Dacă pentru orice $x\in A$ există un întreg $n=n(x)>1$ cu $x^{\,n(x)}=x$, atunci $A$ este comutativ.
\end{teorema}

Inelele cu această proprietate se numesc \emph{potente}. Cazul $n$ fixat (orice inel cu $x^n=x$ pentru un $n\ge2$ fix este comutativ) conține, pentru $n=2$, Teorema~\ref{teo:boole}, iar pentru $n=3$ se demonstrează elementar, astfel. Fie $A$ cu $x^3=x$ pentru orice $x$. \emph{(1)} $A$ este redus: din $x^2=0$ rezultă $x=x^3=x\cdot x^2=0$. \emph{(2)} Pentru orice $x$, $x^2$ este idempotent ($x^4=x\cdot x^3=x^2$), deci central (Lema~\ref{teo:machale}). \emph{(3)} $(x+1)^2=x^2+2x+1$ este central (ca pătrat), iar $x^2$ și $1$ sunt centrale, deci $2x\in Z(A)$. \emph{(4)} Din $(x+1)^3=x+1$ și $x^3=x$ rezultă $3x^2+3x=0$, deci $3x=-3x^2\in Z(A)$. \emph{(5)} $x=3x-2x\in Z(A)$ pentru orice $x$: $A$ este comutativ. Teorema lui Jacobson generalizează teorema lui Wedderburn.

\begin{teorema}[Herstein]\label{teo:herstein}
Dacă pentru orice $x\in A$ există $n=n(x)>1$ astfel încât $x^{\,n(x)}-x$ să fie central (adică în $Z(A)$), atunci $A$ este comutativ.
\end{teorema}

Este o generalizare a teoremei lui Jacobson: dacă $x^{n(x)}=x$, atunci $x^{n(x)}-x=0\in Z(A)$. Reciproca teoremei este banală (într-un inel comutativ orice element este central), deci condiția lui Herstein caracterizează exact inelele comutative.

\begin{teorema}[Kaplansky]\label{teo:kaplansky}
Fie $D$ un corp (inel cu diviziune, nu neapărat comutativ) cu centrul $Z=Z(D)$.
\emph{(i)} Dacă pentru orice $x\in D$ o putere $x^{\,n(x)}$ ($n(x)\ge1$) aparține lui $Z$, atunci $D$ este comutativ.
\emph{(ii)} Dacă grupul cât $D^*/Z^*$ este de torsiune (orice element are ordin finit), atunci $D$ este comutativ.
\end{teorema}

\begin{teorema}[exponenți de parități diferite]\label{teo:parity}
Dacă pentru orice $x\in A$ există întregi $n(x)>m(x)>1$ \textbf{de parități diferite} cu $x^{\,n(x)}=x^{\,m(x)}$, atunci $A$ este comutativ.
\end{teorema}

Aceasta generalizează Propoziția~\ref{prop:redus} (coborârea exponentului în inele reduse), fără a mai cere ca inelul să fie redus. Ipoteza de paritate este esențială: fără ea există contraexemple necomutative.

\begin{teorema}[Nicholson--Yaqub, 1979]\label{teo:nicholson}
Fie $j,k$ două numere naturale \textbf{prime între ele} și $A$ un inel (sau $G$ un grup) în care puterile de exponent $j$, respectiv $k$, comută două câte două:
\[
x^{\,j}y^{\,j}=y^{\,j}x^{\,j}\quad\text{și}\quad x^{\,k}y^{\,k}=y^{\,k}x^{\,k}\qquad\text{pentru orice }x,y.
\]
Atunci $A$ este comutativ (respectiv $G$ este abelian).
\end{teorema}

Rezultatul face parte din familia teoremelor de comutativitate ,,prin aplicația de ridicare la putere`` (W.\,K.~Nicholson, A.~Yaqub, \emph{Canad. Math. Bull.} \textbf{22} (1979); vezi și H.\,E.~Bell, 1978). Dacă unul dintre $j,k$ este $1$, concluzia este banală; interesul este pentru $j,k\ge2$.

\begin{propozitie}[trucul $1-ab$ / $1-ba$]\label{prop:halmos}
Într-un inel $A$, dacă $1-ab$ este inversabil, atunci $1-ba$ este inversabil și
\[
(1-ba)^{-1}=1+b\,(1-ab)^{-1}a.
\]
\end{propozitie}

\begin{proof}
Fie $u=(1-ab)^{-1}$, deci $u-abu=u-uab=1$. Punem $v=1+bua$ și verificăm:
\[
(1-ba)v=1-ba+bua-b(ab\,u)a=1-ba+b(u-abu)a=1-ba+ba=1,
\]
\[
v(1-ba)=1-ba+bua-bu(ab)a=1-ba+b(u-uab)a=1-ba+ba=1.\qedhere
\]
\end{proof}

\begin{observatie}
Formula pare ,,scoasă din pălărie``, dar provine din dezvoltarea formală $(1-ba)^{-1}=1+ba+baba+\dots=1+b(1+ab+abab+\dots)a$; la concurs se prezintă doar verificarea, care este riguroasă.
\end{observatie}

\subsection*{Corpul $U(A)\cup\{0\}$ (S. Rădulescu, M. Piticaru)}

Grupăm aici câteva rezultate de tip olimpiadă despre inele finite în care unitățile plus $0$ formează un corp; ca și celelalte teoreme grele de mai sus (Wedderburn, Jacobson, Kaplansky), se citează.

\begin{propozitie}\label{prop:rp1}
Fie $A$ un inel integru de caracteristică $p\in\NN^*$ cu un număr \emph{finit} de elemente inversabile. Atunci operațiile lui $A$ induc pe $L=U(A)\cup\{0\}$ o structură de corp.
\end{propozitie}

\begin{propozitie}\label{prop:rp2}
Dacă $A$ este un inel finit astfel încât $\big(U(A)\cup\{0\},+,\cdot\big)$ este corp, atunci $A$ nu are elemente nilpotente nenule.
\end{propozitie}

\begin{propozitie}[S. Rădulescu, M. Piticaru]\label{prop:rp3}
Fie $A$ un inel finit de caracteristică $p\ge3$ (număr impar). Dacă $\big(U(A)\cup\{0\},+,\cdot\big)$ este corp, atunci $A$ este corp.
\end{propozitie}

\begin{propozitie}[S. Rădulescu, M. Piticaru]\label{prop:rp4}
Fie $A$ un inel finit de caracteristică $p\ge3$ (număr impar). Dacă pentru orice $x\in U(A)$ avem $1+x\in U(A)\cup\{0\}$, atunci $A$ este corp.
\end{propozitie}

\section{Polinoame peste corpuri finite}\label{sec:polcorpfinite}

Fixăm un prim $p$. Studiem inelul de polinoame $K[X]$ peste un corp comutativ $K$ de caracteristică $p$ (în particular $K=\ZZ_p$), unde apar fenomene specifice caracteristicii, în care morfismul lui Frobenius (Teorema~\ref{teo:frobenius}) joacă rolul principal. \emph{Convenție}: un polinom $f\in K[X]$ de grad $\ge1$ este \emph{ireductibil} dacă nu se scrie ca produs $f=gh$ cu $\deg g,\deg h\ge1$; teoria generală a lui $K[X]$ (grad, împărțire cu rest, factorizare unică în ireductibile, derivată formală) este dezvoltată în Partea~III și o folosim aici ca atare.

\begin{teorema}[Artin--Schreier]\label{teo:artinschreier}
Fie $K$ un corp comutativ de caracteristică $p$ și $a\in K$. Polinomul \emph{Artin--Schreier}
\[
f=X^p-X-a
\]
este ireductibil în $K[X]$ \textbf{dacă și numai dacă} $a$ nu este de forma $\lambda^p-\lambda$ cu $\lambda\in K$.
\end{teorema}

\begin{proof}
Dacă $a=\lambda^p-\lambda$ cu $\lambda\in K$, atunci $f(\lambda)=0$, deci $(X-\lambda)\mid f$ și $f$ este reductibil.

Reciproc, presupunem $f$ reductibil și arătăm că $a=\lambda^p-\lambda$ pentru un $\lambda\in K$. Observăm mai întâi, prin Frobenius pe inelul $K[X]$ (care are tot caracteristica $p$),
\[
f(X+1)=(X+1)^p-(X+1)-a=(X^p+1)-X-1-a=X^p-X-a=f(X).
\]
Cum $f'=pX^{p-1}-1=-1$, un eventual factor pătratic $g^2\mid f$ ar da $g\mid f'=-1$, imposibil pentru $\deg g\ge1$; deci $f$ este \textbf{liber de pătrate}: $f=g_1g_2\cdots g_r$, cu $g_j$ ireductibile monice \emph{distincte} (Teorema~\ref{teo:descompunere}), iar $r\ge2$ prin reductibilitate.

Substituția $\sigma:h(X)\mapsto h(X+1)$ este un automorfism al lui $K[X]$ care fixează $f$, deci \emph{permută} factorii $\{g_1,\dots,g_r\}$ (păstrându-le gradul). Cum $p\cdot1=0$ în $K$, avem $\sigma^p=\mathrm{id}$ ($h(X+p\cdot1)=h(X)$), deci fiecare orbită a lui $\sigma$ pe mulțimea factorilor are cardinalul $1$ sau $p$.

\emph{Nu pot fi toate orbitele de cardinal $1$.} Dacă $\sigma$ ar fixa un factor $g$ de grad $d$, adică $g(X+1)=g(X)$, comparând coeficientul lui $X^{d-1}$ în cei doi membri ($g=X^d+c_{d-1}X^{d-1}+\cdots$ dă, în $g(X+1)$, coeficientul $d+c_{d-1}$) obținem $d=0$ în $K$, adică $p\mid d$; imposibil pentru $1\le d\le\deg f-1<p$.

Rămâne că există o orbită de cardinal $p$: ea conține $p$ factori distincți de același grad $d\ge1$, contribuind cu $pd$ la $\deg f=p$; deci $d=1$ și acei $p$ factori epuizează $f$. Așadar $f$ se descompune în factori liniari peste $K$; în particular are o rădăcină $\lambda\in K$, iar $f(\lambda)=0$ dă $a=\lambda^p-\lambda$.
\end{proof}

\begin{observatie}
Aplicația $\wp:K\to K$, $\wp(\lambda)=\lambda^p-\lambda$, este morfism al grupului aditiv (din nou Frobenius), numită \emph{operatorul Artin--Schreier}; teorema spune că $X^p-X-a$ este ireductibil exact când $a\notin\Ima\wp$. Pe corpul prim $\ZZ_p$ avem $\wp\equiv0$ (Fermat: $\lambda^p=\lambda$), deci $\Ima\wp=\{0\}$: peste $\ZZ_p$, polinomul $X^p-X-a$ este ireductibil pentru \emph{orice} $a\neq0$ (și reductibil pentru $a=0$, căci $X^p-X=\prod_{c\in\ZZ_p}(X-c)$, iarăși prin Fermat).
\end{observatie}

Pentru rezultatele structurale care urmează folosim o construcție standard: dat un polinom monic ireductibil $f\in\ZZ_p[X]$ de grad $d$, \emph{clasele de resturi ale polinoamelor modulo $f$} formează un inel comutativ, notat $\ZZ_p[X]/(f)$ (analog cu $\ZZ_n=\ZZ/n\ZZ$). Fiecare clasă are un unic reprezentant de grad $<d$, deci inelul are $p^{\,d}$ elemente; iar dacă $f$ este ireductibil, orice rest nenul $r$ este inversabil (din $(r,f)=1$ și Bézout polinomial), deci $\ZZ_p[X]/(f)$ este \textbf{corp} --- un corp cu $p^{\,d}$ elemente, în care clasa $\alpha=\widehat X$ este o rădăcină a lui $f$.

\begin{teorema}[ireductibile și $X^{p^n}-X$]\label{teo:xpnx}
Fie $f\in\ZZ_p[X]$ monic ireductibil de grad $d$. Atunci $f\mid X^{p^{\,d}}-X$. Prin urmare, pentru orice polinom ireductibil $f$ există $n$ (anume $n=\deg f$) cu $f\mid X^{p^{\,n}}-X$.
\end{teorema}

\begin{proof}
Fie $L=\ZZ_p[X]/(f)$, corp cu $p^{\,d}$ elemente, și $\alpha=\widehat X\in L$ rădăcină a lui $f$. Grupul multiplicativ $L^*$ are ordinul $p^{\,d}-1$, deci prin Lagrange $\beta^{\,p^{\,d}-1}=1$ pentru orice $\beta\neq0$; adică $\beta^{\,p^{\,d}}=\beta$ pentru \emph{orice} $\beta\in L$ (inclusiv $\beta=0$). În particular $\alpha^{\,p^{\,d}}=\alpha$, deci $\alpha$ este rădăcină a lui $X^{p^{\,d}}-X$. Dar evaluarea $g\mapsto g(\alpha)$ este morfismul de trecere la cât $\ZZ_p[X]\to L$, cu nucleul $(f)$; din $\bigl(X^{p^{\,d}}-X\bigr)(\alpha)=0$ rezultă $X^{p^{\,d}}-X\in(f)$, adică $f\mid X^{p^{\,d}}-X$.
\end{proof}

\begin{propozitie}[subcorpurile unui corp finit]\label{prop:subcorpuri}
Fie $p$ prim, $n\ge1$ și $K=\mathbb F_{p^n}$ un corp finit cu $p^n$ elemente. Pentru fiecare divizor $d\mid n$, $K$ conține \emph{exact un} subcorp cu $p^d$ elemente, anume
\[
K_d=\{x\in K\mid x^{p^d}=x\},
\]
iar acestea sunt \emph{toate} subcorpurile lui $K$.
\end{propozitie}
\begin{proof}
\emph{Orice subcorp are ordinul $p^d$ cu $d\mid n$:} un subcorp $L$ are caracteristica $p$, deci $|L|=p^d$; iar $K$ e spațiu vectorial peste $L$, deci $p^n=(p^d)^{\dim_L K}$, de unde $d\mid n$.

\emph{Existența și forma:} din $d\mid n$ rezultă $p^d-1\mid p^n-1$, deci $X^{p^d-1}-1\mid X^{p^n-1}-1$ și, înmulțind cu $X$, $X^{p^d}-X\mid X^{p^n}-X$. Cele $p^n$ elemente ale lui $K$ sunt exact rădăcinile lui $X^{p^n}-X$ (grupul $K^*$ are ordin $p^n-1$, deci $x^{p^n}=x$ pentru orice $x$). Prin urmare $X^{p^d}-X$ are toate cele $p^d$ rădăcini distincte în $K$, iar mulțimea lor $K_d=\{x:x^{p^d}=x\}$ este un subcorp (mulțimea punctelor fixe ale iteratei $d$-a a morfismului lui Frobenius $x\mapsto x^p$ este închisă la $+,\cdot$ și la inverse). Deci $|K_d|=p^d$.

\emph{Unicitatea:} orice subcorp $L$ cu $p^d$ elemente satisface $x^{p^d}=x$ pentru orice $x\in L$ (grupul $L^*$ are ordin $p^d-1$), deci $L\subseteq K_d$; egalitatea cardinalelor dă $L=K_d$.
\end{proof}

\begin{lema}[factorii ireductibili ai lui $X^{p^n}-X$]\label{lema:ireddivide}
Fie $n\ge1$ și $f\in\ZZ_p[X]$ monic ireductibil de grad $d$. Atunci $f\mid X^{p^{\,n}}-X$ dacă și numai dacă $d\mid n$.
\end{lema}

\begin{proof}
($\Leftarrow$) Prin Teorema~\ref{teo:xpnx}, $f\mid X^{p^{\,d}}-X$. Din $d\mid n$ rezultă $p^d-1\mid p^n-1$ (căci $p^n-1=(p^d)^{n/d}-1$), deci $X^{p^d-1}-1\mid X^{p^n-1}-1$ și, înmulțind cu $X$, $X^{p^{\,d}}-X\mid X^{p^{\,n}}-X$.

($\Rightarrow$) Fie $L=\ZZ_p[X]/(f)$, corp comutativ cu $p^{\,d}$ elemente, și $\alpha=\widehat X$; din $f\mid X^{p^{\,n}}-X$ rezultă $\alpha^{p^n}=\alpha$. Morfismul lui Frobenius $F(\beta)=\beta^p$ este morfism de inele al lui $L$ (Teorema~\ref{teo:frobenius}), deci și $F^n(\beta)=\beta^{p^n}$ este; prin urmare mulțimea $S=\{\beta\in L\mid\beta^{p^n}=\beta\}$ este închisă la adunare și înmulțire. Ea conține clasele constantelor din $\ZZ_p$ (Fermat: $c^p=c$) și pe $\alpha$; cum orice element al lui $L$ este de forma $c_0+c_1\alpha+\dots+c_{d-1}\alpha^{d-1}$ cu $c_i\in\ZZ_p$, rezultă $S=L$, adică $F^n=\mathrm{id}_L$. Pe de altă parte, $\beta^{p^d}=\beta$ pentru orice $\beta\in L$ (grupul $L^*$ are ordinul $p^d-1$), adică $F^d=\mathrm{id}_L$. Deci $F$ este o bijecție a lui $L$, iar ordinul ei în grupul bijecțiilor lui $L$ divide atât $d$, cât și $n$, deci divide $g=(d,n)$: $F^g=\mathrm{id}_L$. Așadar toate cele $p^d$ elemente ale lui $L$ sunt rădăcini ale polinomului $X^{p^g}-X$, care are cel mult $p^g$ rădăcini în corpul $L$ (Teorema~\ref{teo:numarradacini}). Rezultă $p^d\le p^g$, deci $d\le g$, adică $d=g\mid n$.
\end{proof}

\begin{teorema}[existența ireductibilelor de orice grad]\label{teo:existaired}
Pentru $d\ge1$ notăm $I(d)$ numărul polinoamelor monice ireductibile de grad $d$ din $\ZZ_p[X]$. Atunci, pentru orice $n\ge1$,
\[
p^{\,n}=\sum_{d\mid n}d\cdot I(d),
\]
și $I(n)\ge1$: există un polinom ireductibil de grad $n$ în $\ZZ_p[X]$.
\end{teorema}

\begin{proof}
Polinomul $h=X^{p^{\,n}}-X$ este monic și \emph{liber de pătrate}: derivata sa este $p^{\,n}X^{p^n-1}-1=-1$, iar dacă $h=g^2k$ cu $\deg g\ge1$, atunci $h'=g\,(2g'k+gk')$ ar fi divizibil cu $g$, imposibil pentru $h'=-1$. Prin descompunerea în factori ireductibili (Teorema~\ref{teo:descompunere}), $h$ este produsul unor polinoame monice ireductibile distincte, iar prin Lema~\ref{lema:ireddivide} acestea sunt exact polinoamele monice ireductibile de grad $d\mid n$. Comparând gradele, $p^{\,n}=\sum_{d\mid n}d\cdot I(d)$.

Aplicând formula pentru fiecare $d$ obținem $d\cdot I(d)\le p^{\,d}$. Divizorii $d<n$ ai lui $n$ satisfac $d\le n/2$, deci, cu $m=\lfloor n/2\rfloor$,
\[
n\cdot I(n)=p^{\,n}-\sum_{\substack{d\mid n\\ d<n}}d\cdot I(d)\ \ge\ p^{\,n}-\sum_{k=1}^{m}p^{\,k}=p^{\,n}-\frac{p^{\,m+1}-p}{p-1}\ >\ p^{\,n}-p^{\,m+1}\ \ge\ 0,
\]
unde am folosit $\frac{p^{m+1}-p}{p-1}\le p^{m+1}-p<p^{m+1}$ (căci $p-1\ge1$) și $m+1\le n$. Așadar $I(n)\ge1$.
\end{proof}

\begin{observatie}\label{obs:existafpn}
\emph{(a)} Dacă $f\in\ZZ_p[X]$ este ireductibil de grad $n$ (există, prin Teorema~\ref{teo:existaired}), atunci $\ZZ_p[X]/(f)$ este un corp cu $p^{\,n}$ elemente (construcția de mai sus). Aceasta demonstrează \emph{existența} corpului $\mathbb F_{p^n}$ pentru orice $p$ prim și orice $n\ge1$.
\emph{(b)} Aplicând inversiunea Möbius (Propoziția~\ref{prop:mobius}(ii)) relației $p^{\,n}=\sum_{d\mid n}d\cdot I(d)$, cu $M_d=d\cdot I(d)$ și $N_n=p^{\,n}$, obținem \emph{formula lui Gauss}
\[
I(n)=\frac1n\sum_{d\mid n}\mu(d)\,p^{\,n/d}.
\]
De exemplu, peste $\ZZ_2$ există $I(4)=\frac14(2^4-2^2)=3$ polinoame ireductibile de grad $4$, iar peste $\ZZ_3$ există $I(2)=\frac12(3^2-3)=3$ polinoame monice ireductibile de grad $2$.
\end{observatie}

\chapter{Polinoame peste un corp}

\section{Polinoame: grad, împărțire cu rest, divizibilitate}\label{sec:polinoame}

Fie $A$ un inel comutativ (de regulă un corp $K$). $A[X]$ este inelul polinoamelor cu coeficienți în $A$: $f=a_nX^n+\dots+a_1X+a_0$. Dacă $a_n\neq0$, \emph{gradul} este $\deg f=n$, $a_n$ este \emph{coeficientul dominant} ($f$ este \emph{unitar} sau \emph{monic} dacă $a_n=1$), iar $a_0$ este \emph{termenul liber}. Convenim $\deg 0=-\infty$.

\begin{propozitie}[gradele]\label{prop:grade}
\emph{(i)} $\deg(f+g)\le\max(\deg f,\deg g)$, cu egalitate dacă $\deg f\neq\deg g$. \emph{(ii)} Dacă $A$ este integru, $\deg(fg)=\deg f+\deg g$. \emph{(iii)} Dacă $A$ este integru, atunci $A[X]$ este integru și $U(A[X])=U(A)$; în particular, pentru un corp $K$, $U(K[X])=K^*$ (constantele nenule).
\end{propozitie}

\begin{proof}
(ii) Coeficientul lui $X^{\deg f+\deg g}$ în $fg$ este produsul coeficienților dominanți, nenul. (iii) $fg=0$ cu $f,g\neq0$ ar contrazice (ii); iar $fg=1$ forțează $\deg f=\deg g=0$, deci $f,g\in U(A)$.
\end{proof}

\begin{observatie}[{construcția riguroasă a lui $R[X]$}]
Inelul $R[X]$ se poate construi \emph{formal} ca mulțimea șirurilor $(a_0,a_1,a_2,\dots)$ de elemente din $R$ cu doar un număr finit de termeni nenuli, cu adunarea pe componente și înmulțirea dată de produsul Cauchy
\[
(a_i)\cdot(b_j)=(c_k),\qquad c_k=\sum_{i+j=k}a_ib_j.
\]
Notând $X=(0,1,0,0,\dots)$ și identificând $a\in R$ cu $(a,0,0,\dots)$, orice șir se scrie unic $(a_0,\dots,a_n,0,\dots)=\sum_{i=0}^n a_iX^i$; astfel „expresia'' $\sum a_iX^i$ capătă sens riguros ca element al acestui inel.
\end{observatie}

\begin{definitie}[ordinul unui polinom]
Pentru $f=\sum_{i=0}^n a_iX^i\in R[X]$ nenul, \emph{ordinul} $\ord f$ este cel mai mic $j$ cu $a_j\neq0$ (cel mai mic exponent care apare efectiv). Convenim $\ord 0=+\infty$.
\end{definitie}

\begin{propozitie}[ordinul: sumă și produs]\label{prop:ordinpol}
Pentru $f,g\in R[X]$:
\[
\ord(f+g)\ge\min\{\ord f,\ord g\},\qquad \ord(fg)\ge\ord f+\ord g.
\]
Dacă $R$ este domeniu de integritate, atunci $\ord(fg)=\ord f+\ord g$.
\end{propozitie}
\begin{proof}
Prima inegalitate: termenul de exponent $<\min\{\ord f,\ord g\}$ e nul în ambele, deci în sumă. Pentru produs: coeficientul lui $X^k$ în $fg$ este $\sum_{i+j=k}a_ib_j$, care e nul pentru $k<\ord f+\ord g$ (orice pereche are $i<\ord f$ sau $j<\ord g$). Când $R$ e domeniu, coeficientul lui $X^{\ord f+\ord g}$ este exact $a_{\ord f}\,b_{\ord g}\neq0$, deci egalitate. (Rezultat simetric celui pentru grad de mai sus.)
\end{proof}

\begin{observatie}[polinom vs.\ funcție polinomială]\label{obs:functie}
Polinomul $f$ și funcția $\tilde f:A\to A$, $\tilde f(x)=f(x)$, sunt obiecte diferite. Peste un corp \textbf{finit} ele nu se determină reciproc: în $\ZZ_p[X]$, polinomul $X^p-X$ este nenul, dar funcția sa este identic nulă (Fermat). Peste un inel integru \textbf{infinit}, egalitatea funcțiilor implică egalitatea polinoamelor (Corolarul~\ref{cor:identitate}).
\end{observatie}

\begin{teorema}[împărțirea cu rest]\label{teo:impartire}
Fie $K$ un corp comutativ și $f,g\in K[X]$, $g\neq0$. Există și sunt \textbf{unice} $q,r\in K[X]$ cu
\[
f=gq+r,\qquad \deg r<\deg g.
\]
\end{teorema}

\begin{proof}
\emph{Unicitatea}: din $gq_1+r_1=gq_2+r_2$ rezultă $g(q_1-q_2)=r_2-r_1$; dacă $q_1\neq q_2$, membrul stâng are grad $\ge\deg g$, cel drept $<\deg g$ --- contradicție. Deci $q_1=q_2$, apoi $r_1=r_2$.
\emph{Existența}, prin inducție după $\deg f$: dacă $\deg f<\deg g$, luăm $q=0$, $r=f$. Altfel, cu $a,b$ coeficienții dominanți ai lui $f,g$ și $n=\deg f\ge m=\deg g$, polinomul $f_1=f-\frac ab X^{\,n-m} g$ are grad $<n$; aplicăm ipoteza de inducție lui $f_1$ și recompunem.
\end{proof}

\begin{observatie}
Aceeași demonstrație arată că peste \textbf{orice} inel comutativ $A$ împărțirea cu rest funcționează dacă împărțitorul $g$ este \emph{unitar} (sau are coeficient dominant inversabil) --- util pentru $\ZZ[X]$.
\end{observatie}

\begin{definitie}
$g\mid f$ dacă $f=gq$ pentru un $q\in K[X]$. Polinoamele $f$ și $cf$ ($c\in K^*$) se numesc \emph{asociate}. Cel mai mare divizor comun $(f,g)$ este divizorul comun de grad maxim, ales unitar (pentru unicitate).
\end{definitie}

\begin{teorema}[algoritmul lui Euclid; Bézout]\label{teo:bezout}
Pentru $f,g\in K[X]$, nu ambele nule, c.m.m.d.c.\ $d=(f,g)$ există, se obține ca ultimul rest nenul din algoritmul lui Euclid (împărțiri cu rest succesive) și există $u,v\in K[X]$ cu
\[
d=uf+vg.
\]
\end{teorema}

\begin{proof}[Schiță]
Ca în $\ZZ$: dacă $f=gq+r$, atunci divizorii comuni ai perechii $(f,g)$ coincid cu cei ai perechii $(g,r)$; gradele resturilor scad strict, deci algoritmul se oprește. Scriind fiecare rest ca o combinație $u_if+v_ig$ (recurență), ultimul rest nenul are forma cerută.
\end{proof}

\begin{definitie}\label{def:ireductibil}
$f\in K[X]$ cu $\deg f\ge1$ este \emph{ireductibil peste $K$} dacă nu se poate scrie $f=gh$ cu $g,h\in K[X]$, $\deg g,\deg h\ge1$; altfel este \emph{reductibil}.
\end{definitie}

\begin{lema}[Euclid pentru polinoame]\label{lema:euclidpol}
Dacă $p\in K[X]$ este ireductibil și $p\mid fg$, atunci $p\mid f$ sau $p\mid g$.
\end{lema}

\begin{proof}
Dacă $p\nmid f$, atunci $(p,f)=1$ (divizorii lui $p$ sunt constante sau asociați ai lui $p$), deci $1=up+vf$ (Bézout). Atunci $g=upg+v(fg)$ este divizibil cu $p$.
\end{proof}

\begin{teorema}[descompunerea în factori ireductibili]\label{teo:descompunere}
Orice $f\in K[X]$ cu $\deg f\ge1$ se scrie ca produs de polinoame ireductibile peste $K$; scrierea este unică până la ordinea factorilor și înmulțirea cu constante nenule.
\end{teorema}

\begin{proof}[Schiță]
Existența: inducție tare după grad (dacă $f$ nu e ireductibil, se sparge în factori de grade mai mici). Unicitatea: din două descompuneri, Lema~\ref{lema:euclidpol} arată că fiecare factor ireductibil dintr-o parte divide (deci este asociat cu) un factor din cealaltă; simplificăm și repetăm --- exact ca la teorema fundamentală a aritmeticii.
\end{proof}

\section{Rădăcini. Bézout. Derivata formală}

\begin{teorema}[teorema restului și a factorului]\label{teo:rest}
Fie $f\in A[X]$, $A$ comutativ, și $a\in A$. Restul împărțirii lui $f$ la $X-a$ este $f(a)$:
\[
f=(X-a)\,q+f(a).
\]
În particular \emph{(Bézout)}: $a$ este rădăcină a lui $f$ $\iff(X-a)\mid f$.
\end{teorema}

\begin{proof}
Împărțim la polinomul unitar $X-a$: $f=(X-a)q+r$ cu $\deg r\le0$, deci $r$ constant; punând $X=a$ obținem $r=f(a)$.
\end{proof}

\begin{observatie}
\emph{Schema lui Horner} calculează simultan câtul $q$ și valoarea $f(a)$ prin recurența $b_{n-1}=a_n$, $b_{k-1}=a_k+ab_k$; este metoda rapidă de testat rădăcini la concurs.
\end{observatie}

\begin{definitie}
$a$ este rădăcină cu \emph{ordinul de multiplicitate} $m\ge1$ pentru $f$ dacă $(X-a)^m\mid f$ și $(X-a)^{m+1}\nmid f$; echivalent, $f=(X-a)^m g$ cu $g(a)\neq0$.
\end{definitie}

\begin{teorema}[numărul rădăcinilor]\label{teo:numarradacini}
Fie $A$ un inel \textbf{integru} și $f\in A[X]$, $f\neq0$, $\deg f=n$. Atunci $f$ are cel mult $n$ rădăcini în $A$, \emph{numărate cu multiplicități}.
\end{teorema}

\begin{proof}
Inducție după $n$. Dacă $f$ nu are rădăcini, gata. Altfel, fie $a$ o rădăcină de multiplicitate $m$: $f=(X-a)^mg$, $g(a)\neq0$, $\deg g=n-m$. Orice altă rădăcină $b\neq a$ dă $0=(b-a)^m g(b)$; cum $A$ este integru și $b-a\neq0$, rezultă $g(b)=0$, cu aceeași multiplicitate în $g$ ca în $f$ (același argument de integritate). Prin ipoteza de inducție, $g$ are cel mult $n-m$ rădăcini cu multiplicități, deci $f$ are cel mult $m+(n-m)=n$.
\end{proof}

\begin{exemplu}[integritatea este esențială]
În $\ZZ_8[X]$, polinomul $X^2-\cl1$ de gradul $2$ are \textbf{patru} rădăcini: $\cl1,\cl3,\cl5,\cl7$.
\end{exemplu}

\begin{corolar}[principiul identității]\label{cor:identitate}
Fie $A$ integru. \emph{(i)} Dacă $\deg f,\deg g\le n$ și $f,g$ iau valori egale în $n+1$ puncte distincte, atunci $f=g$. \emph{(ii)} Dacă $A$ este infinit și $\tilde f=\tilde g$ ca funcții, atunci $f=g$ ca polinoame.
\end{corolar}

\begin{proof}
$f-g$ are grad $\le n$ și cel puțin $n+1$ rădăcini (respectiv o infinitate), deci este polinomul nul.
\end{proof}

\begin{definitie}
\emph{Derivata formală} a lui $f=\sum_{k=0}^{n}a_kX^k$ este $f'=\sum_{k=1}^{n}ka_kX^{k-1}$ --- definiție pur algebrică, fără limite.
\end{definitie}

\begin{propozitie}\label{prop:derivata}
$(f+g)'=f'+g'$, $(cf)'=cf'$, $(fg)'=f'g+fg'$ și $\bigl((X-a)^m\bigr)'=m(X-a)^{m-1}$.
\end{propozitie}

\begin{proof}[Schiță]
Primele două sunt imediate. Regula produsului se verifică pe monoame $f=X^i$, $g=X^j$ (ambii membri dau $(i+j)X^{i+j-1}$) și se extinde prin biliniaritate; ultima rezultă prin inducție din regula produsului.
\end{proof}

\begin{teorema}[criteriul rădăcinii multiple]\label{teo:multipla}
Fie $K$ corp comutativ, $f\in K[X]$, $a\in K$. Atunci $a$ este rădăcină multiplă (de ordin $\ge2$) a lui $f$ $\iff f(a)=f'(a)=0$.
\end{teorema}

\begin{proof}
Dacă $f(a)=0$, scriem $f=(X-a)g$; atunci $f'=g+(X-a)g'$, deci $f'(a)=g(a)$. Rădăcina $a$ este multiplă $\iff g(a)=0\iff f'(a)=0$.
\end{proof}

\begin{teorema}[multiplicitatea prin derivate succesive]\label{teo:derivatesucc}
Fie $f\in K[X]$ cu $K\in\{\QQ,\RR,\CC\}$. Atunci $a$ are ordinul de multiplicitate \textbf{exact} $m$ $\iff$
\[
f(a)=f'(a)=\dots=f^{(m-1)}(a)=0\quad\text{și}\quad f^{(m)}(a)\neq0.
\]
\end{teorema}

\begin{proof}
Cheia: dacă $f=(X-a)^mg$ cu $g(a)\neq0$ și $m\ge1$, atunci
$f'=(X-a)^{m-1}\bigl[mg+(X-a)g'\bigr]$, iar paranteza ia în $a$ valoarea $mg(a)\neq0$ (aici folosim $\car K=0$). Deci multiplicitatea lui $a$ în $f'$ este exact $m-1$; o inducție încheie demonstrația în ambele sensuri.
\end{proof}

\begin{observatie}
\textbf{Atenție în caracteristică $p$}: în $\ZZ_p[X]$, $f=X^p$ are $f'=pX^{p-1}=0$, deci criteriul cu derivate succesive eșuează; Teorema~\ref{teo:multipla} (cu $f$ și $f'$) rămâne însă valabilă peste orice corp.
\end{observatie}

\section{Relațiile lui Viète și sumele lui Newton}

\begin{teorema}[Viète]\label{teo:viete}
Fie $f=a_nX^n+\dots+a_0\in\CC[X]$, $a_n\neq0$, cu rădăcinile $x_1,\dots,x_n$ (cu multiplicități). Notând sumele simetrice fundamentale $e_k=\sum_{i_1<\dots<i_k}x_{i_1}\cdots x_{i_k}$, avem
\[
e_k=(-1)^k\,\frac{a_{n-k}}{a_n},\qquad k=1,\dots,n.
\]
\end{teorema}

\begin{proof}
Din descompunerea $f=a_n(X-x_1)\cdots(X-x_n)$ (Teorema~\ref{teo:tfa}), dezvoltăm produsul: coeficientul lui $X^{n-k}$ este $a_n(-1)^ke_k$; identificăm cu $a_{n-k}$.
\end{proof}

\begin{exemplu}
$n=2$: $x_1+x_2=-\frac{a_1}{a_2}$, $x_1x_2=\frac{a_0}{a_2}$.\quad
$n=3$: pentru $aX^3+bX^2+cX+d$,
\[
x_1+x_2+x_3=-\frac ba,\qquad x_1x_2+x_2x_3+x_3x_1=\frac ca,\qquad x_1x_2x_3=-\frac da.
\]
\end{exemplu}

\begin{observatie}[reciproca]
Numerele $x_1,\dots,x_n$ cu sumele simetrice $e_1,\dots,e_n$ date sunt exact rădăcinile polinomului
$X^n-e_1X^{n-1}+e_2X^{n-2}-\dots+(-1)^ne_n$. Aceasta transformă sisteme simetrice în ecuații polinomiale --- tipar frecvent la concurs.
\end{observatie}

\begin{teorema}[sumele de puteri]\label{teo:newton}
Fie $f=X^n+c_{n-1}X^{n-1}+\dots+c_0$ unitar, cu rădăcinile $x_1,\dots,x_n$, și $s_k=x_1^k+\dots+x_n^k$ (cu $s_0=n$). Atunci:
\emph{(i)} pentru orice $k\ge n$,
\[
s_k+c_{n-1}s_{k-1}+c_{n-2}s_{k-2}+\dots+c_0\,s_{k-n}=0;
\]
\emph{(ii)} \emph{(identitățile lui Newton, admise)} pentru $1\le k\le n$,
\[
s_k+c_{n-1}s_{k-1}+\dots+c_{n-k+1}s_1+k\,c_{n-k}=0.
\]
\end{teorema}

\begin{proof}[Demonstrație pentru (i)]
Fiecare rădăcină satisface $x_i^n+c_{n-1}x_i^{n-1}+\dots+c_0=0$. Înmulțim cu $x_i^{\,k-n}$ și sumăm după $i$. (Identitățile (ii) se verifică direct pentru $n$ mic; de exemplu $k=1$: $s_1+c_{n-1}=0$, adică Viète.)
\end{proof}

\begin{aplicatie}[identitatea celor trei cuburi]\label{apl:treicuburi}
Pentru $n=k=3$, cu $c_2=-e_1$, $c_1=e_2$, $c_0=-e_3$, relația (ii) dă $s_3=e_1s_2-e_2s_1+3e_3$. Cum $s_2=e_1^2-2e_2$, rezultă
\[
x^3+y^3+z^3-3xyz=(x+y+z)\bigl(x^2+y^2+z^2-xy-yz-zx\bigr).
\]
\end{aplicatie}

\begin{definitie}[inelul de polinoame în mai multe variabile]
Inelul polinoamelor în $n$ nedeterminate cu coeficienți în $R$ se definește inductiv prin
\[
R[X_1,\dots,X_{n+1}]:=\big(R[X_1,\dots,X_n]\big)[X_{n+1}].
\]
Elementele sale se scriu în mod unic ca sume finite $\sum a_{i_1\dots i_n}X_1^{i_1}\cdots X_n^{i_n}$, cu $a_{i_1\dots i_n}\in R$.
\end{definitie}

\begin{definitie}[polinom simetric]
$f\in R[X_1,\dots,X_n]$ se numește \emph{simetric} dacă $f(X_{\sigma(1)},\dots,X_{\sigma(n)})=f(X_1,\dots,X_n)$ pentru orice permutare $\sigma\in S_n$. Polinoamele simetrice fundamentale sunt $e_1=\sum X_i$, $e_2=\sum_{i<j}X_iX_j$, \dots, $e_n=X_1\cdots X_n$ (notația $s_k$ este rezervată sumelor de puteri, Teorema~\ref{teo:newton}).
\end{definitie}

\begin{propozitie}[funcțiile simetrice de rădăcini rămân în inelul de bază]\label{prop:simroots}
Fie $R\subseteq S$ inele integre, $f\in R[X]$ \emph{monic} de grad $n$, cu toate rădăcinile $x_1,\dots,x_n$ în $S$. Atunci pentru orice polinom simetric $g\in R[X_1,\dots,X_n]$ avem $g(x_1,\dots,x_n)\in R$.
\end{propozitie}
\begin{proof}
Prin Teorema fundamentală a polinoamelor simetrice (de mai jos), $g=G(e_1,\dots,e_n)$ cu $G\in R[X_1,\dots,X_n]$. Cum $f=(X-x_1)\cdots(X-x_n)$ (monic), relațiile lui Viète (Teorema~\ref{teo:viete}) dau $e_k(x_1,\dots,x_n)=(-1)^k a_{n-k}\in R$ (coeficienții lui $f$). Deci $g(x_1,\dots,x_n)=G\big(e_1(x),\dots,e_n(x)\big)$ este o expresie polinomială cu coeficienți în $R$ evaluată în elemente din $R$, deci aparține lui $R$.
\end{proof}

\begin{definitie}[corpul funcțiilor raționale]
Pentru un corp comutativ $K$, corpul de fracții al domeniului $K[X_1,\dots,X_n]$ se notează $K(X_1,\dots,X_n)$ și se numește \emph{corpul funcțiilor raționale} în $X_1,\dots,X_n$ cu coeficienți în $K$.
\end{definitie}

\begin{teorema}[teorema fundamentală a polinoamelor simetrice]\label{teo:polsim}
Orice polinom simetric în $x_1,\dots,x_n$ se exprimă în mod \textbf{unic} ca polinom în polinoamele simetrice fundamentale $e_1,\dots,e_n$.
\end{teorema}

La concurs se citează ca atare; calculul efectiv al unei expresii simetrice date se face de obicei cu relațiile lui Newton de mai sus.

\section{Rădăcini pentru polinoame din $\ZZ[X]$, $\QQ[X]$, $\RR[X]$, $\CC[X]$}

\begin{teorema}[rădăcinile raționale]\label{teo:rationale}
Fie $f=a_nX^n+\dots+a_0\in\ZZ[X]$, $a_n\neq 0$, și $\frac pq\in\QQ$ o rădăcină, cu fracția ireductibilă. Atunci $p\mid a_0$ și $q\mid a_n$. În particular, dacă $f$ este unitar, orice rădăcină rațională este întreagă și divide $a_0$.
\end{teorema}

\begin{proof}
Din $a_np^n+a_{n-1}p^{n-1}q+\dots+a_0q^n=0$: toți termenii în afara ultimului sunt divizibili cu $p$, deci $p\mid a_0q^n$, iar $(p,q)=1$ dă $p\mid a_0$. Simetric pentru $q$.
\end{proof}

\begin{propozitie}\label{prop:aminusb}
Dacă $f\in\ZZ[X]$ și $a,b\in\ZZ$, atunci $(a-b)\mid f(a)-f(b)$.
\end{propozitie}

\begin{proof}
$f(a)-f(b)=\sum_k a_k(a^k-b^k)$ și $(a-b)\mid a^k-b^k=(a-b)(a^{k-1}+\dots+b^{k-1})$.
\end{proof}

\begin{observatie}[testul de paritate]
Dacă $f\in\ZZ[X]$ și $f(0),f(1)$ sunt ambele impare, atunci $f$ nu are rădăcini întregi: pentru $k$ par, $f(k)\equiv f(0)\equiv1\pmod2$; pentru $k$ impar, $f(k)\equiv f(1)\equiv1\pmod 2$.
\end{observatie}

\begin{aplicatie}[ciclu de lungime 3 --- problemă clasică]\label{apl:ciclu}
Nu există $f\in\ZZ[X]$ și întregi \emph{distincți} $a,b,c$ cu $f(a)=b$, $f(b)=c$, $f(c)=a$.

\emph{Soluție.} Din Propoziția~\ref{prop:aminusb}: $(a-b)\mid f(a)-f(b)=b-c$, apoi $(b-c)\mid c-a$ și $(c-a)\mid a-b$. Deci $|a-b|\le|b-c|\le|c-a|\le|a-b|$, adică toate cele trei module sunt egale cu un $d>0$. Dar $(a-b)+(b-c)+(c-a)=0$, iar o sumă de trei termeni egali cu $\pm d$ nu poate fi $0$ (ar trebui $\pm d\pm d\pm d=0$, imposibil pentru $d\neq0$). Contradicție.
\end{aplicatie}

\begin{teorema}[rădăcini complexe conjugate]\label{teo:conjugata}
Fie $f\in\RR[X]$ și $z\in\CC\setminus\RR$ o rădăcină. Atunci $\bar z$ este rădăcină, \textbf{cu același ordin de multiplicitate}.
\end{teorema}

\begin{proof}
Conjugarea este compatibilă cu operațiile ($\overline{u+v}=\bar u+\bar v$, $\overline{uv}=\bar u\,\bar v$) și fixează coeficienții reali, deci $f(\bar z)=\overline{f(z)}=0$. Pentru multiplicitate: $q=(X-z)(X-\bar z)=X^2-2\mathrm{Re}(z)X+|z|^2\in\RR[X]$; împărțind în $\RR[X]$, $f=q\,h+r$ cu $\deg r\le1$, iar $r(z)=0$ cu $z\notin\RR$ forțează $r=0$ (un polinom real de grad $\le1$ cu rădăcină nereală este nul). Deci $f=qh$ cu $h\in\RR[X]$, iar multiplicitățile lui $z$ și $\bar z$ în $f$ sunt cu $1$ mai mari decât în $h$; inducția încheie.
\end{proof}

\begin{corolar}
Orice polinom din $\RR[X]$ de grad impar are cel puțin o rădăcină reală. (Rădăcinile nereale se grupează în perechi conjugate cu multiplicități egale, deci sunt în număr par.)
\end{corolar}

\begin{teorema}[conjugata pătratică]\label{teo:conjpatratica}
Fie $f\in\QQ[X]$, $a,b,d\in\QQ$ cu $\sqrt d\notin\QQ$. Dacă $a+b\sqrt d$ este rădăcină a lui $f$, atunci și $a-b\sqrt d$ este rădăcină (cu același ordin de multiplicitate).
\end{teorema}

\begin{proof}
Prin inducție, $(a+b\sqrt d)^k=A_k+B_k\sqrt d$ și $(a-b\sqrt d)^k=A_k-B_k\sqrt d$, cu $A_k,B_k\in\QQ$ (pasul de inducție: înmulțirea cu $a\pm b\sqrt d$ păstrează simetria semnului). Sumând cu coeficienții lui $f$: $f(a+b\sqrt d)=P+Q\sqrt d$ și $f(a-b\sqrt d)=P-Q\sqrt d$, cu $P,Q\in\QQ$. Din $P+Q\sqrt d=0$: dacă $Q\neq0$, atunci $\sqrt d=-P/Q\in\QQ$, contradicție; deci $Q=0$, apoi $P=0$, adică $f(a-b\sqrt d)=0$. Multiplicitatea: se împarte la $\bigl(X-(a+b\sqrt d)\bigr)\bigl(X-(a-b\sqrt d)\bigr)=X^2-2aX+(a^2-db^2)\in\QQ[X]$ și se raționează ca la Teorema~\ref{teo:conjugata}.
\end{proof}

\begin{teorema}[teorema fundamentală a algebrei; d'Alembert--Gauss]\label{teo:tfa}
\emph{(Admisă.)} Orice polinom din $\CC[X]$ de grad $\ge1$ are cel puțin o rădăcină în $\CC$. Prin urmare, orice $f\in\CC[X]$ de grad $n\ge 1$ se descompune
\[
f=a_n(X-x_1)^{m_1}\cdots(X-x_r)^{m_r},\qquad m_1+\dots+m_r=n.
\]
\end{teorema}

\begin{corolar}[{descompunerea în $\RR[X]$}]\label{cor:descR}
Orice $f\in\RR[X]$ de grad $\ge1$ se descompune în $\RR[X]$ ca produs de factori de gradul întâi și factori de gradul al doilea cu discriminant negativ.
\end{corolar}

\begin{proof}
Descompunem în $\CC[X]$ și grupăm rădăcinile nereale în perechi conjugate (Teorema~\ref{teo:conjugata}): $(X-z)(X-\bar z)\in\RR[X]$ cu $\Delta<0$.
\end{proof}

\begin{teorema}[corp de descompunere]\label{teo:corpdescompunere}
Fie $K$ un corp comutativ și $f\in K[X]$ un polinom neconstant de grad $n$. Atunci există un corp $L\supseteq K$ în care $f$ are $n$ rădăcini (adică $f$ se descompune complet în factori de gradul întâi peste $L$).
\end{teorema}
\begin{proof}
Procedăm prin inducție după $n=\deg f$, arătând că putem \emph{adăuga o rădăcină} și apoi itera. Fie $p$ un factor ireductibil al lui $f$ în $K[X]$. Inelul cât $L_1=K[X]/(p)$ este \emph{corp}, deoarece $(p)$ este ideal maximal ($p$ ireductibil), iar $K$ se scufundă în $L_1$ prin $a\mapsto a+(p)$. Clasa $\alpha=X+(p)$ satisface $p(\alpha)=0$, deci și $f(\alpha)=0$: peste $L_1$, $f$ are rădăcina $\alpha$, deci $f=(X-\alpha)f_1$ cu $f_1\in L_1[X]$, $\deg f_1=n-1$. Prin ipoteza de inducție aplicată lui $f_1$ peste $L_1$, există $L\supseteq L_1\supseteq K$ în care $f_1$ (deci și $f$) are toate rădăcinile. (Este construcția lui Kronecker; pentru $K=\ZZ_p$ ea este exact mecanismul folosit mai jos la corpuri finite, $\ZZ_p[X]/(p)$.)
\end{proof}

\section{Ireductibilitate: criterii practice}

\begin{propozitie}[gradele mici]\label{prop:gradmic}
Peste un corp $K$: \emph{(i)} orice polinom de gradul $1$ este ireductibil; \emph{(ii)} un polinom de gradul $2$ sau $3$ este ireductibil $\iff$ nu are rădăcini în $K$.
\end{propozitie}

\begin{proof}
(ii) O descompunere netrivială a unui polinom de grad $2$ sau $3$ conține obligatoriu un factor de gradul $1$, $aX+b$, care dă rădăcina $-b/a\in K$; reciproc, o rădăcină $\alpha$ dă factorul $X-\alpha$ (Bézout).
\end{proof}

\begin{observatie}[capcană la gradul 4]
Pentru grad $\ge4$, lipsa rădăcinilor \textbf{nu} implică ireductibilitatea:
\[
X^4+X^2+1=(X^2+X+1)(X^2-X+1)
\]
nu are nicio rădăcină reală (este $>0$), dar este reductibil peste $\QQ$. Verificarea: $(X^2+1)^2-X^2$.
\end{observatie}

\begin{corolar}[din TFA]
Peste $\CC$, ireductibile sunt exact polinoamele de gradul $1$; peste $\RR$, cele de gradul $1$ și cele de gradul $2$ cu $\Delta<0$ (Corolarul~\ref{cor:descR}). Chestiunea ireductibilității este deci interesantă peste $\QQ$ (și peste $\ZZ_p$).
\end{corolar}

\begin{definitie}
Pentru $f\in\ZZ[X]$ nenul, \emph{conținutul} $c(f)$ este c.m.m.d.c.\ (pozitiv) al coeficienților; $f$ este \emph{primitiv} dacă $c(f)=1$. Orice $f\in\ZZ[X]$ nenul se scrie $f=c(f)\,f_0$ cu $f_0$ primitiv.
\end{definitie}

\begin{lema}[Gauss]\label{lema:gauss}
Produsul a două polinoame primitive din $\ZZ[X]$ este primitiv.
\end{lema}

\begin{proof}
Fie $g,h$ primitive și presupunem că un prim $p$ divide toți coeficienții lui $gh$. Reducem coeficienții modulo $p$: în $\ZZ_p[X]$ avem $\bar g\cdot\bar h=\overline{gh}=0$. Dar $\ZZ_p[X]$ este inel integru (Propoziția~\ref{prop:grade}(iii), $\ZZ_p$ corp), deci $\bar g=0$ sau $\bar h=0$, adică $p$ divide toți coeficienții lui $g$ sau ai lui $h$ --- contradicție cu primitivitatea.
\end{proof}

\begin{teorema}[teorema lui Gauss]\label{teo:gauss}
Fie $f\in\ZZ[X]$, $\deg f\ge1$. Dacă $f=gh$ cu $g,h\in\QQ[X]$ de grade $\ge1$, atunci există $g_1,h_1\in\ZZ[X]$ cu $f=g_1h_1$, $\deg g_1=\deg g$, $\deg h_1=\deg h$. Echivalent: \emph{ireductibil peste $\ZZ$} (în sensul gradelor) $\Rightarrow$ \emph{ireductibil peste $\QQ$}.
\end{teorema}

\begin{proof}
Alegem $m,n\in\NN^*$ care elimină numitorii: $G=mg\in\ZZ[X]$, $H=nh\in\ZZ[X]$, deci $mn\,f=GH$. Scriem $G=c(G)G_0$, $H=c(H)H_0$ cu $G_0,H_0$ primitive; atunci $mn\,f=c(G)c(H)\,G_0H_0$, iar $G_0H_0$ este primitiv (Lema~\ref{lema:gauss}). Luând conținutul ambilor membri: $mn\,c(f)=c(G)c(H)$. Prin urmare
\[
f=\frac{c(G)c(H)}{mn}\,G_0H_0=c(f)\,G_0H_0=\bigl(c(f)G_0\bigr)\cdot H_0,
\]
o descompunere în $\ZZ[X]$ cu aceleași grade.
\end{proof}

\begin{teorema}[criteriul lui Eisenstein]\label{teo:eisenstein}
Fie $f=a_nX^n+\dots+a_0\in\ZZ[X]$ și $p$ un prim cu:
\[
p\nmid a_n,\qquad p\mid a_i\ (0\le i\le n-1),\qquad p^2\nmid a_0.
\]
Atunci $f$ este ireductibil peste $\QQ$.
\end{teorema}

\begin{proof}
Prin Teorema~\ref{teo:gauss} este suficient să arătăm că $f$ nu se scrie $f=gh$ cu $g,h\in\ZZ[X]$, $\deg g=m\ge1$, $\deg h=k\ge1$. Fie $g=b_mX^m+\dots+b_0$, $h=c_kX^k+\dots+c_0$. Din $a_0=b_0c_0$, $p\mid a_0$ și $p^2\nmid a_0$, exact unul dintre $b_0,c_0$ este divizibil cu $p$; zicem $p\mid b_0$, $p\nmid c_0$. Din $a_n=b_mc_k$ și $p\nmid a_n$ rezultă $p\nmid b_m$; fie deci $r=\min\{i\mid p\nmid b_i\}$, cu $1\le r\le m<n$. Atunci
\[
a_r=b_rc_0+\underbrace{b_{r-1}c_1+\dots+b_0c_r}_{\text{divizibil cu }p},
\]
iar $p\mid a_r$ (deoarece $r<n$), deci $p\mid b_rc_0$ --- contradicție, căci $p\nmid b_r$ și $p\nmid c_0$.
\end{proof}

\begin{exemplu}
$X^n-2$ și, mai general, $X^n-p$ sunt ireductibile peste $\QQ$ pentru orice $n\ge1$ (Eisenstein cu primul $p$). Consecință: există polinoame ireductibile peste $\QQ$ de orice grad.
\end{exemplu}

\begin{aplicatie}[polinomul ciclotomic $\Phi_p$]\label{apl:ciclotomic}
Pentru $p$ prim, $f=X^{p-1}+X^{p-2}+\dots+X+1=\dfrac{X^p-1}{X-1}$ este ireductibil peste $\QQ$.

\emph{Soluție.} Substituția $X\mapsto X+1$ păstrează (i)reductibilitatea (este inversabilă, prin $X\mapsto X-1$). Calculăm
\[
f(X+1)=\frac{(X+1)^p-1}{X}=X^{p-1}+\binom p1X^{p-2}+\dots+\binom{p}{p-2}X+\binom{p}{p-1}.
\]
Coeficientul dominant este $1$, toți ceilalți coeficienți $\binom pk$ ($1\le k\le p-1$) sunt divizibili cu $p$ (demonstrat la Teorema~\ref{teo:frobenius}), iar termenul liber este $\binom p{p-1}=p$, nedivizibil cu $p^2$. Eisenstein încheie.
\end{aplicatie}

\begin{teorema}[reducerea modulo $p$]\label{teo:redmodp}
Fie $f\in\ZZ[X]$ și $p$ un prim cu $p\nmid a_n$ (coeficientul dominant). Dacă polinomul redus $\bar f\in\ZZ_p[X]$ este ireductibil peste $\ZZ_p$, atunci $f$ este ireductibil peste $\QQ$.
\end{teorema}

\begin{proof}
Presupunem $f$ reductibil peste $\QQ$; prin Teorema~\ref{teo:gauss}, $f=gh$ în $\ZZ[X]$ cu $\deg g,\deg h\ge1$. Coeficienții dominanți satisfac $a_n=(\text{dom }g)(\text{dom }h)$, deci $p$ nu divide niciunul; prin urmare $\deg\bar g=\deg g$ și $\deg\bar h=\deg h$, iar $\bar f=\bar g\,\bar h$ este o descompunere netrivială în $\ZZ_p[X]$ --- contradicție.
\end{proof}

\begin{observatie}
(a) Reciproca este falsă: un exemplu clasic este $X^4+1$, ireductibil peste $\QQ$, dar reductibil modulo \emph{orice} prim $p$. Deci eșecul criteriului pentru un $p$ nu spune nimic. (b) Ireductibilitatea \textbf{depinde de corp}: $X^2-2$ este ireductibil peste $\QQ$, reductibil peste $\RR$; $X^2+1$ este ireductibil peste $\RR$ și peste $\ZZ_3$, dar reductibil peste $\CC$ și peste $\ZZ_5$ ($X^2+\cl1=(X-\cl2)(X+\cl2)$ în $\ZZ_5[X]$).
\end{observatie}

\section{Polinoame peste $\ZZ_p$. Rădăcini primitive}

\begin{teorema}\label{teo:XpX}
În $\ZZ_p[X]$ ($p$ prim) are loc descompunerea
\[
X^p-X=\prod_{\cl a\in\ZZ_p}(X-\cl a),\qquad\text{deci}\qquad X^{p-1}-\cl1=\prod_{\cl a\in\ZZ_p^*}(X-\cl a).
\]
\end{teorema}

\begin{proof}
Prin Fermat (Teorema~\ref{teo:euler}), fiecare dintre cele $p$ elemente ale lui $\ZZ_p$ este rădăcină a lui $X^p-X$; prin Bézout, produsul $\prod(X-\cl a)$ divide $X^p-X$ (factorii sunt primi între ei doi câte doi). Ambele polinoame sunt unitare de grad $p$, deci sunt egale. A doua egalitate: simplificăm cu $X$ (adică factorul $X-\cl 0$).
\end{proof}

\begin{corolar}[Wilson, a doua demonstrație]
Identificând termenii liberi în $X^{p-1}-\cl 1=\prod_{\cl a\neq\cl 0}(X-\cl a)$:
$-\cl1=(-1)^{p-1}\cdot\cl1\cdot\cl2\cdots\cl{p-1}$, deci pentru $p$ impar $\overline{(p-1)!}=-\cl1$ (iar $p=2$ direct).
\end{corolar}

\begin{lema}[identitatea lui Gauss pentru $\varphi$]\label{lema:sumaphi}
$\displaystyle\sum_{d\mid n}\varphi(d)=n$ pentru orice $n\ge1$.
\end{lema}

\begin{proof}
Împărțim $\{1,2,\dots,n\}$ după c.m.m.d.c.\ cu $n$: pentru fiecare divizor $e\mid n$, numerele $k$ cu $(k,n)=e$ sunt $k=ej$ cu $1\le j\le n/e$ și $(j,n/e)=1$, în număr de $\varphi(n/e)$. Sumând după $e$ și schimbând $d=n/e$: $n=\sum_{e\mid n}\varphi(n/e)=\sum_{d\mid n}\varphi(d)$.
\end{proof}

\begin{teorema}[ciclicitate; rădăcini primitive]\label{teo:ciclicitate}
Orice subgrup finit $G$ al grupului multiplicativ $(K^*,\cdot)$ al unui corp comutativ $K$ este ciclic. În particular, $(\ZZ_p^*,\cdot)$ este ciclic: există $g$ (numit \emph{rădăcină primitivă modulo $p$}) cu $\ZZ_p^*=\{\cl 1,\cl g,\dots,\cl g^{\,p-2}\}$.
\end{teorema}

\begin{proof}
Fie $|G|=n$ și, pentru $d\mid n$, fie $\psi(d)$ numărul elementelor de ordin $d$ din $G$ (ordinele divid $n$, prin Lagrange). Dacă $\psi(d)>0$, alegem $x$ de ordin $d$. Toate cele $d$ elemente din $\langle x\rangle$ satisfac $y^d=1$, iar în corpul $K$ ecuația polinomială $y^d=1$ are \textbf{cel mult} $d$ soluții (Teorema~\ref{teo:numarradacini}); deci mulțimea soluțiilor este exact $\langle x\rangle$. Orice element de ordin $d$ este o astfel de soluție, deci aparține lui $\langle x\rangle$ și este generator al acestuia: sunt exact $\varphi(d)$ asemenea generatori (Propoziția~\ref{prop:generatori}). Așadar $\psi(d)\in\{0,\varphi(d)\}$ pentru fiecare $d\mid n$. Dar
\[
n=\sum_{d\mid n}\psi(d)\le\sum_{d\mid n}\varphi(d)=n\qquad(\text{Lema~\ref{lema:sumaphi}}),
\]
deci peste tot avem egalitate: $\psi(d)=\varphi(d)$ pentru orice $d\mid n$. În particular $\psi(n)=\varphi(n)\ge1$: există un element de ordin $n$, care generează $G$.
\end{proof}

\begin{observatie}
Demonstrația dă mai mult: în $\ZZ_p^*$ există exact $\varphi(d)$ elemente de ordin $d$, pentru fiecare $d\mid p-1$. Împreună cu Propoziția~\ref{prop:subgrupC}, obținem și structura subgrupurilor finite ale lui $\CC^*$.
\end{observatie}

\begin{observatie}[a doua demonstrație a ciclicității]
Corolarul~\ref{cor:lcm} dă o demonstrație alternativă: fie $m$ ordinul maxim al unui element din $G$, atins de $a$. Pentru orice $x\in G$, grupul abelian $G$ conține un element de ordin $[\,\ord x,\,m\,]$; prin maximalitate, $[\,\ord x,\,m\,]\le m$, deci $\ord x\mid m$. Atunci toate cele $|G|$ elemente sunt rădăcini ale polinomului $X^m-1$, care are cel mult $m$ rădăcini în $K$ (Teorema~\ref{teo:numarradacini}); așadar $|G|\le m=\ord a\le|G|$, deci $a$ generează $G$.
\end{observatie}


\begin{propozitie}[criteriul lui Euler]\label{prop:eulercrit}
Fie $p$ prim impar și $a$ cu $p\nmid a$. Atunci $a^{\frac{p-1}2}\equiv\pm1\pmod p$, iar
\[
a^{\frac{p-1}{2}}\equiv1\pmod p\iff a\text{ este rest pătratic mod }p\ (\exists\,x:\ x^2\equiv a).
\]
\end{propozitie}

\begin{proof}
Fie $g$ o rădăcină primitivă și $\cl a=\cl g^{\,t}$. Elementul $h=\cl g^{\,\frac{p-1}2}$ satisface $h^2=\cl 1$ și $h\neq\cl1$ (ordinul lui $g$ este $p-1$), deci $h=-\cl1$ (singurele soluții ale lui $x^2=\cl1$ în corp). Atunci $\cl a^{\frac{p-1}2}=(-\cl 1)^t$. Dacă $t$ este par, $a=g^t=(g^{t/2})^2$ este pătrat; reciproc, $\cl a=\cl x^2$ cu $\cl x=\cl g^{\,s}$ dă $t\equiv 2s\pmod{p-1}$, deci $t$ par ($p-1$ fiind par).
\end{proof}

\begin{corolar}\label{cor:minus1}
$-1$ este rest pătratic modulo primul impar $p$ $\iff p\equiv1\pmod4$. (Se aplică criteriul: $(-1)^{\frac{p-1}2}=1\iff\frac{p-1}2$ par.)
\end{corolar}

\begin{propozitie}[sumele de puteri modulo $p$]\label{prop:sumeputeri}
Pentru $p$ prim și $k\ge 1$:
\[
1^k+2^k+\dots+(p-1)^k\equiv
\begin{cases}
-1\pmod p,&\text{dacă }(p-1)\mid k,\\
\ \ 0\pmod p,&\text{altfel}.
\end{cases}
\]
\end{propozitie}

\begin{proof}
Dacă $(p-1)\mid k$, fiecare termen este $\equiv1$ (Fermat), iar suma este $p-1\equiv-1$. Altfel, fie $g$ o rădăcină primitivă; suma este $S=\sum_{j=0}^{p-2}\cl g^{\,jk}$, iar
\[
(\cl g^{\,k}-\cl 1)\,S=\cl g^{\,k(p-1)}-\cl1=\cl0.
\]
Cum $(p-1)\nmid k$, avem $\cl g^{\,k}\neq\cl1$, deci $\cl g^{\,k}-\cl1$ este inversabil în corpul $\ZZ_p$ și $S=\cl0$.
\end{proof}

\section{Interpolarea Lagrange și polinoame auxiliare}

\begin{teorema}[interpolarea Lagrange]\label{teo:lagrangeint}
Fie $K$ un corp comutativ, $a_0,\dots,a_n\in K$ distincte și $b_0,\dots,b_n\in K$. Există un \textbf{unic} polinom $L\in K[X]$ cu $\deg L\le n$ și $L(a_i)=b_i$ pentru toți $i$, anume
\[
L(X)=\sum_{i=0}^{n}b_i\prod_{j\neq i}\frac{X-a_j}{a_i-a_j}.
\]
\end{teorema}

\begin{proof}
Fiecare produs $\prod_{j\neq i}\frac{X-a_j}{a_i-a_j}$ ia valoarea $1$ în $a_i$ și $0$ în celelalte noduri, deci $L(a_i)=b_i$; gradul este evident $\le n$. Unicitatea: diferența a două soluții are grad $\le n$ și $n+1$ rădăcini, deci este nulă (Corolarul~\ref{cor:identitate}).
\end{proof}

\begin{aplicatie}[polinomul auxiliar --- problemă clasică]\label{apl:auxiliar}
Fie $f\in\RR[X]$ de grad $n$ cu $f(k)=\dfrac{k}{k+1}$ pentru $k=0,1,\dots,n$. Calculați $f(n+1)$.

\emph{Soluție.} Polinomul $g(X)=(X+1)f(X)-X$ are gradul $n+1$ și se anulează în $0,1,\dots,n$, deci
$g(X)=c\,X(X-1)\cdots(X-n)$. Evaluăm în $X=-1$: $g(-1)=0\cdot f(-1)+1=1$, iar produsul dă $c\,(-1)^{n+1}(n+1)!$, de unde $c=\dfrac{(-1)^{n+1}}{(n+1)!}$. Atunci $g(n+1)=c\,(n+1)!=(-1)^{n+1}$ și
\[
f(n+1)=\frac{g(n+1)+(n+1)}{n+2}=
\begin{cases}
1,&n\text{ impar},\\[2pt]
\dfrac{n}{n+2},&n\text{ par}.
\end{cases}
\]
\emph{Morala:} când se dau valorile lui $f$ în noduri, construiți un polinom auxiliar ale cărui rădăcini sunt nodurile; gradul și un nod suplimentar determină constanta.
\end{aplicatie}

\vspace{6pt}
\begin{center}
\rule{0.5\textwidth}{0.4pt}\\[4pt]
\small\emph{Succes la olimpiadă!}
\end{center}









